% Generated from Paul/Mathematica/Gaudin/notes-src/ on 2026-10-16; may be edited directly thereafter (the generator is optional).
\documentclass[11pt]{article}
\usepackage{lmodern}
\usepackage[T1]{fontenc}
\usepackage[utf8]{inputenc}
\usepackage{amsmath,amssymb}
\usepackage{booktabs,longtable,array}
\usepackage[margin=1in]{geometry}
\usepackage{xcolor}
\usepackage[colorlinks=true,linkcolor=blue!50!black,urlcolor=blue!50!black,citecolor=blue!50!black]{hyperref}
\usepackage{microtype}
\setlength{\parskip}{0.4em}
\setlength{\parindent}{0pt}
\setlength{\LTcapwidth}{\textwidth}
\allowdisplaybreaks
\emergencystretch=2em
\sloppy


\title{Clebsch--Gordan coefficients of $\mathrm{gl}(3)$ from separation of variables in the two-site Gaudin model: working notes on rounds 1--4}

\author{Paul Ryan (project notes, compiled from the work of the 2026-10-15/16 agent rounds)}

\date{16 October 2026}

\begin{document}
\maketitle

\begin{abstract}
We want the Clebsch--Gordan coefficients of $\mathrm{gl}(3)$ for $V_\Lambda\otimes V_\Lambda$, and we want them as explicit expressions in Q-functions, obtained from separation of variables (SoV) in the two-site rational Gaudin model.
These notes collect everything that four rounds of computation (chain side in Wolfram, Q side and kernels in Python, about fourteen agents) have established, and say honestly what has not been established.

The main finding is that the overlap of a CG-complete two-vacua covector family with an eigenstate of the Gaudin Hamiltonians is a Q-number: a bilinear form, with a state-independent kernel, in the Plücker coordinates of the order-$(\lambda_1+1)$ jets at $\theta_1$ and $\theta_2$ of the dual Q-functions.
For the adjoint the kernel is a $5\times5$ matrix of polynomials in $(\omega_1,\omega_1',1/\theta_{12})$; it was found as the formal $h\to0$ limit of the XXX ten-product formula, re-found by a direct fit, and derived from first principles from the Gaudin jet basis.
Writing the Plücker coordinates through the Frobenius solution at $\theta_s$ shows that the overlap depends only on the Laurent data $(H_s,a_s)$ of the Hamiltonians at the two sites.
Taking the untwisting limit $\zeta\to0$ of this expression gives, for all six highest-weight branches of the adjoint, the CG generating function of $\mathrm{adj}\otimes\mathrm{adj}$ as an exact polynomial identity, once the $\epsilon$-series of the Laurent data are known; those series are exact, polynomial in the twist direction, and algebraic with a square root on the $(4,1,1)$ and $(3,3,0)$ branches.
For general $\Lambda$ we have closed forms for the one-site kernel, a block theorem reducing every two-row representation to the one-row master kernel, and the master kernel to first order in $1/\theta_{12}$ (all orders for $[l,l,0]$).
No per-sector cleaning of the XXX kind is needed anywhere.

What is open: the master kernel beyond first order, the Laurent-data formulas and the limit for $\lambda_1\ge3$, a derivation of the $\epsilon$-series, and the norm of the untwisted eigenvectors from the Q side.
Appendices hold all the data (kernels, series, CG tables, untwisted Laurent data) and a reproducibility map.
\end{abstract}

\tableofcontents
\bigskip

\setcounter{section}{-1}

\section{Overview: what the program is, what we found, what is open}\label{sec:overview}

\textbf{Punchline.} For the two-site $\mathrm{gl}(3)$ Gaudin model, the overlap of a suitable pair of rotated vacua with an eigenstate is a Q-number that depends on the eigenstate only through the Laurent data of the Gaudin Hamiltonians at the two sites; sending the twist to zero turns this Q-number into the CG generating function.
We have this fully for the adjoint, with the main structural pieces (kernel, Frobenius reduction, representation-independent rules) in place for $[3,1,0]$, $[3,2,0]$ and, at the level of kernels, many more.

\subsection{The program in a nutshell}\label{sec:nutshell}

The goal is simple to state.
Take two copies of an irreducible $\mathrm{gl}(3)$ representation $V_\Lambda$, $\Lambda=(\lambda_1,\lambda_2,0)$, and ask for the Clebsch--Gordan coefficients $\langle e_\nu\otimes e_1|\Psi_0\rangle$ of the highest-weight (HW) vectors $\Psi_0$ in $V_\Lambda\otimes V_\Lambda$. The idea is to realise $V_\Lambda\otimes V_\Lambda$ as the Hilbert space of a two-site Gaudin chain with inhomogeneities $\theta_1,\theta_2$ and a diagonal twist $g=\mathrm{diag}(\zeta_1,\zeta_2,\zeta_3)$. The Gaudin Hamiltonians are the Laurent coefficients of the quantum characteristic polynomial of the Lax matrix $L_g(u)=g+\sum_sE^{(s)}/(u-\theta_s)$. Their eigenvectors for generic $\zeta$ are labelled by Q-functions, and as $\zeta\to0$ (in a controlled direction $\zeta=\epsilon\phi$) the eigenvectors that are regular become the HW vectors of the untwisted $V_\Lambda\otimes V_\Lambda$. SoV tells us how to write the overlap of any covector with such an eigenvector as a determinant or bilinear form of Q-functions.
If we pick covectors whose overlaps with the untwisted HW kets are exactly the CG coefficients (suitably rotated), the CG coefficients become Q-expressions.

This is the ``XXX program'' of the earlier notes (Lax matrix $u+hE$, companion twist, SoV bases, localising transfer matrices), transplanted to the rational Gaudin model, which is the $h\to0$ limit of the XXX chain with twists $z=e^{h\zeta}$. We use the XXX chain as a scaffold (Section~\ref{sec:scaffold}) but the deliverable is Gaudin.\footnote{Paul's steer of 2026-10-16: ``the goal of the project is Gaudin anyway, not XXX; I assumed XXX would be simpler''; and later ``project target is gl(3) CG coefficients from SoV in the Gaudin model''. \texttt{HANDOFF.\allowbreak{}md}, entries ``ROUND 2'' and ``ROUND 4''. The repository's \texttt{CLAUDE.\allowbreak{}md} still states a su(4) goal; that sentence is outdated and was deliberately not edited by the notes writer.}

\subsection{What we found, in eight points}\label{sec:eight}

\begin{enumerate}
\item \textbf{The chain side works without cleaning.} The family $\Phi_\infty(\omega_1,\omega_1')=\Omega^\infty[\zeta,\omega_1]\otimes\Omega^\infty[\zeta,\omega_1']$ of products of two rotated vacua is CG-complete, its exact undo is $\langle\varphi_u\otimes e_1|$ with $\epsilon$-free Gauss parameters, and the contamination ratio $K/S$ is $O(\epsilon)$ on every HW branch of $[2,1,0]$, $[3,1,0]$, $[3,2,0]$, including the multiplicity-2 sectors that needed per-copy cleaning in XXX (Section~\ref{sec:chain}).
\item \textbf{The overlap is a Q-number (the jet formula).} Up to a frame-tie scalar $c$ and a vacuum monomial $M$, the overlap is a bilinear form, with a state-independent kernel $K$, in Plücker coordinates of the jets of the dual Q-functions at $\theta_1$ and $\theta_2$, equation~(\ref{eq:jet0}) below (Section~\ref{sec:jet}).
\item \textbf{Where the kernel comes from.} The adjoint kernel was obtained three independent ways: as the analytic $h\to0$ limit of the XXX ten-product formula, by a direct fit, and as the Gaudin K tensor in the jet basis that replaces the SoV basis when $h\to0$ (Sections~\ref{sec:jet} and~\ref{sec:B}).
\item \textbf{The overlap depends only on the Laurent data} $(H_s,a_s)$ at $\theta_1,\theta_2$, because every Plücker ratio is fixed by the Frobenius recursion at $\theta_s$ (Section~\ref{sec:nu}).
\item \textbf{The kernel has closed forms.} One site: a decomposable bivector, explicitly a dressed Wronskian.
Two sites: a block theorem relating every $[\lambda_1,\lambda_2,0]$ to the one-row master kernel, the first-order term in $1/\theta_{12}$ for all pairs, and all orders for $[l,l,0]$ (Section~\ref{sec:kernel}).
\item \textbf{The untwisting limit is established for the adjoint}, as a polynomial identity in the rotation parameters, for all six HW branches; the CG coefficients come out $\phi$-independent to 77 digits (Section~\ref{sec:untwist}).
\item \textbf{The $\epsilon$-series of the Laurent data are exact} and partly algebraic (Section~\ref{sec:series}), and the untwisted model has a rich structure of its own: Fuchsian at infinity, three polynomial solutions, Casimir closed forms (Section~\ref{sec:untwisted}).
\item \textbf{Several of our own early expectations were wrong}, and we collect them in Section~\ref{sec:surprises} because they are as instructive as the successes.
\end{enumerate}

The formula of point 2, written out once so that the reader has it in front of them, is

\begin{equation}
c\,\frac{\langle\Phi_\infty|\Psi\rangle}{\langle x_0^G|\Psi\rangle}\,M(\theta_1)M(\theta_2)=\sum_{(kl),(k'l')}K_{(kl),(k'l')}(\omega_1,\omega_1',\theta_{12})\,(P^3p_{kl})(\theta_1)\,(P^3p_{k'l'})(\theta_2),\label{eq:jet0}
\end{equation}

with $M=P^3p_{0,\lambda_2+1}$. Here $\Phi_\infty$ is the two-vacua covector, $\Psi$ any eigenstate, $x_0^G$ a reference covector, $P^2,P^3$ two of the three dual Q-functions of $\Psi$, $p_{kl}=P^{2(k)}P^{3(l)}-P^{3(k)}P^{2(l)}$ their Plücker coordinates at $\theta_s$, and $(kl)$ runs over the $\binom{\lambda_1+2}{2}-\binom{\lambda_2+1}{2}$ pairs $0\le k<l\le\lambda_1+1$ with $l\ge\lambda_2+1$. Every symbol is defined properly in Sections~\ref{sec:setting} to~\ref{sec:jet}.

\subsection{What is open}\label{sec:open}

The master kernel $R^{\lambda_1}$ for all four jet indices beyond first order; the Laurent-data (Frobenius) form of the overlap and the untwisting limit for $\lambda_1\ge3$; a derivation, not a fit, of the $\epsilon$-series of the Laurent data (the degenerate solver for the top state and the octets); the norm $\langle\Psi_0|\Psi_0\rangle$ from the Q side, which would remove the last piece of chain input; the multiplicity-3 sector $[4,2,0]$; and the write-up as a paper.
Section~\ref{sec:outlook} turns this into the round-5 plan.

\subsection{How to read these notes}\label{sec:reading}

Each section opens with its punchline in two or three sentences, then tells the story, then gives the data.
The main text works one example, the adjoint $[2,1,0]$, all the way through; bulk data live in the appendices (A conventions, B kernels, C one-site kernel, D Laurent series and CG tables, E untwisted data, F reproducibility).
A reader who skips all appendices still gets the whole story.
Every number carries a provenance footnote: a file and a section or cell.
Where sources disagreed, the later round, or the explicit reconciliation in \texttt{HANDOFF.\allowbreak{}md}, wins, and we say so in a footnote.
Table~\ref{tab:status} is the map of what has been done for which representation.

{\scriptsize
\begin{longtable}{lrrrcccccc}
\caption{Status map.
``Chain'' = chain-side checks (K/S, CG at leading order, Section~\ref{sec:chain}); ``Q'' = Q-side rules (oper, Wronskians, $\det N$, Section~\ref{sec:q}); ``jet'' = jet formula verified; ``kernel'' = closed forms (Section~\ref{sec:kernel}); ``limit'' = analytic untwisting limit (Section~\ref{sec:untwist}); ``exact'' = exact untwisted reference data (Appendices D, E).
All entries are ours to the precision stated in the text; ``fit'' = determined by a unique exact fit, not derived.}\label{tab:status}\\
\toprule
$\Lambda$ & $\dim V$ & states & HW branches & chain & Q & jet & kernel & limit & exact \\
\midrule
\endfirsthead
\toprule
$\Lambda$ & $\dim V$ & states & HW branches & chain & Q & jet & kernel & limit & exact \\
\midrule
\endhead
\bottomrule
\endlastfoot
$[1,1,0]$ & 3 & 9 & 2 & -- & -- & fit & exact & -- & yes \\
$[2,1,0]$ (adjoint) & 8 & 64 & 6 & yes & yes & derived & exact & yes & yes \\
$[2,2,0]$ & 6 & 36 & 3 & -- & -- & fit & exact & -- & yes \\
$[3,1,0]$ & 15 & 225 & 10 & yes & yes & fit & laws & -- & yes \\
$[3,2,0]$ & 15 & 225 & 10 & yes & yes & fit & laws & -- & yes \\
$[4,1,0]$ & 24 & 576 & 14 & -- & -- & base point & block thm & -- & yes \\
$[l,0,0]$, $l\le4$; $[3,3,0]$, $[4,4,0]$ & 3--15 & 9--225 & -- & -- & -- & fit & exact & -- & -- \\
\end{longtable}}


\section{Setting and conventions}\label{sec:setting}

\textbf{Punchline.} The commuting family is the set of Laurent coefficients of a column-ordered quantum determinant; on any state it is described by six numbers $(H_s,a_s,b_s)_{s=1,2}$, with a scalar double pole $c_2=\lambda_2(\lambda_1+1)$ and no triple pole.
The bra vacua are $e_1^T\rho(n)$ with explicit Gauss parameters, and they live in a frame tied to the diagonal frame of the eigenstates by an explicit scalar $\kappa^G$.

\subsection{The model and its Hamiltonians}\label{sec:model}

We work with $\mathfrak{gl}_3$ and a two-site chain; both sites carry the same irreducible representation $V_\Lambda$, $\Lambda=(\lambda_1,\lambda_2,\lambda_3)$ with $\lambda_3=0$ throughout.\footnote{All representations are written in the Gelfand--Tsetlin (GT) basis with plain, non-Hermitian rational generators $E_{ij}$ and the invariant diagonal metric $G$ ($E_{ij}^TG=GE_{ji}$); this is the convention of \texttt{Paul/\allowbreak{}Python/\allowbreak{}gl3\_\allowbreak{}reps.\allowbreak{}py} and of the notebooks' \texttt{JJ}. Raising operators are $E_{ij}$ with $i<j$; $e_1$ is the highest-weight ket and covectors are rows acted on from the right, so $e_1^TE_{ij}\neq0$ only for $i\le j$. Sources: \texttt{Gaudin-\allowbreak{}sl3-\allowbreak{}CG-\allowbreak{}plan.\allowbreak{}md} section 1; \texttt{RESULTS-\allowbreak{}Qside.\allowbreak{}md} section 0.} On site $s$ the generators are $E^{(1)}=E\otimes1$, $E^{(2)}=1\otimes E$, they obey $[E_{ij},E_{kl}]=\delta_{jk}E_{il}-\delta_{li}E_{kj}$, and the inhomogeneities are $\theta=(\theta_1,\theta_2)$, $\theta_{12}=\theta_1-\theta_2$. There is no $h$ anywhere in the Gaudin model.
The twist is a Lie-algebra element $g=\mathrm{diag}(\zeta_1,\zeta_2,\zeta_3)$ and the Lax matrix is

\begin{equation}
L_g(u)_{ij}=\zeta_i\,\delta_{ij}+\sum_{s=1}^2\frac{E^{(s)}_{ij}}{u-\theta_s}.\label{eq:Lax}
\end{equation}

The commuting family is generated by the column-ordered quantum determinant of $\partial_u-L_g(u)$, with the derivative moved to the right ($\partial f=f\partial+f'$):

\begin{equation}
\mathrm{cdet}(\partial_u-L_g)=\sum_{\sigma\in S_3}\mathrm{sgn}\,\sigma\; M_{\sigma(1)1}M_{\sigma(2)2}M_{\sigma(3)3}=\partial^3-S_1(u)\,\partial^2+S_2(u)\,\partial-S_3(u),\qquad M=\partial_u\mathbf 1-L_g.\label{eq:cdet}
\end{equation}

Writing $l_{ij}=L_g(u)_{ij}$ and $l'_{ij}=\partial_ul_{ij}$, $l''_{ij}=\partial_u^2l_{ij}$, and keeping operator products in the order written, the coefficients are

\begin{equation}
\begin{aligned}
S_1&=l_{11}+l_{22}+l_{33},\\
S_2&=\sum_{i<j}\bigl(l_{ii}l_{jj}-l_{ji}l_{ij}\bigr)-l'_{22}-2\,l'_{33},\\
S_3&=l''_{33}-l_{11}l'_{33}-l_{22}l'_{33}-l'_{22}l_{33}+l_{31}l'_{13}+l_{32}l'_{23}+l'_{32}l_{23}\\
&\quad+l_{11}l_{22}l_{33}-l_{11}l_{32}l_{23}-l_{21}l_{12}l_{33}+l_{21}l_{32}l_{13}+l_{31}l_{12}l_{23}-l_{31}l_{22}l_{13}.
\end{aligned}\label{eq:S123}
\end{equation}

We checked that $[S_a(u),S_b(v)]=0$ exactly (exact rationals, 3$\times$3 grids of rational $u,v$, for $[2,1,0]$, $[3,1,0]$, $[3,2,0]$), that $S_1$ is a scalar and that the weights are conserved.\footnote{Agents W1 (\texttt{Gaudin-\allowbreak{}L2-\allowbreak{}chain.\allowbreak{}wb}, summary cell), P1 (\texttt{RESULTS-\allowbreak{}Qside.\allowbreak{}md} section 0), P3, P4.
The convention was not tuned: it worked at the first attempt.
The row-ordered product with $E_{ij}$ and the column-ordered product with $E_{ji}$ do \emph{not} commute (274 non-zero entries in $[S_2,S_3]$); the doubly flipped variant is just the transposed chain.}

\subsection{Laurent data}\label{sec:laurent}

Every non-scalar information about a state is contained in the Laurent coefficients at the two sites:

\begin{equation}
\begin{aligned}
S_1&=\sigma+|\Lambda|\sum_s\frac1{u-\theta_s},\qquad \sigma=\zeta_1+\zeta_2+\zeta_3,\\
S_2&=e_2(\zeta)+\sum_s\Bigl[\frac{H_s}{u-\theta_s}+\frac{c_2}{(u-\theta_s)^2}\Bigr],\\
S_3&=e_3(\zeta)+\sum_s\Bigl[\frac{b_s}{u-\theta_s}+\frac{a_s}{(u-\theta_s)^2}\Bigr],
\end{aligned}\label{eq:laurent}
\end{equation}

where $e_k(\zeta)$ are elementary symmetric polynomials, $|\Lambda|=\lambda_1+\lambda_2+\lambda_3$, and

\begin{equation}
c_2=\lambda_2(\lambda_1+1)\label{eq:c2}
\end{equation}

is a scalar. $S_3$ has no triple pole.
On an eigenstate the six operators $H_1,H_2,a_1,a_2,b_1,b_2$ become six numbers; this is all the state-dependent data that the Gaudin Hamiltonians carry.
They are not independent: three no-log conditions per site (Section~\ref{sec:nolog}) cut the solution set down to a finite one.

The value $c_2=\lambda_2(\lambda_1+1)$ deserves a remark, because we got it wrong twice.
It equals $|\Lambda|$ for $[2,1,0]$ and $[3,1,0]$, so those two representations cannot tell the two apart; the first guess, half the quadratic Casimir, gives the right $3$ for $[2,1,0]$ but $6$ for $[3,1,0]$, whose true value is $4$. The scan over representations is in Table~\ref{tab:c2}.
For one-row representations ($\lambda_2=0$) the double pole is the scalar $0$.\footnote{Agent P3 (\texttt{RESULTS-\allowbreak{}Qside-\allowbreak{}310.\allowbreak{}md} section 0) found $c_2=\lambda_2(\lambda_1+1)$ on $(1,1,0),(2,1,0),(3,1,0),(2,2,0),(3,2,0)$ and reported the double pole as ``not scalar'' for $(n,0,0)$; agent G1 (\texttt{RESULTS-\allowbreak{}general-\allowbreak{}kernel.\allowbreak{}md} K2) later checked $[1,0,0],[2,0,0],[3,0,0]$ on $9/36/100$ states and found the scalar $0$ every time.
We take G1's statement, as the later round. $(2,1,1)$ has $c_2=8$ but an $S_3$ triple pole and was not used.}

{\small
\begin{longtable}{lcccc}
\caption{The scalar double pole $c_2$ of $S_2$ and the induced exponents.
Dual exponents at $\theta_s$: $\{0,\lambda_2+1,\lambda_1+2\}$; direct exponents $\{0,\lambda_1-\lambda_2+1,\lambda_1+2\}$. Sources: P3 (\texttt{RESULTS-\allowbreak{}Qside-\allowbreak{}310.\allowbreak{}md}), P4 (\texttt{RESULTS-\allowbreak{}Qside-\allowbreak{}320.\allowbreak{}md}), W2/W4 (\texttt{Gaudin-\allowbreak{}L2-\allowbreak{}chain-\allowbreak{}2.\allowbreak{}wb}, \texttt{Gaudin-\allowbreak{}untwisted-\allowbreak{}exact.\allowbreak{}wb}).}\label{tab:c2}\\
\toprule
$\Lambda$ & $c_2=\lambda_2(\lambda_1+1)$ & $|\Lambda|$ & dual exponents & direct exponents \\
\midrule
\endfirsthead
\toprule
$\Lambda$ & $c_2=\lambda_2(\lambda_1+1)$ & $|\Lambda|$ & dual exponents & direct exponents \\
\midrule
\endhead
\bottomrule
\endlastfoot
$[1,1,0]$ & 2 & 2 & $\{0,2,3\}$ & $\{0,1,3\}$ \\
$[2,1,0]$ & 3 & 3 & $\{0,2,4\}$ & $\{0,2,4\}$ \\
$[3,1,0]$ & 4 & 4 & $\{0,2,5\}$ & $\{0,3,5\}$ \\
$[2,2,0]$ & 6 & 4 & $\{0,3,4\}$ & $\{0,1,4\}$ \\
$[3,2,0]$ & 8 & 5 & $\{0,3,5\}$ & $\{0,2,5\}$ \\
$[l,0,0]$ & 0 & $l$ & $\{0,1,l+2\}$ & $\{0,l+1,l+2\}$ \\
\end{longtable}}


The first Laurent coefficient has an explicit form, and its sign cost us a round-1 correction.
With $C_\zeta^{(s)}=\sum_i\zeta_iE^{(s)}_{ii}$ and $W=\sum_k\zeta_k(E^{(1)}+E^{(2)})_{kk}$,

\begin{equation}
H_1=-\Bigl[\sum_i\zeta_iE^{(1)}_{ii}+\frac1{\theta_{12}}\sum_{i,k}E^{(1)}_{ik}E^{(2)}_{ki}\Bigr]+|\Lambda|\sigma+\frac{|\Lambda|^2}{\theta_{12}},\qquad H_1+H_2=-W+2|\Lambda|\sigma,\label{eq:H1}
\end{equation}

and $H_2$ is the same with $1\leftrightarrow2$ and $\theta_{12}\to-\theta_{12}$. The sign is $-$ in front of the bracket, not $+$ as the first plan had it: it follows from the residue at $\theta_1$ of $\sum_{i<j}(l_{ii}l_{jj}-l_{ji}l_{ij})$.\footnote{Agent W1, \texttt{Gaudin-\allowbreak{}L2-\allowbreak{}chain.\allowbreak{}wb} summary cell (``Sign difference to the plan''), exact on $[2,1,0]$; agents P3 and P4 report the same explicit form for $[3,1,0]$ and $[3,2,0]$ (exact fits, 0 mismatches).
The plan \texttt{Gaudin-\allowbreak{}sl3-\allowbreak{}CG-\allowbreak{}plan.\allowbreak{}md} section 1 carries the corrected row.} The same formula at $\zeta=0$ gives $H_1=\bigl(|\Lambda|^2-\tfrac12[C_2(\mu)-2C_2(\Lambda)]\bigr)/\theta_{12}$ on a state of total weight $\mu$ (Section~\ref{sec:casimir}).

\subsection{The dictionary to XXX}\label{sec:dict}

The Gaudin family is the limit of the XXX transfer matrices at twist $z=e^{h\zeta}$. Expanding $t_1(u)/\nu(u)=3+hA_1+h^2A_2+h^3A_3+\dots$ and $t_2(u)/(\nu(u)\nu(u-h))=3+hB_1+h^2B_2+\dots$ with $\nu(u)=\Pi(u)=(u-\theta_1)(u-\theta_2)$ one finds, exactly,

\begin{equation}
\begin{aligned}
&A_1=S_1,\quad B_1=2S_1,\quad B_2=A_2+\tfrac12S_1^2+\tfrac12S_1',\\
&S_2=-A_2+\tfrac12S_1^2+\tfrac12S_1',\quad S_3=A_3-B_3+\tfrac16S_1^3+\tfrac12S_1S_1'+\tfrac12S_1'',
\end{aligned}\label{eq:dict}
\end{equation}

with $S_k(u\to\infty)=e_k(\zeta)$. The relations $A_1=S_1$ and $B_1=2S_1$ are forced by the Baxter equations at orders $h^1$ and $h^2$; nothing in the dictionary needs more than $S_1$ and its derivatives.\footnote{Agent P2, \texttt{RESULTS-\allowbreak{}xxx-\allowbreak{}limit.\allowbreak{}md} section X2: exact rational series in $h$, checked at four values of $u$ and three twist sets, and at finite $h$ to $O(h^3)$ residuals.
P2 also found the vacuum, frame and Baxter-to-oper statements used below.} This is why the cdet convention could be pinned independently of the chain-side commutation test: the plan's convention is the only one that fits.

\subsection{Frames, vacua and the reference covector}\label{sec:frames}

Eigenstates $\Psi_B$ of the diagonal-twist Hamiltonians live in what we call the \emph{diagonal frame}.
The SoV construction naturally lives in the \emph{companion frame}, where the twist is replaced by its Lie-algebra companion matrix, built from $e_k=e_k(\zeta)$,

\begin{equation}
g_c(\zeta)=\begin{pmatrix}e_1&-e_2&e_3\\1&0&0\\0&1&0\end{pmatrix},\qquad P_{\rm aux}(\zeta)\ \text{with rows}\ (1,\ \zeta_i-e_1,\ \zeta_i^2-e_1\zeta_i+e_2),\label{eq:gc}
\end{equation}

so that $P_{\rm aux}g_cP_{\rm aux}^{-1}=\mathrm{diag}(\zeta)$. The physical-space image of the frame change is $P^G_{\rm phys}(\zeta)=\rho\bigl((P_{\rm aux}(\zeta)^{-1})^T\bigr)$, and the reference covector is $x_0^G=x(\zeta)\otimes x(\zeta)$ with $x(\zeta)=e_1^T\bigl(P^G_{\rm phys}\bigr)^{-1}$ (one site).
For the adjoint, in notebook order,

\begin{equation}
x(\zeta)=\bigl(\zeta_2-\zeta_1,\ \zeta_3-\zeta_1,\ \zeta_2-\zeta_1,\ 2\zeta_3-\zeta_1-\zeta_2,\ 3(\zeta_2-\zeta_1),\ 3(\zeta_3-\zeta_1),\ 2(\zeta_3-\zeta_2),\ 3(\zeta_3-\zeta_2)\bigr),\label{eq:x0}
\end{equation}

a covector whose components are homogeneous of degree $\lambda_2$ in $\zeta$ for every representation we have looked at.\footnote{Exactly linear: agent D2, \texttt{RESULTS-\allowbreak{}twovac-\allowbreak{}limit.\allowbreak{}md} L1(b), with the XXX $x_0(z)$ exactly linear in $z$ at every order in $h$; $x_0(e^{h\zeta})=h\,x(\zeta)+O(h^2)$. Degree $\lambda_2$ in general: agent W4, \texttt{Gaudin-\allowbreak{}untwisted-\allowbreak{}exact.\allowbreak{}wb} E5 and the \texttt{XVEC} records of the data files ($[1,1,0]$: $(z_2-z_1,z_3-z_1,z_3-z_2)$; $[2,2,0]$: quadratic).} We expected $x_0^G$ not to be parallel to the $h\to0$ limit of the XXX $x_0$; it is, to $10^{-31}$ (norms $\sim h$ per site).

The twist of the bra vacuum is conjugated into the diagonal frame, $g'=P_{\rm aux}(\zeta)\,g_c(\omega)\,P_{\rm aux}(\zeta)^{-1}$, where $\omega=(\omega_1,\omega_2,\omega_3)$ are the three eigenvalues of the vacuum's own twist.
The bra vacuum is

\begin{equation}
\Omega_{\zeta\omega}=e_1^T\rho(n),\qquad n=\exp(r_1E_{12})\exp(r_2E_{13})\exp(r_3E_{23}),\label{eq:vac}
\end{equation}

\begin{equation}
r_1=\frac{\omega_1-\zeta_1}{\omega_1-\zeta_2},\quad r_2=\frac{(\omega_1-\zeta_1)(\zeta_1-\omega_2)(\zeta_2-\zeta_3)}{(\zeta_1-\zeta_2)(\omega_1-\zeta_3)(\omega_2-\zeta_3)},\quad r_3=\frac{(\omega_1-\zeta_2)(\zeta_2-\omega_2)(\zeta_1-\zeta_3)}{(\zeta_1-\zeta_2)(\omega_1-\zeta_3)(\zeta_3-\omega_2)}.\label{eq:rk}
\end{equation}

We verified that $\Omega_{\zeta\omega}\otimes\Omega_{\zeta\omega}$ is an exact left eigencovector of all coefficient operators of $S_2,S_3$ built with the twist $g'$, for $[2,1,0]$, $[3,1,0]$, $[3,2,0]$ at three twist points; the four other orientations of the conjugation, the transposes, and the naive control $\mathrm{diag}(\omega)$ all fail.\footnote{Agents P1, P2 (independent), W1; \texttt{RESULTS-\allowbreak{}Qside.\allowbreak{}md} section 0, \texttt{RESULTS-\allowbreak{}xxx-\allowbreak{}limit.\allowbreak{}md} X1, \texttt{Gaudin-\allowbreak{}L2-\allowbreak{}chain.\allowbreak{}wb}. The $r_k$ are the same degree-zero homogeneous functions as in XXX: $r_k(z,w)\to r_k(\zeta,\omega)$ with $O(h)$ corrections, and that is why the Gaudin vacuum is \emph{exactly} $\epsilon$-independent in the regime $\zeta=\epsilon\phi$, $\omega=\epsilon\hat\omega$. At $\zeta=(3/2,5/4,2/3)$, $\omega=(5/2,7/3,3/5)$ the Gauss parameters are $r=(4/5,-7/11,65/44)$, rep-independent.} The covector $\Omega_{\zeta\omega}$ is the vacuum of a \emph{different} twist than the ket's; its overlaps with the $\zeta$-eigenstates are the objects the SoV machinery turns into Q-functions.

At $\omega_2\to\infty$ the vacuum has a finite limit $\Omega^\infty[\zeta,\omega_1]=e_1^T\rho(n^\infty)$ with

\begin{equation}
r_1=\frac{\omega_1-\zeta_1}{\omega_1-\zeta_2},\quad r_2=-\frac{(\omega_1-\zeta_1)(\zeta_2-\zeta_3)}{(\zeta_1-\zeta_2)(\omega_1-\zeta_3)},\quad r_3=\frac{(\omega_1-\zeta_2)(\zeta_1-\zeta_3)}{(\zeta_1-\zeta_2)(\omega_1-\zeta_3)},\label{eq:rkinf}
\end{equation}

and it depends on only two parameters in the two-site family $\Phi_\infty=\Omega^\infty[\zeta,\omega_1]\otimes\Omega^\infty[\zeta,\omega_1']$. In the companion frame the same vacuum has Gauss parameters $(1/\omega_1,\,-1/(\omega_1\omega_2),\,(\omega_1+\omega_2)/(\omega_1\omega_2))$, i.e. exactly the XXX ones with $w\to\omega$, and at $\omega_2=\infty$ it is $\Omega_c^\infty[\omega_1]=e_1^T\rho(e^{aE_{12}}e^{aE_{23}})$ with $a=1/\omega_1$. The two frames are tied by

\begin{equation}
\Omega_D[\zeta,\omega]\,P^G_{\rm phys}(\zeta)=\kappa^G(\zeta,\omega)\,\Omega_c[\omega],\qquad \kappa^G=\Delta_1^{\lambda_1-\lambda_2}\Delta_2^{\lambda_2-\lambda_3}\Delta_3^{\lambda_3},\label{eq:kappaG}
\end{equation}

with $\Delta_1=\omega_1^2/((\zeta_2-\omega_1)(\zeta_3-\omega_1))$ and, at $\omega_2=\infty$, $\Delta_2=\omega_1/((\zeta_1-\zeta_2)(\zeta_3-\omega_1))$. This gives the scalar that appears throughout the kernel normalisation,

\begin{equation}
c=\frac{(\omega_1\omega_1')^{2\lambda_1-\lambda_2}}{\kappa^G(\zeta,\omega_1)\,\kappa^G(\zeta,\omega_1')}=(\zeta_2-\omega_1)^{\lambda_1-\lambda_2}(\zeta_3-\omega_1)^{\lambda_1}(\zeta_2-\omega_1')^{\lambda_1-\lambda_2}(\zeta_3-\omega_1')^{\lambda_1}(\zeta_1-\zeta_2)^{2\lambda_2}.\label{eq:cfactor}
\end{equation}

It was derived for the adjoint ($c=(\zeta_2-\omega_1)(\zeta_3-\omega_1)^2(\ldots)(\zeta_1-\zeta_2)^2$), $[3,1,0]$ and $[3,2,0]$, with the other exponent choices excluded by the fits.\footnote{Agents W2/W3 (frame tie, \texttt{Gaudin-\allowbreak{}L2-\allowbreak{}chain-\allowbreak{}2.\allowbreak{}wb} task A, \texttt{Gaudin-\allowbreak{}L2-\allowbreak{}chain-\allowbreak{}320.\allowbreak{}wb} A2) and P4 (\texttt{RESULTS-\allowbreak{}Qside-\allowbreak{}320.\allowbreak{}md} T0): with the correct $c$ the per-$\zeta$ fits reproduce the stacked kernel to $10^{-71}$ to $10^{-74}$; the wrong exponents $(1,2,2),(1,3,1),(0,3,2),(2,3,2),(1,3,3),(1,4,2)$ give residuals $0.18$ to $0.69$.}

\subsection{The XXX scaffold in one paragraph}\label{sec:scaffold}

All of the above was first understood in the XXX chain with $L=2$ sites in the new conventions (Lax $u\delta_{ij}+hE_{ij}$, companion twist on the right, $B$ and $C$ built from the untwisted $T$).
There, the SoV basis consists of the eigencovectors of $B(u)$, with simple spectrum and labels $(a|bc)$ per site, and the overlap of the vacuum $\Omega_{zw}=e_1^T\rho(n(z,w))$ with an eigenstate is a determinant or a sum of products of ladder brackets of the Q-functions.
For the family of two rotated vacua $\Omega_w\otimes\Omega_{w'}$ the untwisting $z,w,w'\to1$ was clean: $K/S=O(\epsilon)$ on every HW branch, all CG data at leading order, and for the adjoint $(\infty,\infty)$ family the Q-number is a sum of ten products of $2\times2$ bracket determinants.\footnote{\texttt{Paul/\allowbreak{}Mathematica/\allowbreak{}new-\allowbreak{}conventions/\allowbreak{}Two-\allowbreak{}site-\allowbreak{}rotation-\allowbreak{}brackets-\allowbreak{}notes.\allowbreak{}md}, sections 0 and 5 (Stages D and C), notebooks \texttt{L2-\allowbreak{}two-\allowbreak{}vacua-\allowbreak{}untwist.\allowbreak{}wb} and \texttt{L2-\allowbreak{}210-\allowbreak{}two-\allowbreak{}vacua-\allowbreak{}brackets.\allowbreak{}wb}; the XXX Q-number reproduces the chain overlap to $7.9\cdot10^{-111}$ and the six adjoint CG data to $10^{-18}$.} The XXX program also needed, for the single-vacuum approach, a ``cleaning'' step for some sectors (\texttt{Clean-\allowbreak{}CG-\allowbreak{}combinations.\allowbreak{}md}).
The Gaudin model is the $h\to0$ limit, equivalently $\theta_{12}/h\to\infty$ at fixed twist, because of the scaling symmetry $(u-\theta)\to(u-\theta)/\lambda$, $g\to\lambda g$. The XXX companion matrix is not the exponential of a Lie-algebra companion matrix (at $z=1$ it is a Jordan block), which is why the Gaudin companion frame had to be defined separately; the Jordan block turned out to be harmless, since $(G'-1)/h\to g'$ at $O(h)$.

\subsection{Representations, bases and numerical points}\label{sec:points}

We label the GT basis of $V_\Lambda$ in \emph{notebook order}: index $1$ is the highest weight, patterns sorted in descending order.
The Python modules use the \emph{\texttt{gl3\_\allowbreak{}reps} order}, which starts at the lowest-weight vector, so that the HW ket has the last index ($7$ for the adjoint, $14$ for $[3,1,0]$ and $[3,2,0]$).
The map is a permutation, not a reversal: for the adjoint, notebook index $a=0,\dots,7$ corresponds to \texttt{gl3\_\allowbreak{}reps} index $(7,5,6,4,2,1,3,0)$.\footnote{Agent D3, \texttt{RESULTS-\allowbreak{}gaudin-\allowbreak{}B.\allowbreak{}md} header (\texttt{NB2G}), also D1 \texttt{RESULTS-\allowbreak{}twovac-\allowbreak{}fit.\allowbreak{}md} section 0 (the $x_0^G$ direction of P2 mapped with \texttt{NB2G} is reproduced to $10^{-81}$).
Agent W4's data files use 1-based notebook order; P1/P3/P4, D1, U2 use \texttt{gl3\_\allowbreak{}reps} order.
In this text, slices and indices are in notebook order unless said otherwise.} The two-site product basis has index $(i,j)\mapsto(i-1)\dim V+j$, site 1 slow, site 2 fast.

The numerical points reused across all agents are these.
Inhomogeneities $\theta=(2/7,-3/5)$, $\theta_{12}=31/35$. Twist points $P:\ \zeta=(3/2,5/4,2/3)$ and $P':\ \zeta=(2,5/3,-1/2)$, with extra points $(7/5,-3/4,1/6)$, $(-4/3,2/9,5/7)$, $(1/3,-2/7,5/11)$. Vacuum twists $\omega=(5/2,7/3,3/5)$, $\omega'=(7/4,-2/5,9/7)$ and the two-parameter pair $(\omega_1,\omega_1')=(5/2,7/4)$. The untwisting regime is $\zeta=\epsilon\phi$, $\omega=\epsilon\hat\omega$, $\omega'=\epsilon\hat\omega'$ with $\phi=(1/3,-2/7,5/11)$ (control direction $\phi'=(-2/5,3/8,7/11)$), and $\epsilon\in\{1,2,4\}\cdot10^{-6}$ for the chain-side checks and $\epsilon_k=10^{-6}2^{-k}$, $k=0,\dots,11$, for the series fits; precision is 200 digits on the chain side and 60 to 120 digits on the Q side.
The full list of grids is in Appendix~\ref{app:conv}.

\section{Q-functions, Wronskians, \texorpdfstring{$\det N$}{detN}}\label{sec:q}

\textbf{Punchline.} Every eigenstate is a triple of polynomial-times-exponential solutions of one third-order ODE (the oper built from $S_1,S_2,S_3$); the ODE has local exponents $\{0,\lambda_2+1,\lambda_1+2\}$ at each $\theta_s$, which forces linear relations among the jets of the solutions and fixes the Frobenius constant $\kappa_s=a_s/c_2$. The overlap of two states is detected by a $6\times6$ determinant of residue brackets whose columns come from a Lagrange identity.

\subsection{The oper and its Q-functions}\label{sec:oper}

The dual oper is

\begin{equation}
\bigl(\partial^3-S_1\partial^2+S_2\partial-S_3\bigr)P=0,\qquad P^j=e^{\zeta_ju}\,p_j(u),\ \ \deg p_j=\mu_j,\ p_j\ \text{monic},\label{eq:oper}
\end{equation}

where $\mu$ is the weight of the state.
For each of the three exponentials there is exactly one polynomial solution (kernel dimension 1), and $\sum\deg=|\mu|$. The XXX dual Baxter equation becomes exactly this oper at order $h^3$, with $Q^j=z_j^{u/h}q^j\to e^{\zeta_ju}p_j$.\footnote{Agents P1 (\texttt{RESULTS-\allowbreak{}Qside.\allowbreak{}md} Q1; all 64 states of $[2,1,0]$, ODE residual $3.9\cdot10^{-55}$) and P2 (\texttt{RESULTS-\allowbreak{}xxx-\allowbreak{}limit.\allowbreak{}md} X4, sympy, symbolic $h$).}

The direct Q-functions are a little subtler, and this was our first surprise.
One expects them to solve the formal adjoint of (\ref{eq:oper}).
They do not: the plain adjoint has \emph{no} solutions of the form $e^{-\zeta u}\times$polynomial.
What is true is that $\psi_i=\Pi^{-p}Q_i$ solves the adjoint, with $\Pi(u)=(u-\theta_1)(u-\theta_2)$,

\begin{equation}
Q_i=e^{-\zeta_iu}q_i,\quad \deg q_i=2\lambda_1-\mu_i,\qquad \text{direct oper}=\Pi^{p}\circ(\text{formal adjoint})\circ\Pi^{-p},\ \ p=\lambda_1-\lambda_3.\label{eq:direct}
\end{equation}

The XXX Baxter equation at order $h^3$ gives the adjoint acting on $Q_i/\nu^2$, which is the same statement for the adjoint.
For $[2,1,0]$ the gauge power $2$ is forced by the degrees; for $[3,1,0]$ and $[3,2,0]$ a scan over $p=0,\dots,5$ gives no solution at $p=0,1,2$, exactly the right degrees $6-\mu_i$ at $p=3$ in all 36 sectors, and degrees raised by $2(p-3)$ above.
All our data have $\lambda_3=0$, so $p=\lambda_1-\lambda_3$ and $p=\lambda_1$ are not separated by data; the exponent argument below selects the former.\footnote{Agents P1 (gauge $2$, \texttt{RESULTS-\allowbreak{}Qside.\allowbreak{}md} section 0), P3 (\texttt{RESULTS-\allowbreak{}Qside-\allowbreak{}310.\allowbreak{}md} Q1, gauge scan table), P4 (\texttt{RESULTS-\allowbreak{}Qside-\allowbreak{}320.\allowbreak{}md} Q1), P2 (X4: the direct Baxter equation is $\nu^2\circ$adjoint$\circ\nu^{-2}$ exactly).
Wrong-sign and wrong-gauge controls have no solutions to degree 8--10.} Table~\ref{tab:degrees} lists the degrees for the three representations at the top state.

{\small
\begin{longtable}{lcccc}
\caption{Degrees of the Q-functions. $\mu$ is the state weight; dual degrees are $\mu_j$, direct degrees $2\lambda_1-\mu_i$. The vacuum has $\mu=(2\lambda_1,2\lambda_2,0)$. Sources: P1, P3, P4 (all 64/225/225 states: degrees as stated in every sector, kernel dimension 1).}\label{tab:degrees}\\
\toprule
$\Lambda$ & sectors & dual $\deg p_j$ (vacuum) & direct $\deg q_i$ (vacuum) & gauge power $p$ \\
\midrule
\endfirsthead
\toprule
$\Lambda$ & sectors & dual $\deg p_j$ (vacuum) & direct $\deg q_i$ (vacuum) & gauge power $p$ \\
\midrule
\endhead
\bottomrule
\endlastfoot
$[2,1,0]$ & 19 & $(4,2,0)$ & $(0,2,4)$ & 2 \\
$[3,1,0]$ & 36 & $(6,2,0)$ & $(0,4,6)$ & 3 \\
$[3,2,0]$ & 36 & $(6,4,0)$ & $(0,2,6)$ & 3 \\
\end{longtable}}


\subsection{Local exponents and Frobenius constants}\label{sec:frob}

At $\theta_s$ the operator has a regular singular point with indicial polynomial $e(e-1)(e-2)-|\Lambda|e(e-1)+c_2e$, whose roots are

\begin{equation}
\{0,\ \lambda_2+1,\ \lambda_1+2\}\quad\text{(dual)},\qquad\{0,\ \lambda_1-\lambda_2+1,\ \lambda_1+2\}\quad\text{(direct)}.\label{eq:exponents}
\end{equation}

This was confirmed in $450/450$ (state, site) cases for each of $[2,1,0]$, $[3,1,0]$, $[3,2,0]$ from the vanishing orders of the span of the three solutions; for $[3,2,0]$ the direct exponents are $\{0,2,5\}$, not $\{0,3,5\}$, and duality swaps the two patterns.\footnote{Agents P3, P4.
The duality $[3,2,0]\cong[3,1,0]^*\otimes\det^3$ exchanges $\lambda_2+1\leftrightarrow\lambda_1-\lambda_2+1$ between dual and direct opers.} A consequence of the first exponent being $\lambda_2+1$ is that the jets of the dual Q-functions at $\theta_s$ obey linear relations.
For every solution $P$ of the oper,

\begin{equation}
P'(\theta_s)=\kappa_sP(\theta_s),\qquad \kappa_s=\frac{a_s}{c_2},\label{eq:kappa}
\end{equation}

with the same $\kappa_s$ for all three solutions, because the coefficient of $(u-\theta_s)^{-2}$ in $(\partial^3-S_1\partial^2+S_2\partial-S_3)P$ reads $\bigl(I(1)c_1-a_sc_0\bigr)$ with $I(1)=c_2$, and no regular part enters.
For $\lambda_2\ge2$ the second derivative is also tied, $P''(\theta_s)=\kappa_{2,s}P(\theta_s)$ with $\kappa_{2,s}=(a_s^2-H_sa_s+8b_s)/24$ for $[3,2,0]$. We checked the relations on all states and both sites, to $2.7\cdot10^{-78}$ ($\kappa$, adjoint, 80 digits) and $3.1\cdot10^{-57}$, $3.6\cdot10^{-57}$ ($\kappa,\kappa_2$, $[3,2,0]$).\footnote{Agents D1 (\texttt{RESULTS-\allowbreak{}twovac-\allowbreak{}fit.\allowbreak{}md} F1(b): the constant is $a_s/3$ and involves neither $H_s$ nor the regular parts; the spec's expectation that it would was refuted) and P4 (\texttt{RESULTS-\allowbreak{}Qside-\allowbreak{}320.\allowbreak{}md} Q1).}

The jet matrix $[P^{j(k)}(\theta_s)/k!]_{j,k}$ then has rank $1$ for $k\le\lambda_2$ and rank $2$ up to $k=\lambda_1+1$, jumping to $3$ at order $\lambda_1+2$ (measured for $[3,2,0]$: $3\times3$ rank 1, $3\times5$ rank 2, $3\times6$ rank 3, in $450/450$ cases).
Hence the Plücker coordinates

\begin{equation}
p_{kl}=P^{2(k)}P^{3(l)}-P^{3(k)}P^{2(l)}\ \ \text{(at }\theta_s)\label{eq:pluecker}
\end{equation}

vanish for $0\le k<l\le\lambda_2$ (``forced'' pairs), and each non-forced $p_{kl}$ is the $2\times2$ determinant $D_s=c_2d_3-c_3d_2$ of the two non-trivial solutions' leading components times a universal function of the Frobenius coefficients.
For the adjoint $p_{01}=W(\theta_s)=0$ to $1.9\cdot10^{-78}$.

\subsection{The no-log conditions: Gaudin's Bethe equations}\label{sec:nolog}

The three polynomial solutions must be free of logarithms at $\theta_s$. Resonances between the exponents $0,\lambda_2+1,\lambda_1+2$ give three independent conditions per site, six in all, on the six numbers $(H_s,a_s,b_s)$. For $\lambda_2=1$ the first of them is

\begin{equation}
a_s^2-H_sa_s+c_2\,b_s=0,\qquad\text{i.e.}\qquad b_s=\frac{a_s(H_s-a_s)}{c_2},\label{eq:nolog1}
\end{equation}

($c_2=3$ for the adjoint, $4$ for $[3,1,0]$); the other two involve the Taylor coefficients of the regular parts of $S_1,S_2,S_3$ and therefore the other site's data.
For $[3,2,0]$ the first condition arises at order $x^3$ and is cubic.
All hold on every state and both sites to $10^{-56}$ to $10^{-60}$ ($[3,1,0]$: $1.3\cdot10^{-58}$; $[3,2,0]$: $3\cdot10^{-59}$).\footnote{Agents P3 (\texttt{RESULTS-\allowbreak{}Qside-\allowbreak{}310.\allowbreak{}md} Q1, \texttt{frobenius\_\allowbreak{}polys()} in \texttt{qside\_\allowbreak{}310.\allowbreak{}py}) and P4 (\texttt{RESULTS-\allowbreak{}Qside-\allowbreak{}320.\allowbreak{}md} Q1); for the adjoint, U1 (\texttt{RESULTS-\allowbreak{}untwisted-\allowbreak{}Q.\allowbreak{}md} T2, to $10^{-148}$ at 200 digits) and U2.
The number of solutions of the system versus the number of states was not counted.} We call these the Gaudin Bethe equations in Laurent-data form.
In the untwisted model they are what constrains the eigenvalue data in the Casimir language of Section~\ref{sec:casimir}, and perturbatively in $\epsilon$ they determine the series of Section~\ref{sec:series} (except in two degenerate cases).

\subsection{Wronskians and quantum Weyl}\label{sec:wronsk}

The three solutions are linked by Wronskian (QQ) relations with a representation-dependent power of $\Pi$:

\begin{equation}
\begin{aligned}
P^aP^{b\prime}-P^bP^{a\prime}&=(\zeta_b-\zeta_a)\,e^{\sigma u}\,\Pi^{\lambda_2}\,Q_c,\\
Q_aQ_b'-Q_bQ_a'&=(\zeta_a-\zeta_b)\,e^{-\sigma u}\,\Pi^{\lambda_1-\lambda_2}\,P^c,
\end{aligned}\label{eq:QQ}
\end{equation}

for $(a,b,c)$ a permutation of $(1,2,3)$ and monic $p_j,q_i$, so that the constants are exactly $\zeta_b-\zeta_a$ and $\zeta_a-\zeta_b$. The triple Wronskian is $W=\det[P,P',P'']=\prod_{i<j}(\zeta_j-\zeta_i)\,e^{\sigma u}\,\Pi^{|\Lambda|}$, state-independent.
Relative remainders of the polynomial divisions are $\le10^{-48}$ to $10^{-54}$ on all states.\footnote{Agents P1 (adjoint, $F=\Pi$ both sides), P3 ($[3,1,0]$: direct $F=\Pi^2$, dual $F'=\Pi$), P4 ($[3,2,0]$: direct $F=\Pi$, dual $F'=\Pi^2$); plan tables of rounds 2--3 had the dual power wrong for $[3,2,0]$ and were corrected.
The constants and exponential factors are the $h\to0$ images of the XXX $(1/z_b-1/z_a)\chi_3^{-u/h}$ at $z=e^{h\zeta}$.} The rule is $F=\Pi^{\lambda_1-\lambda_2}$ for the direct side, $F'=\Pi^{\lambda_2}=\Pi^{|\Lambda|-p}$ for the dual side.
In particular $W(P^2,P^3)(\theta_s)=0$, and $W'(\theta_s)=(\zeta_3-\zeta_2)e^{\sigma\theta_s}\Pi'(\theta_s)Q_1(\theta_s)$ for the adjoint.

The oper itself is recovered from the three solutions by the quantum-Weyl formulas

\begin{equation}
S_1=\frac{\det[P,P',P''']}{W},\qquad S_2=-\frac{\det[P,P''',P'']}{W},\qquad S_3=\frac{\det[P''',P',P'']}{W},\label{eq:weyl}
\end{equation}

(rows labelled by $j$), verified on all states to $4\cdot10^{-47}$ ($[3,1,0]$) and $1.8\cdot10^{-45}$ ($[3,2,0]$).

\subsection{Orthogonality and the \texorpdfstring{$\det N$}{detN} overlap}\label{sec:detN}

The SoV scalar product of an eigenstate with a vacuum is detected by a $6\times6$ determinant, and this is how we first saw the Q-structure on the Gaudin side.
The orthogonality relation is not of fixed-coefficient type: it follows from the Lagrange identity $\psi\,D_B[P]-P\,D_B^\dagger[\psi]=C(P,\psi)'$ with $P=P^j$ of ket $B$ and $\psi=\Pi^{-p}Q_i$ of bra $A$. Because $D_B^\dagger\psi=\Delta D^\dagger\psi=-(\Delta S_2\psi)'-\Delta S_3\psi$ (with $\Delta S_k=S_k^B-S_k^A$; $S_1$ and all constants cancel), taking the residue at each $\theta_s$ (the residue of a derivative vanishes) gives

\begin{equation}
\sum_m\Delta_m\,N\bigl[(s,i,j),m\bigr]=0,\qquad \Delta=(\Delta H_1,\Delta H_2,\Delta a_1,\Delta a_2,\Delta b_1,\Delta b_2),\label{eq:lagrange}
\end{equation}

\begin{equation}
N[\cdot,H_t]=-\mathop{\mathrm{Res}}_{\theta_s}\frac{\psi P'}{u-\theta_t},\quad N[\cdot,a_t]=\mathop{\mathrm{Res}}_{\theta_s}\frac{\psi P}{(u-\theta_t)^2},\quad N[\cdot,b_t]=\mathop{\mathrm{Res}}_{\theta_s}\frac{\psi P}{u-\theta_t},\label{eq:Ncols}
\end{equation}

$t=1,2$. The six non-scalar Laurent data differences $\Delta_m$ are state-dependent, so no fixed linear combination of brackets works; that is why the spec's search for fixed coefficients could never have succeeded.
The identity holds to $6\cdot10^{-57}$ ($[2,1,0]$, 300 pairs), $5.5\cdot10^{-56}$ ($[3,1,0]$, 600 pairs), $8.4\cdot10^{-57}$ ($[3,2,0]$, 600 pairs).

The stacked $18\times6$ matrix (rows $(s,i,j)$ for the nine pairs $(i,j)$, two sites) has rank $5$ for generic $A\ne B$, never $6$; it has rank $6$ for $A=B$. The $6\times6$ determinant $\det N$, built on the rows $(s,(i,j))$ with $(i,j)\in\{(1,3),(1,2),(2,3)\}$, vanishes for $A\neq B$ at the same twist ($\le1.6\cdot10^{-65}$ over 100 pairs, $2.2\cdot10^{-62}$ for $[3,2,0]$) and is non-zero for $A=B$ (minimum $2\cdot10^{-37}$).
For a vacuum of twist $\omega$, with $R_B(\omega)=\langle\Omega_{\zeta\omega}\otimes\Omega_{\zeta\omega}|\Psi_B\rangle$ for \emph{any} normalisation of $\Psi_B$ and $D_B(\omega)$ the $\det N$ built from the ket's dual $P^j$ and the vacuum's direct $Q_i^{\rm vac,\omega}$,

\begin{equation}
\frac{R_B(\omega)}{R_B(\omega')}=\frac{D_B(\omega)}{D_B(\omega')}\cdot\frac{D_{B_0}(\omega')}{D_{B_0}(\omega)},\qquad B_0=e_1\otimes e_1,\label{eq:ratio}
\end{equation}

on all 64 ($6.4\cdot10^{-55}$), 225 ($6.8\cdot10^{-52}$) and 225 ($1.6\cdot10^{-51}$) states of $[2,1,0]$, $[3,1,0]$, $[3,2,0]$; the identity is $x_0$-free.
In the untwisting regime it holds on the HW branches to $5.4\cdot10^{-67}$, $6.6\cdot10^{-52}$, $2.2\cdot10^{-52}$ at $\epsilon=10^{-6}$ with $|\det N|\sim\epsilon^{12}$ for $(4,2,0),(4,1,1),(3,3,0)$, $\epsilon^{11}$ for the octets and $\epsilon^{10}$ for the singlet in $[2,1,0]$; for $[3,1,0]$, $\epsilon^{16}$ on $(6,2,0),(6,1,1),(5,3,0),(4,4,0)$, $\epsilon^{15}$ on both $(5,2,1)$ and both $(4,3,1)$, $\epsilon^{14}$ on $(4,2,2),(3,3,2)$, and the same pattern transported by duality for $[3,2,0]$. Of the 84 triples of pairs, 72 work, 6 are degenerate (all $i$ or all $j$ equal) and the 6 permutation triples fail; the XXX triple works.
Controls (a$_2$ column replaced by a triple pole, swapped orientation, wrong gauge $\Pi^{-2},\Pi^{-4}$) fail at deviation $\sim1$ to $2$.\footnote{Agents P1 (\texttt{RESULTS-\allowbreak{}Qside.\allowbreak{}md} Q3), P3 (\texttt{RESULTS-\allowbreak{}Qside-\allowbreak{}310.\allowbreak{}md} Q3), P4 (\texttt{RESULTS-\allowbreak{}Qside-\allowbreak{}320.\allowbreak{}md} Q3).} Why the same six equal-twist columns also close the mixed-twist overlaps, with no $\Delta S_1,\Delta e_2,\Delta e_3$ columns, was observed and not derived.

\subsection{The rank-4 rule}\label{sec:rank4}

About 20\% of different-sector pairs have rank $4$ instead of $5$. The rule is clean: the rank is a function of $\Delta\mu=\mu_B-\mu_A$ only, and \textbf{rank $4$ if and only if $\Delta\mu$ has exactly one positive and two negative components} (all $4032$ ordered pairs of $[2,1,0]$: $3318$ of rank 5, $714$ of rank 4; $133$ rank-4 pairs among $600$ for $[3,1,0]$, $130$ for $[3,2,0]$).
The mechanism: for every $k$ with $\Delta\mu_k\le-1$ the product $\psi_k^AP_B^k$ is a rational function, its exponentials cancelling, and it decays at infinity, so its residues at $\theta_1$ and $\theta_2$ sum to zero: row$(1,k,k)+$row$(2,k,k)=0$ in all six columns (violation $\le5.4\cdot10^{-58}$).
Each site's nine rows span a 3-dimensional space $S_s$, the rank is $6-\dim(S_1\cap S_2)$, each negative component of $\Delta\mu$ adds one vector to $S_1\cap S_2$, and two negative components give rank $4$. The second null vector has no closed form (open, minor).\footnote{Agent P3, \texttt{RESULTS-\allowbreak{}Qside-\allowbreak{}310.\allowbreak{}md} section R.}

\subsection{Vacuum Q-functions and vacuum Laurent data}\label{sec:vacQ}

The vacuum covector is itself an eigencovector, with Q-functions of twist $\omega$: direct $Q_i^{\rm vac}=e^{-\omega_iu}q_i$ of degrees $2\lambda_1-\mu^{\rm vac}_i$ as in Table~\ref{tab:degrees}, with $q_1=1$ (the vacuum has $\mu^{\rm vac}=(2\lambda_1,2\lambda_2,0)$).
Its Laurent data are known in closed form and independent of $\zeta$:

\begin{equation}
H_s^{\rm vac}=\sum_i(|\Lambda|-\lambda_i)\,\omega_i\ \pm\ \frac{2e_2(\Lambda)}{\theta_{12}},\qquad a_s^{\rm vac}=c_2\,\omega_3,\qquad b_s^{\rm vac}=\omega_3\bigl(H_s^{\rm vac}-|\Lambda|\omega_3\bigr),\label{eq:vacdata}
\end{equation}

with the upper sign for $s=1$, $e_2(\Lambda)=\lambda_1\lambda_2+\lambda_1\lambda_3+\lambda_2\lambda_3$. The middle statement is where $|\Lambda|$ and $c_2$ were conflated before $[3,2,0]$: $a^{\rm vac}_s=8\omega_3$ for $[3,2,0]$, not $5\omega_3$. The consistency check is $\kappa^{\rm vac}=a^{\rm vac}/c_2=\omega_3$, the exponent of $P^3$. The vacuum eigenvalue functions for the three representations are tabulated in Appendix~\ref{app:conv}.\footnote{Agents P1/P3/P4 (exact eigenvalue functions from the exact rational vacua, \texttt{RESULTS-\allowbreak{}Qside*.\allowbreak{}md} section 0), W2 task C (closed forms, $10$ representations) and W3/P4 (correction for $[3,2,0]$: $a_s^{\rm vac}=8\omega_3$ verified at $\omega_3=3/5,9/7$: $24/5$, $72/7$).}

\section{Which operators localise: the answer to "which \texorpdfstring{$T$}{T}'s?"}\label{sec:localise}

\textbf{Punchline.} In Gaudin, every Laurent coefficient of $S_2,S_3$ at $\theta_1$ localises the two-vacua covector on site 1 by itself: $\{1,H_1,S3c[1,2]\}$ each separately for every representation, plus one extra quadratic combination $X_4$ once $\lambda_1=3$. The number of localisers is $\ell_1=\lambda_1+1$, as in XXX, but the fused shapes of XXX are not needed.

\subsection{The question}\label{sec:Tq}

Paul's question at the start of the project was ``I don't know which $T$'s to use''. In XXX the answer was specific: the shapes $T_{(1,1)}(\theta_1-2h)$ and $T_{(2,1)}(\theta_1-2h)$ (for the adjoint), and in general $T_{(k,k)}$ ($k\le d$) and $T_{(k,d)}$ ($d<k\le\lambda_1$) at $\theta_1-\lambda_1h$ with $d=\lambda_1-\lambda_2$, which acting on $\langle\Omega\otimes\Omega|$ produce $f_\lambda\langle\Omega T^{(1)}_\lambda|\otimes\langle\Omega|$: a covector on site 1 times the unchanged vacuum on site 2.
The set of such site-1 covectors is the ``localising span'' $L_1$, and the CG data are read off from $L_1$ as soon as it reaches every weight with rank equal to the multiplicity.

In Gaudin there are no shifts $\theta_s-mh$; the commuting family is the Laurent coefficients themselves.
We test every operator $X$ in the algebra they generate: $X$ \emph{localises} if $\langle\Omega\otimes\Omega|X=v\otimes\Omega_1$ for a covector $v$ on site 1 (the rank-1 scan: $\langle\Omega\otimes\Omega|X$ viewed as an $\dim V\times\dim V$ matrix has rank $1$ and second factor $\Omega$).

\subsection{The scan}\label{sec:scan}

Let $S3c[s,k]$ denote the coefficient of $(u-\theta_s)^{-k}$ in $S_3$ (so $S3c[s,2]=a_s$ and $S3c[s,1]=b_s$), and

\begin{equation}
D_1=\mathrm{span}\{1,\ W,\ H_1,\ S3c[1,2],\ S3c[1,1],\ S3c[2,2],\ S3c[2,1]\},\qquad D_2=D_1+\text{quadratic monomials},\label{eq:D12}
\end{equation}

with $\dim D_1=7$ (independent basis), $\dim D_2=24$ for $[2,1,0]$ and $26$ for $[3,1,0]$, $[3,2,0]$; the cubic span has rank $48,68,68$ (not $64$, hence not vacuous).
The result, exact at $P$ and $P'$, is as follows.

\textbf{Adjoint.} $\ell_1=\ell_2=3$ already inside $D_1$, with no products needed.
Site 1: $\{1,H_1,S3c[1,2]\}$ localise, each separately,

\begin{equation}
\langle\Omega\otimes\Omega|H_1=\bigl[\Omega_1(c_H-C_\zeta)\bigr]\otimes\Omega_1,\qquad \langle\Omega\otimes\Omega|S3c[1,2]=\bigl[\Omega_1S3^{\rm 1\text{-}site}[-2]\bigr]\otimes\Omega_1,\label{eq:localise}
\end{equation}

where $C_\zeta=\sum_k\zeta_kE_{kk}$, $S3^{\rm 1\text{-}site}[-2]$ is the $(u-\theta_1)^{-2}$ coefficient of the one-site cdet with twist $\mathrm{diag}\,\zeta$, and $c_H$ is a constant (below).
Hence the localising span is $L_1=\Omega_1\cdot\mathrm{span}\{1,C_\zeta,S3^{\rm 1\text{-}site}[-2]\}$, of full rank on every weight of $V$. Site 2: $\{1,\ W+H_1(=-H_2+\text{scalar}),\ S3c[2,2]\}$ with $+C_\zeta$.\footnote{Agent W1, \texttt{Gaudin-\allowbreak{}L2-\allowbreak{}chain.\allowbreak{}wb}, task G3 (cells ``Which T's localise'' and the summary cell).
The prediction by agent P2 (\texttt{RESULTS-\allowbreak{}xxx-\allowbreak{}limit.\allowbreak{}md} X3) from the XXX limit was $T_{(1,1)}(\theta_1-2h)=3h^2\theta_{12}^2[1+h\,X_1/\ldots]$ with $X_1=-\tfrac{4631}{124}\mathbf 1+3H_1-S3c[1,2]$ exactly (the constant at one numerical point only) and $\langle\Omega\otimes\Omega|X_1=u\otimes\Omega_1$, $u=\Omega_1(\tfrac{431}{62}-3C_\zeta-S3^{\rm 1\text{-}site}[-2])$; W1 found this to be a combination of the three separate localisers.}

\textbf{$[3,1,0]$ and $[3,2,0]$.} $\ell_1=4=\lambda_1+1$: three over $D_1$ (the same three), four over $D_2$ (and over $D_1+H_1D_1$ at site 1).
No single quadratic monomial supplies the fourth; the minimal supports are $\{H_1^2,H_1S3c[1,2]\}$ and $\{H_1^2,S3c[1,2]^2\}$, and the localiser is

\begin{equation}
X_4=H_1^2-\tfrac12H_1\,S3c[1,2]-\frac{\gamma}{\theta_{12}}\,W-S3c[1,1]\quad(\mathrm{mod}\ 1,H_1,S3c[1,2]),\label{eq:X4}
\end{equation}

where $\gamma=2$ for $[3,1,0]$ and $\gamma=1$ for $[3,2,0]$, and with $\zeta$- and $\omega$-independent coefficients (checked at 9 values of $\theta_2$ at $P$ and 2 at $P'$ for $[3,2,0]$).
Its covector $v_4=\Omega_1(c_0+c_1C_\zeta+c_2S3[-2]-S3[-1]+C_\zeta^2+\tfrac12C_\zeta S3[-2])$ is $\theta_{12}$-independent.
The coefficient $\gamma$ changed from $2$ to $1$; it could be $\lambda_1-\lambda_2$ or $2/\lambda_2$, and two data points cannot decide.
The three-dimensional $L_1$ already reaches every weight (multiplicities $\le2$), so the fourth is needed to match the count $\ell=\lambda_1+1$ and not for reach.\footnote{Agents W2 (\texttt{Gaudin-\allowbreak{}L2-\allowbreak{}chain-\allowbreak{}2.\allowbreak{}wb} task B; $X_4$ for $[3,1,0]$), W3 (\texttt{Gaudin-\allowbreak{}L2-\allowbreak{}chain-\allowbreak{}320.\allowbreak{}wb} summary: $\gamma=1$, ``not $c_\Lambda$'').
The XXX analogue is presumably the third fused shape $T_{(3,2)}$; the identification is a guess from degrees.}

The quadratic subspace has exact relations that make $D_2$ smaller than $28$ by two.
For $[3,1,0]$ they are $S3c[s,1]=\tfrac14S3c[s,2](H_s-S3c[s,2])$ for $s=1,2$. For $[3,2,0]$ there is no scalar $c_\Lambda$ (none of $1/4,1/3,1/5,1/6,1/8,1/(\lambda_1+1),1/|\Lambda|$) and the relation is an inhomogeneous quadratic: with $a=S3c[1,2]$, $s=\sum\zeta$,

\begin{equation}
S3c[1,1]=3H_1^2-\tfrac74H_1a+\tfrac14a^2-\frac3{\theta_{12}}W-\beta H_1-\gamma'a-\delta,\label{eq:D2rel}
\end{equation}

$\beta=78/\theta_{12}+15s$, $\gamma'=-45/(2\theta_{12})-\tfrac92s$, $\delta=-504/\theta_{12}^2-210s/\theta_{12}-18s^2-3e_2(\zeta)$, an exact $225\times225$ operator identity at six combined $(\zeta,\theta)$ points; site 2 is the mirror $\theta_{12}\to-\theta_{12}$, $H_1\to H_2=-H_1-W+10s$. Its quadratic part factorises, $12H^2-7Ha+a^2=(3H-a)(4H-a)$.\footnote{Agent W3, \texttt{Gaudin-\allowbreak{}L2-\allowbreak{}chain-\allowbreak{}320.\allowbreak{}wb} cells 63--66 and the summary table (verdict ``a different relation'' from the plan's $c_\Lambda$).}

\subsection{Spectators and closed forms}\label{sec:spect}

Two things accompany each localiser: the constant $c_H$ in (\ref{eq:localise}) and the \emph{spectator} factor $f(X)$, defined by $\langle\Omega\otimes\Omega|X|e_1\otimes e_1\rangle$ with $v$ normalised to first component 1 (the analogue of $f_\lambda$ in XXX).
All of them have closed forms:

\begin{equation}
c_H=|\Lambda|\sigma+\frac{2e_2(\Lambda)}{\theta_{12}},\qquad f(H_1)=|\Lambda|\sigma-\zeta\cdot\Lambda+\frac{2e_2(\Lambda)}{\theta_{12}},\qquad f(S3c[s,2])=c_2\,\zeta_3,\label{eq:spect}
\end{equation}

with $f(H_2)$ having the opposite sign of the $1/\theta_{12}$ term, and $e_2(\Lambda)=\lambda_1\lambda_2+\lambda_1\lambda_3+\lambda_2\lambda_3$. The constant $c_H$ has a one-line origin: contracting the $\Omega$'s one has $\sum_{ik}(\Omega E_{ik})\otimes(\Omega E_{ki})=(\sum_i\lambda_i^2)\,\Omega\otimes\Omega$ (the Casimir-type identity on the vacuum), and with $|\Lambda|^2-\sum\lambda_i^2=2e_2(\Lambda)$ the $1/\theta_{12}$ coefficient follows.
We initially guessed $2\lambda_1$ for it (from the two data points $4,6$ of $[2,1,0],[3,1,0]$) and were wrong: it is $2e_2(\Lambda)$, equal to $2\lambda_1\lambda_2=4,6,12$ for $[2,1,0],[3,1,0],[3,2,0]$. Table~\ref{tab:spect} lists spectators at the two standard points.\footnote{$c_H$: agent W2 task C, ten representations (the conjecture $2\lambda_1$ refuted); W3 cell 69.
Spectators: W1, W2, W3 (Laurent polynomials in $\theta_{12}$ fitted on six values with three checks).}

The last formula in (\ref{eq:spect}) is the place where $c_2=\lambda_2(\lambda_1+1)$ and not $|\Lambda|$ appears: $f(S3c[s,2])=8\zeta_3$ for $[3,2,0]$, not $5\zeta_3$; the same correction as for the vacuum's $a^{\rm vac}_s=c_2\omega_3$ (Section~\ref{sec:vacQ}), and for the same reason, since $\langle\Omega\otimes\Omega|S3c[1,2]=\Omega_1S3^{\rm 1\text{-}site}[-2]\otimes\Omega_1$ and the one-site double pole is $c_2\zeta_3$.

{\small
\begin{longtable}{lcccc}
\caption{Spectator factors at the two standard twist points ($P=(3/2,5/4,2/3)$, $P'=(2,5/3,-1/2)$).
Sources: W1, W2, W3 summary cells.}\label{tab:spect}\\
\toprule
$\Lambda$ & $f(H_1)$ at $P$ & $f(H_1)$ at $P'$ & $f(S3c[1,2])$ at $P$ & at $P'$ \\
\midrule
\endfirsthead
\toprule
$\Lambda$ & $f(H_1)$ at $P$ & $f(H_1)$ at $P'$ & $f(S3c[1,2])$ at $P$ & at $P'$ \\
\midrule
\endhead
\bottomrule
\endlastfoot
$[2,1,0]$ & $6+4/\theta_{12}$ & $\tfrac{23}6+4/\theta_{12}$ & $2$ & $-\tfrac32$ \\
$[3,1,0]$ & $\tfrac{95}{12}+6/\theta_{12}$ & $5+6/\theta_{12}$ & $\tfrac83$ & $-2$ \\
$[3,2,0]$ & $\tfrac{121}{12}+12/\theta_{12}$ & $\tfrac{13}2+12/\theta_{12}$ & $\tfrac{16}3$ & $-4$ \\
\end{longtable}}


\subsection{Why Gaudin differs from XXX}\label{sec:whydiff}

In XXX only the specific fused shapes $T_{(1,1)},T_{(2,1)},\dots$ localise; in Gaudin every Laurent coefficient in $D_1$ at $\theta_1$ localises on its own.
This was the round-1 surprise, and we have no derivation of it.
Our reading, which has not been tested, is that in the $h\to0$ limit the shifted transfer matrices at $\theta_s-mh$ collapse onto jets at $\theta_s$, so that the products needed for a fused shape reduce to combinations of Laurent coefficients, and the spectators become Laurent polynomials in $\theta_{12}$ with point-dependent constants instead of ladder products.

\section{The two-vacua family and the untwist: the chain side}\label{sec:chain}

\textbf{Punchline.} The covector $\Omega_{\zeta\omega}\otimes\Omega_{\zeta\omega'}$ with $\omega\neq\omega'$, or its two-parameter $(\infty,\infty)$ sub-family, is CG-complete: after an exact undo it becomes $\langle\varphi_u\otimes e_1|$ with $\epsilon$-free Gauss parameters, it has $K/S=O(\epsilon)$ on every HW branch of $[2,1,0]$, $[3,1,0]$, $[3,2,0]$, and every CG coefficient is read off at leading order.
No cleaning step is needed anywhere.

\subsection{Why a two-vacua family}\label{sec:why2}

With a \emph{single} vacuum and localised transfer matrices the XXX approach had to be cleaned: the localising covector is pure top weight at $\epsilon=0$, its $O(\epsilon)$ lower-weight part meets the $O(\epsilon^{-k})$ raised corrections of the HW-limit eigenstate, and the contamination ratio comes out $K/S=\kappa-1=O(1)$ (for example $\kappa=(t+1)/2$ for $[3,3,0]$).
``Sector-wide cleaning'' removed $K$ for a whole sector by one combination of insertions; ``per-copy cleaning'' was needed where the clean space lacked HW rank ($[4,3,1]$ in $[3,1,0]^{\otimes2}$, $[5,3,2]$ in $[3,2,0]^{\otimes2}$, $[6,4,2]$ in $[4,2,0]^{\otimes2}$).\footnote{\texttt{Paul/\allowbreak{}Mathematica/\allowbreak{}new-\allowbreak{}conventions/\allowbreak{}Clean-\allowbreak{}CG-\allowbreak{}combinations.\allowbreak{}md}; recapitulated in \texttt{Gaudin-\allowbreak{}sl3-\allowbreak{}round3-\allowbreak{}plan.\allowbreak{}md} section 0 (question Q1).} The two-vacua family removed the need already in XXX (Stage D).
Paul asked whether Gaudin avoids cleaning too; the answer is yes, and it is the content of this section.

\subsection{The family, the exact undo and CG-completeness}\label{sec:undo}

The family is $\Phi=\Omega_D[\zeta,\omega]\otimes\Omega_D[\zeta,\omega']$ in general and $\Phi_\infty(\omega_1,\omega_1')$ for $\omega_2,\omega_2'\to\infty$. On every HW ket $h$ of the untwisted algebra the product of two rotated vacua reduces to a rotation on site 1 only, because $\rho(n)\otimes\rho(n)$ with $n$ unipotent acts trivially on HW kets:

\begin{equation}
\langle\Phi|h\rangle=\langle\varphi_u\otimes e_1^T|h\rangle,\qquad \varphi_u=e_1^T\rho(n'),\quad n'=e^{a'E_{12}}e^{g'E_{13}}e^{c'E_{23}}.\label{eq:undo}
\end{equation}

We checked this identity exactly on every HW ket of $[2,1,0]$ (at 11 points), $[3,1,0]$ and $[3,2,0]$ (9 points).
In the regime $\zeta=\epsilon\phi$, $\omega_1=\epsilon\hat\omega_1$, $\omega_1'=\epsilon\hat\omega_1'$ and for the $(\infty,\infty)$ family the Gauss parameters are $\epsilon$-independent,

\begin{equation}
a'=\frac{(\phi_1-\phi_2)(\hat\omega_1-\hat\omega_1')}{(\phi_2-\hat\omega_1)(\phi_2-\hat\omega_1')},\quad
c'=\frac{(\phi_1-\phi_3)(\phi_2-\phi_3)(\hat\omega_1-\hat\omega_1')}{(\phi_1-\phi_2)(\phi_3-\hat\omega_1)(\phi_3-\hat\omega_1')},\quad
g'=\frac{(\phi_3-\phi_1)(\phi_3-\phi_2)(\hat\omega_1-\hat\omega_1')}{(\phi_3-\hat\omega_1)(\phi_2-\hat\omega_1')(\phi_3-\hat\omega_1')},\label{eq:gauss}
\end{equation}

equal to the XXX leading coefficients exactly (the XXX values are the same rational functions).
The Jacobian of the map from the direction parameters to $(a',c',g')$ has rank $3$ for the full family and rank $2$ for $(\infty,\infty)$, and the per-weight ranks of the family of covectors $\varphi_u$ over a grid of $\hat\omega$ equal the weight multiplicities on every weight: the family is CG-complete.
The $\epsilon$-independence of $\Omega_D[\epsilon\phi,\epsilon\hat\omega]$ and of $\varphi_u$ is exact (as exact rationals) and is a gift of Gaudin: in XXX the vacua were $O(1)$ with $O(\epsilon)$ corrections.\footnote{Agents W1 (\texttt{Gaudin-\allowbreak{}L2-\allowbreak{}chain.\allowbreak{}wb} task 4 and summary), W2 (task A: $(\infty,\infty)$ family, closed forms, \texttt{Gaudin-\allowbreak{}L2-\allowbreak{}chain-\allowbreak{}2.\allowbreak{}wb} summary), W3 (\texttt{Gaudin-\allowbreak{}L2-\allowbreak{}chain-\allowbreak{}320.\allowbreak{}wb} A4); also U2, \texttt{RESULTS-\allowbreak{}untwist-\allowbreak{}limit.\allowbreak{}md} ``Chain-side facts''. The $(\infty,\infty)$ vacua span only a 4-dimensional subspace of the 8-dimensional one-site space (D1 \texttt{RESULTS-\allowbreak{}twovac-\allowbreak{}fit.\allowbreak{}md} Surprise 4), yet the family is CG-complete on HW slices.}

\subsection{\texorpdfstring{$K/S$}{K/S} and the CG data}\label{sec:KS}

For each HW-limit branch $\Psi_B=\Psi_0+\epsilon\Psi_1+\dots$ we decompose $\langle\Phi|\Psi_B\rangle=S+K$, with $S=\langle\varphi_u\otimes e_1^T|P_{\rm HW}\Psi_B\rangle$ the ``signal'' and $K$ the contamination from the non-HW part of $\Psi_B$. The result is that $K/S=O(\epsilon)$ on every branch, with $S=O(1)$: the maximum over the 24-point grid at $\epsilon=10^{-6}$ doubles at $2\cdot10^{-6}$ (ratios $2.000$) and $K\equiv0$ on the top branch.
Table~\ref{tab:KS} lists the values; the multiplicity-2 sectors (the ones that needed cleaning in XXX) behave like all others.

{\small
\begin{longtable}{llcc}
\caption{$\max|K/S|$ over the grid at $\epsilon=10^{-6}$ (exponent of $\epsilon$ exactly $1.00$ everywhere, doubling ratio $2.000$).
Full family: 24 grid points; $(\infty,\infty)$: its two-parameter sub-family. $[2,1,0]$ and $[3,1,0]$ from W1/W2; $[3,2,0]$ from W3.
Entries ``--'' were not printed separately.}\label{tab:KS}\\
\toprule
$\Lambda$ & branch $\mu$ & full family & $(\infty,\infty)$ \\
\midrule
\endfirsthead
\toprule
$\Lambda$ & branch $\mu$ & full family & $(\infty,\infty)$ \\
\midrule
\endhead
\bottomrule
\endlastfoot
$[2,1,0]$ & $(4,1,1)$ & $4.9\cdot10^{-7}$ & $\le2.5\cdot10^{-6}$ \\
 & $(3,3,0)$ & $2.4\cdot10^{-6}$ & -- \\
 & $(3,2,1)$ copies 1, 2 & $2.4\cdot10^{-6}$, $1.6\cdot10^{-6}$ & -- \\
 & $(2,2,2)$ & $5.4\cdot10^{-6}$ & -- \\
$[3,1,0]$ & all ten branches & $\le7.7\cdot10^{-6}$ [$(3,3,2)$] & $\le4.8\cdot10^{-6}$ \\
$[3,2,0]$ & $(6,3,1)$ & $2.44\cdot10^{-7}$ & $1.27\cdot10^{-6}$ \\
 & $(6,2,2)$ & $9.76\cdot10^{-7}$ & $5.09\cdot10^{-6}$ \\
 & $(5,5,0)$ & $2.42\cdot10^{-6}$ & $2.42\cdot10^{-6}$ \\
 & $(5,4,1)$ copies 1, 2 & $2.44\cdot10^{-6}$, $1.81\cdot10^{-6}$ & $5.75\cdot10^{-7}$, $2.53\cdot10^{-6}$ \\
 & $(5,3,2)$ copies 1, 2 & $1.79\cdot10^{-6}$, $1.42\cdot10^{-6}$ & $2.39\cdot10^{-6}$, $1.30\cdot10^{-6}$ \\
 & $(4,4,2)$ & $5.53\cdot10^{-6}$ & $1.37\cdot10^{-6}$ \\
 & $(4,3,3)$ & $4.16\cdot10^{-6}$ & $3.65\cdot10^{-6}$ \\
\end{longtable}}


The CG coefficients are extracted from the overlap alone: for each branch the overlaps at several grid points give linear equations on the unknown slice $\langle e_\nu\otimes e_1|\Psi_0\rangle$ (one equation per grid point, using the known $\varphi_u$), solved exactly for one-dimensional slices and by $2\times2$ systems for two-dimensional ones.
The raw error is exactly $O(\epsilon)$ (ratio $4.00$ between $4\cdot10^{-6}$ and $10^{-6}$), and after Richardson extrapolation in $\epsilon$ the extracted slices equal the exact G-norm-1 HW-ket slices to $10^{-17}$--$10^{-21}$ (Table~\ref{tab:cgprec}); the exact values are listed in Appendix~\ref{app:CG}.
For the adjoint the one-dimensional slices are $-1/\sqrt2,\ -1/\sqrt2,\ 0.1178511302$ for $(4,1,1),(3,3,0),(2,2,2)$ and the octet slices are $(0.04931160001,0.2088872897)$ and $(0.3379867346,-0.07978784486)$.

{\small
\begin{longtable}{lcc}
\caption{Precision of the chain-side CG extraction after Richardson.
Sources: W1 summary (task 4d), W2 (task A), W3 (A4 cleaning verdict).}\label{tab:cgprec}\\
\toprule
$\Lambda$ & full family & $(\infty,\infty)$ \\
\midrule
\endfirsthead
\toprule
$\Lambda$ & full family & $(\infty,\infty)$ \\
\midrule
\endhead
\bottomrule
\endlastfoot
$[2,1,0]$ & $10^{-17}$ to $10^{-19}$ & $10^{-18}$ to $10^{-21}$ \\
$[3,1,0]$ & $10^{-17}$ to $10^{-19}$ & $10^{-18}$ to $10^{-21}$ \\
$[3,2,0]$ & $\le1.5\cdot10^{-18}$ (chain Richardson $\le1.6\cdot10^{-20}$) & $\le1.9\cdot10^{-18}$; $(4,4,2)$: $1.3\cdot10^{-17}$ \\
\end{longtable}}


\subsection{Multiplicity-2 sectors and universal mixing}\label{sec:mixing}

In a multiplicity-2 sector the HW space is two-dimensional and the untwisted copies of the HW ket are fixed by the eigenvectors of $S3c[1,2]=a_1$ restricted to it.
We expected this to depend on $\theta$. It does not: $\theta_{12}\,S3c[1,2]|_{\rm HW}$ is a $\theta$-independent matrix with no scalar part (exact at $\theta=(2/7,-3/5),(1/5,-4/9),(3/7,-1/4)$), so the copies, and the mixing, are universal.
For each sector we record the matrix $N_0=\theta_{12}\,S3c[1,2]|_{\rm HW}$ and the coordinate-ratio quadratic of its eigenvectors:

{\small
\begin{longtable}{lcccc}
\caption{Universality of the multiplicity-2 mixing. $N_0=\theta_{12}S3c[1,2]|_{\rm HW}$ on the HW space (coordinates in the canonical rational HW basis of the data files), its trace and determinant, the eigenvalues $\alpha_\pm=(\mathrm{tr}\mp\sqrt{\mathrm{disc}})/2$ with $\mathrm{disc}=\mathrm{tr}^2-4\det$, and the quadratic satisfied by the ratio $r$ of the coordinates of the two eigenvectors.
Sources: W2 task C, W3 cells 33, 35, 37; the adjoint matrix at $\theta_{12}=31/35$ is $(105/62)[[3,-1],[-1,7]]$.}\label{tab:univ}\\
\toprule
rep, sector & $N_0$ & tr, det & eigenvalues & $r$-quadratic \\
\midrule
\endfirsthead
\toprule
rep, sector & $N_0$ & tr, det & eigenvalues & $r$-quadratic \\
\midrule
\endhead
\bottomrule
\endlastfoot
$[2,1,0]$ $(3,2,1)$ & $\tfrac32\begin{pmatrix}3&-1\\-1&7\end{pmatrix}$ & $15,\ 45$ & $(15\mp3\sqrt5)/2$ & $r^2+4r-1=0$ \\
$[3,1,0]$ $(5,2,1)$ & $\tfrac13\begin{pmatrix}19&-40\\-2&44\end{pmatrix}$ & $21,\ 84$ & $(21\mp\sqrt{105})/2$ & $40r^2+25r-2=0$ \\
$[3,1,0]$ $(4,3,1)$ & $\tfrac43\begin{pmatrix}4&-1\\-2&14\end{pmatrix}$ & $24,\ 96$ & $12\mp4\sqrt3$ & $r^2+10r-2=0$ \\
$[3,2,0]$ $(5,4,1)$ & $\begin{pmatrix}13&-6\\-1&22\end{pmatrix}$ & $35,\ 280$ & $(35\mp\sqrt{105})/2$ & $6r^2+9r-1=0$ \\
$[3,2,0]$ $(5,3,2)$ & $\tfrac43\begin{pmatrix}16&-1\\-2&26\end{pmatrix}$ & $56,\ 736$ & $28\mp4\sqrt3$ & $r^2+10r-2=0$ \\
\end{longtable}}


The octet mixing parameter $R=(e_1E_{13}\cdot C)/(e_1E_{23}E_{12}\cdot C)$ on the weight-$(1,1,1)$ slice satisfies $R^2-5R-5=0$ for the adjoint, with roots $5.854101966250$ and $-0.854101966250$; the Gaudin-extracted values agree to $1.5\cdot10^{-18}$ and $5.7\cdot10^{-20}$. For $[3,1,0]$ $(5,2,1)$ the roots are $4.311737691490$ and $-0.811737691490$ (agreement $5\cdot10^{-19}$), and exactly the same numbers appear for $[3,2,0]$ $(5,4,1)$, as duality requires.
Notice that $(5,3,2)$, the sector that needed per-copy cleaning in XXX, has literally the same quadratic as $[3,1,0]$'s $(4,3,1)$, namely $r^2+10r-2=0$. The XXX values at finite $t=\theta_{12}/h$ were different (for example $-1.694$ and $1.096$ for the adjoint octet); the comparison with the $t\to\infty$ limit of those numbers was proposed but not carried out.
For $[3,1,0]$ agent W1 did not compute the mixing in the $(4,3,1)$ sector (the functionals $e_1E_{13}$ and $e_1E_{23}E_{12}$ do not suffice there); the $[3,2,0]$ extractions of W3 solve the two-dimensional slices from pairs of grid points.

\subsection{Duality \texorpdfstring{$[3,2,0]\cong[3,1,0]^*\otimes\det^3$}{[3,2,0][3,1,0]*det}}\label{sec:duality}

The HW sectors of $[3,2,0]^{\otimes2}$ are exactly the dual images $\mu^*=(6-\mu_3,6-\mu_2,6-\mu_1)$ of those of $[3,1,0]^{\otimes2}$, with the same multiplicities: $(6,2,0)\to(6,4,0)$, $(6,1,1)\to(5,5,0)$, $(5,3,0)\to(6,3,1)$, $(5,2,1)\times2\to(5,4,1)\times2$, $(4,4,0)\to(6,2,2)$, $(4,3,1)\times2\to(5,3,2)\times2$, $(4,2,2)\to(4,4,2)$, $(3,3,2)\to(4,3,3)$. The eigenvalues of $\theta_{12}S3c[1,2]$ on a dual sector differ from those of the original by a \emph{sector-dependent scalar} while the traceless parts agree: $(4,3,1)\to(5,3,2)$: $+16$ ($96-24x+x^2\to736-56x+x^2$); $(5,2,1)\to(5,4,1)$: $+7$ ($84-21x+x^2\to280-35x+x^2$); multiplicity one: $(3,3,2)\to(4,3,3)$: $20\to36$, $(4,2,2)\to(4,4,2)$: $18\to30$, $(4,4,0)\to(6,2,2)$: $0\to24$, $(5,3,0)\to(6,3,1)$: $0\to16$, $(6,1,1)\to(5,5,0)$: $8\to0$, $(6,2,0)\to(6,4,0)$: $0\to0$. The scalar shift is new.
The G-orthonormal slice ratios agree for $(4,3,1)\leftrightarrow(5,3,2)$ ($-3.6049$ and $0.06935$ in both) but not for $(5,2,1)\leftrightarrow(5,4,1)$, because the GT bases of $V$ and $V^*$ differ by a rotation in a two-dimensional weight space; the basis-independent mixing $R$ above agrees.
The exponents of $\det N$, $|R(\omega)|$ and the $\epsilon$-orders transport by duality too ($\epsilon^{16,15,14}$).\footnote{Agents W3 (\texttt{Gaudin-\allowbreak{}L2-\allowbreak{}chain-\allowbreak{}320.\allowbreak{}wb}, duality row of the summary table) and P4 (\texttt{RESULTS-\allowbreak{}Qside-\allowbreak{}320.\allowbreak{}md} Q4).}

\subsection{The controls and the ``no cleaning'' answer}\label{sec:controls}

Three controls test what the family needs.

\begin{itemize}
\item $\omega_1'=\omega_1$ (common rotation).
Then $a'=g'=0$, the family reads only the top branches: $(4,2,0),(4,1,1)$ for $[2,1,0]$ (full family), $(6,4,0),(6,3,1),(6,2,2)$ for $[3,2,0]$, and only the top branch in the $(\infty,\infty)$ family.
The others are ``not read''. This is as in XXX.
\item $\omega_1'$ held fixed at a finite value, not untwisted.
In XXX this broke the untwisting.
In Gaudin with \emph{finite} $\omega_2$ it does not: $\Omega_D[\zeta\to0,\omega']$ has a finite direction-dependent limit, $K/S$ is still $O(\epsilon)$ on every branch, and the per-weight ranks over six directions are full.
We expected the opposite and recorded it as a surprise.
\item $\omega_1'$ held fixed in the $(\infty,\infty)$ family (at $\omega_2=\infty$).
This \emph{fails}: $\varphi_u$ has rank $\le1$ on the multiplicity-2 weights, singlet-type branches have $K/S=O(1)$ (for $[3,2,0]$: $0.465$ at $(4,4,2)$ and $0.388$ at $(4,3,3)$, identical to $[3,1,0]$'s duals $(4,2,2)$ and $(3,3,2)$), and all branches with a 2-dimensional slice have $O(1)$ CG errors ($0.39$ to $1.5$; $\sigma_2=1.5\cdot10^{-7},1.7\cdot10^{-8},8.7\cdot10^{-15}$).
Branches with one-dimensional slices still pass.
\end{itemize}

So the answer to Paul's question is: no cleaning is needed anywhere, in either family, for the adjoint, $[3,1,0]$ or $[3,2,0]$. Specifically $(5,3,2)$ has $\max|K/S|=1.8\cdot10^{-6},1.4\cdot10^{-6}$ (full family) and $2.4\cdot10^{-6},1.3\cdot10^{-6}$ ($(\infty,\infty)$) at $\epsilon=10^{-6}$, with the same doubling ratio as every other branch.
What is untested is $[4,2,0]$ (dimension 27, multiplicity-3 sector $(6,4,2)$), the other XXX trouble spot.
The only weakly conditioned $2\times2$ extraction is $(4,4,2)$ in the $(\infty,\infty)$ family ($\sigma_2\sim3\cdot10^{-4}$), still accurate to $1.4\cdot10^{-17}$ at 200 digits; conditioning is a property of the chosen pair of grid points (best pairs have $\sigma_2\approx0.4$ to $2.8$), not of the family.\footnote{Agent W3, \texttt{Gaudin-\allowbreak{}L2-\allowbreak{}chain-\allowbreak{}320.\allowbreak{}wb}, summary ``Cleaning verdict (A4)'', cells 41--57.
The exact slice of $(4,4,2)$ has an exactly vanishing component and the sign convention of the HW kets differs between the exact ket and the chain branch there, so the CG cells align signs and use norm-based errors.}

\section{The jet formula: the overlap as a bilinear form in jets of the Q-functions}\label{sec:jet}

\textbf{Punchline.} For any eigenstate, the normalisation-free overlap $c\,\langle\Phi_\infty|\Psi\rangle/\langle x_0^G|\Psi\rangle$ times $M(\theta_1)M(\theta_2)$ equals a bilinear form in the monomials $P^3p_{kl}$ at the two sites, with a state-independent kernel that is unique, rational, and polynomial in $\omega_1,\omega_1',1/\theta_{12}$. The number of monomials per site is $\binom{\lambda_1+2}{2}-\binom{\lambda_2+1}{2}$, namely $5,9,7$ for $[2,1,0],[3,1,0],[3,2,0]$.

\subsection{The statement}\label{sec:jetstate}

We write $Q_{\rm num}(B)=\langle\Phi_\infty|\Psi_B\rangle/\langle x_0^G|\Psi_B\rangle$; it is independent of the normalisation of $\Psi_B$. Equation~(\ref{eq:jet0}) of the overview says, with $m_{kl}=P^3p_{kl}$ evaluated at $\theta_s$ on the dual Q-functions of the state,

\begin{equation}
c\,Q_{\rm num}\,M(\theta_1)M(\theta_2)=\sum_{(kl),(k'l')}K_{(kl),(k'l')}\,m_{kl}(\theta_1)\,m_{k'l'}(\theta_2),\qquad M=P^3\,p_{0,\lambda_2+1}.\label{eq:jet}
\end{equation}

Several features are worth spelling out.

\begin{itemize}
\item \emph{Which pairs.} $(kl)$ runs over $0\le k<l\le\lambda_1+1$ with not both $k,l\le\lambda_2$. The excluded pairs are those forced to vanish by $P^{(k)}=\kappa_{k,s}P$ for $k\le\lambda_2$ (Section~\ref{sec:frob}).
Their number is $\binom{\lambda_1+2}2-\binom{\lambda_2+1}2$.
\item \emph{The lone factor is always $P^3$.} The monomials $P^{3(c)}p_{kl}$ with $c>0$ are redundant: $P^{3\prime}p_{0l}=P^3p_{1l}$ and $P^{3\prime\prime}p_{0l}=P^3p_{2l}$ exactly (Frobenius).
\item \emph{$M$ is the first non-forced coordinate.} $M=P^3p_{0,\lambda_2+1}\propto\Pi'(\theta_s)Q_1(\theta_s)(\zeta_3-\zeta_2)e^{\sigma\theta_s}P^3$ by the Wronskian relation (adjoint), and it is the vacuum's own monomial.
\item \emph{Normalisation $c$.} The frame-tie scalar (\ref{eq:cfactor}). $\zeta$ enters the formula only through $c$ and the Q-functions; the kernel is $\zeta$-independent.
\item \emph{The kernel.} Unique (the design matrix of products has full rank), rational, and homogeneous.
For the adjoint $K=G/4$ with $G$ the $5\times5$ matrix of Table~\ref{tab:G210}; in terms of the symbols used in the notes there are two equivalent normalisations ($G=4K$, because the vacuum monomial is $M_{\rm vac}=2h^2W'P^3$ per site, and $\Theta=K/(\omega_1\omega_1')^{2\lambda_1-\lambda_2}$ in the companion frame), reconciled in Appendix~\ref{app:conv}.
\end{itemize}

The formula was found, fitted, and derived.
We describe each route.

\subsection{The adjoint in detail}\label{sec:adjoint}

For $[2,1,0]$ the pairs are $(02),(03),(12),(13),(23)$, and since $p_{12}=\kappa p_{02}$, $p_{13}=\kappa p_{03}$ on every state, a basis of five one-site monomials is

\begin{equation}
N=\bigl(P^3p_{02},\ P^{3\prime}p_{02},\ P^3p_{03},\ P^{3\prime}p_{03},\ P^3p_{23}\bigr),\qquad\text{weights }(2,3,3,4,5),\label{eq:N5}
\end{equation}

(the order of D2; $P^{3\prime}p_{0l}=P^3p_{1l}$).
The kernel $G_{ij}(\omega_1,\omega_1',\theta_{12})$ is a $5\times5$ matrix of polynomials listed in full in Table~\ref{tab:G210}; a few entries show the flavour, with $t=1/\theta_{12}$:

\begin{equation}
G_{11}=-4\omega_1^2\omega_1'^2\,(3t^2+t\omega_1-t\omega_1'-\omega_1\omega_1'),\quad G_{33}=\omega_1^2\omega_1'^2,\quad G_{55}=1,\quad G_{45}=-2(2t+\omega_1).\label{eq:Gsome}
\end{equation}

Properties, all exact: $G$ has rank $5$; it is symmetric under site exchange in the form $G_{ji}(\omega_1,\omega_1',\theta_{12})=G_{ij}(\omega_1',\omega_1,-\theta_{12})$; each entry is homogeneous of degree $10-w_i-w_j$ in $(\omega_1,\omega_1',1/\theta_{12})$ jointly, with $w=(2,3,3,4,5)$ (checked by the scaling $u\to2u$, $\zeta,\omega\to\cdot/2$: the refitted kernel equals $G$ at $(\omega/2,2\theta_{12})$ to $5\cdot10^{-72}$); every entry is a polynomial in $(\omega_1,\omega_1')$ of total degree $\le6$ (no poles, none in $\omega_1-\omega_1'$ or $\omega_1-\zeta_j$) and in $1/\theta_{12}$ of degree $\le4$; it depends on $\theta_1,\theta_2$ only through $\theta_{12}$. Table~\ref{tab:Gtw} gives the same kernel in jets of the twisted functions $e^{-\omega_su}P^j$: in those variables the two vacua interact only through $\delta=\omega_1'-\omega_1$ and $t$, the Gaudin image of the XXX statement that they couple only through $s=a'/a$.

Numerically: on all 64 states the formula holds with ratio exactly $1.0$ and maximal relative deviation $2.8\cdot10^{-56}$ at $(\omega_1,\omega_1')=(5/2,7/4)$, and $2\cdot10^{-56}$ to $3\cdot10^{-55}$ at four further parameter points (other $(\omega_1,\omega_1')$, other $\theta$, other $\zeta$).\footnote{Agent D2, \texttt{RESULTS-\allowbreak{}twovac-\allowbreak{}limit.\allowbreak{}md} L3 (60 digits); the constant between left and right sides is exactly $1$, the one predicted from the normalisations, so no leftover scalar exists.}

\subsection{Route 1: the \texorpdfstring{$h\to0$}{h0} limit of the XXX ten-product formula}\label{sec:route1}

Agent D2 took the XXX identity for the adjoint $(\infty,\infty)$ family,

\begin{equation}
\frac{\langle\Phi_c|\Psi_c\rangle}{\langle x_0^c|\Psi_c\rangle}M^{(1)}_{\rm vac}M^{(2)}_{\rm vac}=\sum_{c,c'\in\{0,1\}}Q^3(\theta_1-ch)Q^3(\theta_2-c'h)\sum_{k=1}^{10}\mathrm{Bk2}^{(1)}[F_k,G_k]\,\mathrm{Bk2}^{(2)}[F'_k,G'_k],\label{eq:xxx10}
\end{equation}

with ladder points $\theta_s+mh$, $m=-3,\dots,0$, $2\times2$ bracket determinants $\mathrm{Bk2}$ of the dual Q-functions, $M^{(s)}_{\rm vac}=\mathrm{minor}_s(-3,-2)Q^3(\theta_s-h)$, and kernels rational in $s=a'/a$ and $t=\theta_{12}/h$ (parsed from cell 62 of \texttt{L2-\allowbreak{}210-\allowbreak{}two-\allowbreak{}vacua-\allowbreak{}brackets.\allowbreak{}wb}; the 40 kernel laws and ten products reproduce the printed values with 0 mismatches), and expanded it formally at $h\to0$ in the diagonal frame.
The accounting is as follows.

The XXX normalisation $\kappa\kappa'=h^{-8}/c$ and $\langle x_0^d|\psi\rangle=h^2\langle x_0^G|\psi\rangle$ exactly (since $x_0$ is exactly linear in $z$) give $Q_{\rm num}^{\rm XXX}\sim h^6\,c\,Q_{\rm num}^G$. Each minor is $O(h^2)$ (the $O(h)$ term is $-W(\theta_s)\cdot(m'-m)$ and vanishes), so the naive order of the right-hand side is $h^4$; the actual leading order is $h^{10}$. The cancellations at orders $h^4,\dots,h^9$ are not accidental: they follow from \emph{one} exact relation.
The dual Baxter equation at $x=\theta_s$ has $\nu_3(\theta_s)=0$ and $t_3(\theta_s)=\chi_3\nu_1\nu_2(x-h)\nu_3(x-2h)|_{x=\theta_s}=0$, so

\begin{equation}
Q^j(\theta_s-h)=\rho_s\,Q^j(\theta_s),\qquad\text{hence}\qquad\mathrm{minor}_s(-1,0)=0,\label{eq:minor10}
\end{equation}

with $\rho_s$ independent of $j$, exactly and to all orders in $h$, for every XXX eigenstate: the XXX generalisation of $W(\theta_s)=0$. With only this relation imposed and all other jets random, the right-hand side has valuation exactly $10$ (three random trials).
Consequently the four ladder values of $(Q^2,Q^3)$ at $\theta_s$ lie in a 2-plane, the minors are $D$ times $(-\alpha,\rho\beta,\rho,\beta,1,0)$ with leading values $(\alpha,\beta,\rho)=(-2,3,1)$, and the leading coefficient depends only on the $n=0$ jets $J_{0,0},J_{0,2},J_{0,3}$ (and $J_{0,1}=\kappa J_{0,0}$), through the Plücker coordinates; the $O(h)$ corrections $q_1$ of the XXX Q-functions enter first at order $h^{11}$, so no extra equation for $q_1$ was needed.
Two corrections to the plan came out: the vacuum monomial is $M_{\rm vac}=2h^2W'(\theta_s)P^3(\theta_s)$, not $(5/2)h^2W'P^3$ (the weight-2 minor contains an $n=1$ combination which the exact relation fixes to $p/2$), and only five of the twelve candidate monomials $\{P^{3(c)}p_{kl}\}$ occur.
The result is (\ref{eq:jet}) with the table of $G$.\footnote{Agent D2, \texttt{RESULTS-\allowbreak{}twovac-\allowbreak{}limit.\allowbreak{}md} L0--L2; reproduced exactly for the top state at finite $h=1/2$ ($e_1\otimes e_1$, exact rationals), and independently ($h^6\,c\,Q_{\rm num}$ for the top state converges to $102245/41472$, to $10^{-10}$ by Richardson, against $c\,Q_{\rm num}^G(\mathrm{top})=2.4653983410$).}

A natural shortcut failed, and it is worth recording.
Taking $t\to\infty$ at fixed $a=1/\omega_1$ in the XXX kernels is wrong: of the 40 kernel laws, 25 survive and all equal $s\,u_c[mm']u_{c'}[nn']$, a rank-one $U\otimes U$, and the naive right-hand side divided by the vacuum monomials equals $(1-a)^3(1-a')^3$, state independent.
The Gaudin limit is $a=e^{-h\omega_1}\to1$, not fixed $a$, and all the state dependence sits in the cancellation hierarchy, invisible to a rank-one naive kernel.\footnote{Agent D2, \texttt{RESULTS-\allowbreak{}twovac-\allowbreak{}limit.\allowbreak{}md} L4 (a plan item requested by W2).}

\subsection{Route 2: the direct fit}\label{sec:route2}

Agent D1 fitted the kernel directly from Gaudin data, without reference to XXX. With target $4\,c\,Q_{\rm num}\,(P^3p_{02})_1(P^3p_{02})_2$ and the basis $N$ (25 unknowns) on 320 rows ($5$ twist points), the residual is below the fixed-point floor $10^{-90}$, the rank is 25 (all products independent), the kernel is unique, all entries are rational, and it equals D2's $G$ to $2.8\cdot10^{-74}$; per-$\zeta$ fits agree with the stacked one to $10^{-72}$. A hypothesis ladder led to this: Table~\ref{tab:ladder} lists the levels.
The direct-$Q_1$ bases do \emph{not} fit with a $\zeta$-independent kernel because $p_{03}=A_s[(\Lambda_s-\kappa)Q_1+2Q_1']$ with $\Lambda_s=2\sigma+2/\Pi'(\theta_s)$ contains $\sigma=\sum\zeta$; the Plücker form is the $\zeta$-free one.
The controls fail as they should: lone index $P^2$ or $P^1$ instead of $P^3$ gives residuals $0.99$; jets built from site-1 or site-2 alone $0.37$, $0.58$; jets at wrong points $0.48$; no $M$: $0.97$.\footnote{Agent D1, \texttt{RESULTS-\allowbreak{}twovac-\allowbreak{}fit.\allowbreak{}md} F2; fixed-point 320-bit column-pivoted QR, a $512\times400$ problem in 30 s.
Larger bases (144, 576 unknowns) fit with null spaces whose sparse representative is the $N_5$ kernel.} The kernel laws were then fixed on a $3\times3$ grid of $(\omega_1,\omega_1')$ and five $\theta$ variants: each entry is a polynomial of total degree $\le6$ in $(\omega_1,\omega_1')$ with 28 coefficients, identical (exact rational comparison) to D2's analytic $G_{ij}$ in $25/25$ entries.

{\small
\begin{longtable}{llrrr}
\caption{Hypothesis ladder for the adjoint fit (D1 F2): relative residual $\lVert Ak-y\rVert/\lVert y\rVert$, $0.0$ meaning below the $10^{-90}$ floor; 320 rows (5 twist points); $M=P^3p_{02}$ unless stated.}\label{tab:ladder}\\
\toprule
level & basis per site & unknowns & rank & residual \\
\midrule
\endfirsthead
\toprule
level & basis per site & unknowns & rank & residual \\
\midrule
\endhead
\bottomrule
\endlastfoot
a & $N_5=\{P^3p_{02},P^{3\prime}p_{02},P^3p_{03},P^{3\prime}p_{03},P^3p_{23}\}$ & 25 & 25 & \textbf{0.0} \\
a & direct $A_sQ_1^{(k)}P^{3(c)}$, $k\le2,c\le1$ (6) & 36 & 36 & 0.294 \\
a & direct, $k\le1,c\le3$ (8) & 64 & 64 & 0.562 \\
a & D2's nine direct monomials with $\kappa$ & 81 & 81 & 0.363 \\
b & Plücker $\{p_{kl}P^{3(c)}\}$, $k<l\le3,c\le1$ (12) & 144 & 64 & 0.0 \\
c & Plücker $k<l\le3,c\le3$ (24) & 576 & 256 & 0.0 \\
c & direct $k\le3,c\le3$ (16) & 256 & 256 & 0.0643 \\
ctl & $N_5$ without $M_1M_2$ & 25 & 25 & 0.966 \\
\end{longtable}}


\subsection{Route 3: derivation from the jet basis}\label{sec:route3}

The third derivation (Section~\ref{sec:B}) constructs the Gaudin K tensor of $\Phi_c^\infty$ in the jet basis that replaces the SoV basis when $h\to0$, and finds exactly 25 non-zero entries on $N_5\times N_5$, equal to $G/(4\omega_1^3\omega_1'^3)$ as an exact polynomial identity in $(1/\omega_1,1/\omega_1')$. The ``4'' is the jet normalisation $M_{\rm vac}=2h^2P^3p_{02}$ per site, and $c$ arises from the frame tie (\ref{eq:kappaG}) exactly.

\subsection{General \texorpdfstring{$\Lambda$}{}: \texorpdfstring{$[3,1,0]$}{[3,1,0]} and \texorpdfstring{$[3,2,0]$}{[3,2,0]}}\label{sec:jetgen}

The same fit pipeline, now with the sparse hypothesis tested first, gave the generalisation.
For $[3,2,0]$ (dim 15, 225 states) with $c=\prod_\omega(\zeta_2-\omega)(\zeta_3-\omega)^3(\zeta_1-\zeta_2)^4$ and $M=P^3p_{03}$, the target is exactly (residual $0.0$ at 80 digits) a bilinear form in $P^3$ times the seven non-forced coordinates $\{p_{03},p_{04},p_{13},p_{14},p_{23},p_{24},p_{34}\}$ per site; the kernel is unique (49 unknowns, rank 49 on 1125 rows), rational, and equal for every $\zeta$ point to $\sim10^{-72}$. Dropping any one of the seven leaves a residual ($0.41$ without $p_{03}$ or $p_{04}$, $0.16,0.25,0.078,0.14,0.061$ for $p_{13},p_{14},p_{23},p_{24},p_{34}$).
For $[3,1,0]$ ($c=(\zeta_2-\omega)^2(\zeta_3-\omega)^3(\zeta_1-\zeta_2)^2$ per vacuum, $M=P^3p_{02}$), greedy backward elimination from 18 features ends with the pure Plücker basis of the nine non-forced pairs, 81 unknowns, rank 81, residual $0.0$, unique rational kernel, $K_{99}=1/9$; the adjoint-like five-element basis fails ($0.316$).
The pattern is the count

\begin{equation}
\#\text{monomials}=\binom{\lambda_1+2}2-\binom{\lambda_2+1}2=6-1,\ 10-1,\ 10-3=5,\ 9,\ 7.\label{eq:count}
\end{equation}

In the $\nu$-variables of $[3,2,0]$ the seven monomials are $(1,\kappa,\kappa_2)\times(1,\pi_4)$ and $\pi_{34}$ with $\pi_4=p_{04}/p_{03}$, $\pi_{34}=p_{34}/p_{03}$; monomials of $\nu$-weight $\le3$ (13 of them) fail ($0.0046$), weight $\le4$ fit.\footnote{Agent P4, \texttt{RESULTS-\allowbreak{}Qside-\allowbreak{}320.\allowbreak{}md} section J (hypothesis ladder table, $[3,1,0]$ comparison T7), points $P,P',X_1,X_2,X_3$ built at 80 digits; the wrong exponents of $c$ ($(1,3,2)$ for $[3,1,0]$) give $0.32$.} The kernel is neither symmetric at fixed $\theta_{12}$ nor does $K+K^T$ express site exchange; the exchange symmetry is $K_{ij}(\omega,\omega,1/\theta_{12})=K_{ji}(\omega,\omega,-1/\theta_{12})$ (deviation $5.8\cdot10^{-63}$).

\subsection{Kernel laws}\label{sec:laws}

Table~\ref{tab:laws} collects the laws that each kernel obeys.
They were established on grids of $(\omega_1,\omega_1')$ (9 values, 72 ordered pairs, 5 twist points each) and $\theta_{12}$ (12 distinct values) with exact rational polynomial reconstruction, and the structure was verified at held-out points; Appendix~\ref{app:K} gives the data.
The two sites exchange under $\theta_{12}\to-\theta_{12}$, and the first non-forced coordinate carries weight $0$, with the monomial weight $e_{kl}=k+l-\lambda_2-1$.

{\small
\begin{longtable}{p{2.6cm}p{3.4cm}p{4.4cm}p{4.2cm}}
\caption{Laws of the jet kernel for the three representations: $e$ are the monomial weights relative to $M$; the $\omega$-degree of entry $(i,j)$ is $\le$ total $-e_i-e_j$; $t=1/\theta_{12}$. Sources: D1/D2 (adjoint), P4 ($[3,2,0]$), G1 ($[3,1,0]$, \texttt{RESULTS-\allowbreak{}general-\allowbreak{}kernel.\allowbreak{}md} K0).}\label{tab:laws}\\
\toprule
 & $[2,1,0]$ & $[3,1,0]$ & $[3,2,0]$ \\
\midrule
\endfirsthead
\toprule
 & $[2,1,0]$ & $[3,1,0]$ & $[3,2,0]$ \\
\midrule
\endhead
\bottomrule
\endlastfoot
monomials per site & 5 & 9 & 7 \\
pairs $(kl)$ & $(02)$, $(03)$, $(12)$, $(13)$, $(23)$ & $(02)$, $(03)$, $(04)$, $(12)$, $(13)$, $(14)$, $(23)$, $(24)$, $(34)$ & $(03)$, $(04)$, $(13)$, $(14)$, $(23)$, $(24)$, $(34)$ \\
$e=k+l-\lambda_2-1$ & $0,1,1,2,3$ & $0,1,2,1,2,3,3,4,5$ & $0,1,1,2,2,3,4$ \\
max $\omega$-degree & 3 each & 5 each & 4 each \\
$t$-degree & 4 & 8 & 7 \\
$c$ & $(\zeta_2-\omega)(\zeta_3-\omega)^2$ per vacuum, $(\zeta_1-\zeta_2)^2$ & $(\zeta_2-\omega)^2(\zeta_3-\omega)^3$, $(\zeta_1-\zeta_2)^2$ & $(\zeta_2-\omega)(\zeta_3-\omega)^3$, $(\zeta_1-\zeta_2)^4$ \\
last entry $K_{\rm last}$ & $1/4$ & $1/9$ & $1/9$ \\
monomials in $\Theta$ & 25 entries & 998 & 458 \\
\end{longtable}}


\subsection{What went wrong first}\label{sec:wrong}

The first fits all failed, and the failures are worth a paragraph.
With the plan's vacuum monomial $M=Q_1P^3$ and the assumption that the kernel is $\zeta$-independent, no level of the ladder fit, not even the full jet space $\{J^2_aJ^3_bJ^3_c\}$ with 1295 unknowns (training residual $4\cdot10^{-2}$, hold-out error $8\cdot10^2$).
The cause was normalisation: $p_{02}=A_sQ_1$ with $A_s=(\zeta_3-\zeta_2)e^{\sigma\theta_s}\Pi'(\theta_s)$ is $\zeta$-dependent, and so is the frame-tie scalar.
Once one divides by $c$ and uses $M=P^3p_{02}$, the kernel is exactly $\zeta$-free.
Likewise $\langle x_0^G|\Psi_B\rangle/n_B$ turns out to be, to $2\cdot10^{-77}$, a constant times $Q_1(\theta_1)P^3(\theta_1)Q_1(\theta_2)P^3(\theta_2)$ (exponent $1$ on each, $0$ on the other ten numbers), so the ``x$_0$-free'' target $T(B)=\langle\Phi_\infty|\Psi\rangle D_B(\omega_0)/\langle\Omega_0\otimes\Omega_0|\Psi\rangle$ is not constant times $Q_{\rm num}$ but $c_T(\zeta)\,Q_{\rm num}M_1M_2$: the two targets are the same fit with one free scalar per $\zeta$ point.
And the one-site-fixed-$a$ limit is wrong, as said above.
The vacuum monomial of the plan was off by a factor $5/4$ ($M_{\rm vac}=2h^2W'P^3$, not $\tfrac52h^2W'P^3$), and the twelve candidate monomials reduce to five.

\section{The Gaudin \texorpdfstring{$B$}{B} operator and the jet basis}\label{sec:B}

\textbf{Punchline.} There is no single Gaudin $B$ operator with simple spectrum.
The XXX operator $B(u)$ becomes, as $h\to0$, a hierarchy $\nu^3[1+hb_1+h^2b_2+\dots]$ whose first member is Cartan-like and degenerate; the SoV basis becomes a $\zeta$-independent \emph{jet basis}, defined by a wavefunction rule in terms of eight one-site monomials, and the kernel $G$ of the jet formula is exactly the K tensor of the vacuum family in that basis.
Where the degeneracy lives is the diagonal frame, not $B$.

\subsection{The XXX operator and its eigenbasis}\label{sec:Bxxx}

In the new XXX conventions (gl(3), untwisted) $B(u)=T_{11}(u)\,\mathrm{qmin}_{12,12}(u+h)+T_{21}(u)\,\mathrm{qmin}_{12,13}(u+h)$. It has simple spectrum $\prod_s\prod_k(u-\theta_s-m_kh)$ with labels $(a|bc)$ per site and roots $\theta_s+\{-1-a,-b,-c\}h$, $\lambda_2\le a\le\lambda_1$, $0\le c\le\lambda_2$, $c\le b\le a$; for the adjoint there are eight labels per site, $(2,2,1)$ being the vacuum with roots $\{-3,-2,-1\}$. Its left eigenbasis is the SoV basis, in which the wavefunction is $\langle x|\Psi_c\rangle/\langle x_0|\Psi_c\rangle=N_x\prod_sM^{(s)}_x/M^{(s)}_{\rm vac}$ with $M_{(a|bc)}=\mathrm{minor}_s(-1-a,-b)\,Q^3(\theta_s-ch)$. We verified, exactly at $h=1/2$ on the adjoint, that $B$ has degree $6$ in $u$ and in $h$, $[B(u_0),B(u_1)]=0$, all 64 label-rule eigenvalues are simple and the left eigencovectors are $u$-independent.\footnote{Agent D3, \texttt{RESULTS-\allowbreak{}gaudin-\allowbreak{}B.\allowbreak{}md} B0 (Fractions, $u_0=3/7$, $u_1=-2/9$; nnz of $B_0,\dots,B_6$ are $64,64,232,251,247,178,75$).
The XXX wavefunction check with dual Q's at $h=1/2$ was skipped as allowed; the rule is not used by the derivation, and the formula built on it is reproduced independently.}

As $h\to0$ all roots collapse to $\theta_s$. So the ``Gaudin $B$'' cannot be one operator with simple spectrum; it has to be a hierarchy whose joint generalised eigenbasis is a jet basis, the $h$-expansion coefficients (divided differences) of the SoV covectors.

\subsection{The hierarchy \texorpdfstring{$b_1,b_2,\dots$}{b,b,}}\label{sec:bk}

With $t(x)=T(x)/\nu(x)=1+hL(x)+h^2C(x)$, $\nu=\Pi$, $L_{ij}=\sum_sE^{s}_{ij}/(x-\theta_s)$ the Gaudin Lax matrix and $C_{ij}=\sum_kE^1_{ik}E^2_{kj}/((x-\theta_1)(x-\theta_2))$ the cross-site term, the exact all-order structure is

\begin{equation}
\begin{aligned}
\frac B{\nu^3}=(1+hs_1+h^2s_2)\Bigl\{&t_{11}(u)\bigl[t_{11}(u+h)t_{22}(u)-t_{21}(u+h)t_{12}(u)\bigr]\\
&+t_{21}(u)\bigl[t_{11}(u+h)t_{23}(u)-t_{21}(u+h)t_{13}(u)\bigr]\Bigr\},
\end{aligned}\label{eq:Bexact}
\end{equation}

$s_1=\nu'/\nu$, $s_2=1/((u-\theta_1)(u-\theta_2))$, checked against the chain for $b_0,\dots,b_6$ at three values of $u$ (orders 7, 8 vanish).
Expanding $B/\nu^3=1+hb_1+h^2b_2+\dots$,

\begin{equation}
b_1=2L_{11}+L_{22}+\frac{\nu'}\nu,\qquad b_2=2C_{11}+C_{22}+L_{11}'+L_{11}^2+2L_{11}L_{22}-L_{21}L_{12}+L_{21}L_{23}+\frac{\nu'}\nu(2L_{11}+L_{22})+s_2.\label{eq:b12}
\end{equation}

The first, $b_1$, is diagonal (Cartan) with eigenvalue $\sum_s(1+a+b+c)/(u-\theta_s)$ and has 25 eigenspaces of dimensions $1$ (4 times), $2$ (12 times) and $4$ (9 times): it does not resolve the 64 labels. $b_2$ needs the cross-site $C_{ij}$: it lies outside the span of $\{L_{ij}L_{kl},L',L'',L,1\}$ (171 residual entries) and inside it once the $C_{ij}$ are added.
Further facts, all exact:

\begin{itemize}
\item the $b_k(u)$ do not commute individually, $[b_i(u),b_j(v)]\neq0$ for all $(i,j)$ except $(1,1)$ and $(6,6)$, but $\sum_{i+j=n}[b_i(u),b_j(v)]=0$ for $n=2,\dots,8$;
\item compressed to the $b_1$-eigenspaces, $b_2$ has as eigenvalues the $h^2$-coefficient of the label-rule series (25/25 eigenspaces) and splits every eigenspace completely: the 64 pairs $(b_1,b_2^{\rm compressed})$ are distinct; $b_2$ alone has 62 simple roots and one double root with a Jordan block (the $b_k$ are not individually diagonalisable);
\item $B$ commutes with the grading $H_{\rm tot}=\sum_s(2E^s_{11}+E^s_{22})$, with level block sizes $1,4,8,12,14,12,8,4,1$, and every SoV covector lives in the level $D(x)=\sum_s(6-(1+a+b+c))$;
\item the Lie algebra generated by all coefficient matrices $B^{(m)}_k$ is \emph{solvable} on every level block (derived series: $80\to67\to45\to10\to0$ on the 14-block), so the $b_k(u)$ are simultaneously triangularisable block by block, with the level flag.
\item \textbf{No single polynomial $B_G(u)$ of degree $\le3$ in $L,L'$ triangularises the jet basis} (tried: $b_1,b_2,b_3,b_2-b_1^2/2$ in the jet basis; $b_1$ has only 25 distinct values and is not triangular inside level blocks).
The operator that does resolve all 64 labels is the $b_1$-block compression of $b_2$, which needs the projectors onto $H_1,H_2$ eigenspaces: not a polynomial of small degree in $L$.\footnote{Agent D3, \texttt{RESULTS-\allowbreak{}gaudin-\allowbreak{}B.\allowbreak{}md} B1, B5.}
\end{itemize}

\subsection{Where the degeneracy lives}\label{sec:degeneracy}

We expected degeneracy in $B$: at $h\to0$ all $64$ SoV covectors become proportional.
That is refuted.
In the XXX companion frame the 64 covectors $\langle x|(h)$, normalised by $\langle x|r=1$, are rational in $h$ of degree $\le5$ over $\le5$ (common denominator; degree tally $0:11,\,1:23,\,2:19,\,3:7,\,4:3,\,5:1$), validated at six unseen rational $h$ (64/64), and their \emph{leading terms at $h=0$ are pairwise distinct and linearly independent (rank 64)}.
The degeneracy lives in the diagonal frame, where the Gaudin states live: there the covectors are $f^d_x(h)=f_x(h)(P\otimes P)$ with $P=P_{\rm phys}^{-1}(e^{h\zeta})$, every entry $O(h)$ with singular values $h(1,h,h,h^2,h^2,h^3,h^3,h^4)$. All 64 have valuation $h^2$ and \emph{all 64 leading covectors are proportional to $x^G\otimes x^G$} (rank 1).
The ranks of the span of Taylor coefficients of $f^d_x/h^2$ to order $K=0,\dots,8$ are $1,5,13,25,39,51,59,63,64$. Saturating over $\mathbb{Q}[[h]]$ by echelon reduction (dividing by $h$ when the leading term vanishes) yields the \emph{jet basis} with defects per level $1,4,8,12,14,12,8,4,1$, the coefficients of $(1+2t+2t^2+2t^3+t^4)^2$: per site eight jets of weights $0,1,1,2,2,3,3,4$, the one-site depth of $P_{\rm phys}$.\footnote{Agent D3, \texttt{RESULTS-\allowbreak{}gaudin-\allowbreak{}B.\allowbreak{}md} B2, at $\zeta=P,P',X_1$.}

\subsection{The canonical jet basis and its wavefunction}\label{sec:canjet}

The jet basis is made canonical by the wavefunction rule.
Its one-site monomials are, in terms of the dual Q-functions at $\theta_s$,

\begin{equation}
\begin{aligned}
m_0&=P^3p_{02},& m_1&=P^{3\prime}p_{02},& m_2&=P^3p_{03},& m_3&=P^{3\prime}p_{03},\\
m_4&=P^{3\prime}p_{12},& m_5&=P^3p_{23},& m_6&=P^{3\prime}p_{13},& m_7&=P^{3\prime}p_{23},
\end{aligned}\label{eq:m8}
\end{equation}

with weights $w=c+k+l=2,3,3,4,4,5,5,6$: eight monomials, one per site label ($=\dim V$).
These are what the twelve $P^{3(c)}p_{kl}$, $c\le1$, $l\le3$ reduce to ($p_{01}=0$, $p_{12}=\kappa p_{02}$, $p_{13}=\kappa p_{03}$).
The 64 products $m_i(\theta_1)m_j(\theta_2)$ are independent functions on the 64 states ($\sigma_{\min}/\sigma_{\max}=6\cdot10^{-12}$).
The canonical jet covectors $\langle i,j|^c$ are defined by

\begin{equation}
\frac{\langle i,j|^c\Psi_c\rangle}{\langle x_0|\Psi_c\rangle}=\frac{m_i(\theta_1)\,m_j(\theta_2)}{M_1M_2},\label{eq:wfrule}
\end{equation}

in the Gaudin companion frame ($\Psi_c=(P_{\rm phys}^{-1}\otimes P_{\rm phys}^{-1})|_{z=\zeta}\Psi_d$).
The $64\times64$ matrix $Y$ expressing them in the physical basis is \emph{exactly $\zeta$-independent and rational}: identical exact rationals from the states at $P$, $P'$, $X_1$ (833 non-zero entries, inverse 975).
Only five of the 64 covectors are product covectors (the $8\times8$ coefficient matrix has rank 1): $(0,0),(0,1),(0,2),(1,0),(2,0)$; the other 59 couple the two sites through $1/\theta_{12}$ (ranks $2$ to $8$).
Examples: $\langle00|=e_{00}$, $\langle01|=-e_{01}$, $\langle02|=-2e_{02}+e_{01}-2t\,e_{00}$ with $t=1/\theta_{12}$. The rule (\ref{eq:wfrule}) was verified to $6\cdot10^{-76}$ ($P$), $1.4\cdot10^{-74}$ ($P'$), $1.5\cdot10^{-73}$ ($X_1$) on all $64\times64$ with $Y$ built at $P$ only; the saturation basis is block-triangular with respect to it, the filtration being the jet-weight filtration (level $d=w_i+w_j-4$).\footnote{Agent D3, \texttt{RESULTS-\allowbreak{}gaudin-\allowbreak{}B.\allowbreak{}md} B3.
The unipotent gauge freedom inside each filtration level is irrelevant for $G$ (a bilinear form on jets, not on the basis).
A determined solve at $P$ is the weakest evidence; the independent evidence is rationality, exact triangularity, $\zeta$-independence and the match with $G$.}

\subsection{The K tensor and the derivation of \texorpdfstring{$G$}{G}}\label{sec:Ktensor}

Expanding the companion-frame family $\Phi_c^\infty=\Omega_c^\infty[a]\otimes\Omega_c^\infty[a']$ in the canonical jet basis gives its K tensor $\Theta_\mu=(\Omega_c^\infty[a]\otimes\Omega_c^\infty[a'])Y^{-1}$, with exactly \emph{25 non-zero entries out of 64}, all on $N_5\times N_5$ (in the saturation basis: 36 non-zero entries).
Writing $\Theta_{ij}=\Theta(a,a',1/\theta_{12})$ one has

\begin{equation}
\Theta_{ij}(a,a',t)=\frac{G_{ij}(1/a,1/a',t)}{4\,\omega_1^3\omega_1'^3}\quad\text{(adjoint), i.e. }\ \Theta=\frac{K}{(\omega_1\omega_1')^{2\lambda_1-\lambda_2}},\label{eq:Theta}
\end{equation}

as an exact polynomial identity in $(a,a')$, 25/25 entries, at three values of $\theta_{12}$ ($31/35,29/45,5/6$).
This is the third derivation of the kernel.
The kernel has rank $1$ at $1/\theta_{12}=0$ (the sites decouple, $G$ becomes a product) and rank $5$ otherwise, so the whole site coupling is through $1/\theta_{12}$. For equal twist (finite $\omega_2$) all 64 jet coefficients are non-zero (levels $1,4,8,12,14,12,8,4,1$) and the jet expansion reproduces $\langle\Phi|\psi\rangle/\langle x_0|\psi\rangle$ to $6\cdot10^{-75}$; moreover $Q_{\rm num}^{\rm eq}M_1M_2/\det N_B(\omega)$ is state-independent (spread $1.2\cdot10^{-75}$, value $12.5381319888976$ at $\zeta=P$, $\theta=(2/7,-3/5)$), i.e. the $6\times6$ determinant of Section~\ref{sec:detN} is a Q-number of the same kind.\footnote{Agent D3, \texttt{RESULTS-\allowbreak{}gaudin-\allowbreak{}B.\allowbreak{}md} B4.
The reproduction of the plan's expectation that the kernel is the ``K tensor times wavefunction monomials'' was the main goal of round 3; the origin of $c$ is the frame tie $\Omega_c^\infty[1/\omega]P_{\rm phys}^{-1}(z=\zeta)=\Omega_D^\infty[\zeta,\omega]/\kappa^G$ exactly.}

\section{The overlap as a function of the Laurent data}\label{sec:nu}

\textbf{Punchline.} Every Plücker monomial is $M$ times a rational function of the Laurent data $(H_s,a_s,b_s)$ of the state at the two sites, so $c\,Q_{\rm num}=\sum_{ij}G^s_{ij}\nu_i(\theta_1)\nu_j(\theta_2)$ and neither the Q polynomials nor their normalisation enter: the Q-number is a function of eigenvalue data.

\subsection{The Frobenius reduction (adjoint)}\label{sec:frobred}

Near $\theta_s$, with $x=u-\theta_s$ and exponents $\{0,2,4\}$ for the adjoint, every solution is $P^j=A_jf_0+B_jf_2+C_jf_4$ with Frobenius series

\begin{equation}
f_0=1+\kappa x+\gamma_2x^2+\gamma_3x^3+\dots,\qquad f_2=x^2+d_3x^3+\dots.\label{eq:fseries}
\end{equation}

To derivative order $3$ only $f_0,f_2$ matter, and the Plücker coordinates are $p_{kl}=[A,B]_{23}\,w_{kl}$ with $[A,B]_{23}=A_2B_3-A_3B_2$ and $w_{02}=2$, $w_{03}=6d_3$, $w_{12}=2\kappa$, $w_{23}=12(\gamma_2d_3-\gamma_3)$, $w_{01}=0$. Hence $N_i=M\nu_i$ with

\begin{equation}
\nu=(1,\ \kappa,\ \pi,\ \kappa\pi,\ \varepsilon_s),\qquad \pi=\frac{p_{03}}{p_{02}}=3d_3,\qquad \varepsilon_s=\frac{p_{23}}{p_{02}}=6(\gamma_2d_3-\gamma_3),\label{eq:nu}
\end{equation}

and $M$ cancels from the left-hand side of (\ref{eq:jet}):

\begin{equation}
c\,Q_{\rm num}=\sum_{ij}G^s_{ij}(\omega_1,\omega_1',\theta_{12})\,\nu_i(\theta_1)\,\nu_j(\theta_2),\qquad G^s=G/4.\label{eq:nuform}
\end{equation}

The Frobenius recursion, with $\delta=\theta_s-\theta_t$ ($t\ne s$), $\sigma_1=\sum\zeta$, $e_2=e_2(\zeta)$, $e_3=e_3(\zeta)$, gives exactly (sympy)

\begin{equation}
\begin{aligned}
\kappa&=\frac{a_s}3,\qquad \pi=2H_s-a_s-2\sigma_1-\frac6\delta,\\
\varepsilon_s&=-\frac23\Bigl[\frac{H_ta_s}\delta-a_sb_s+a_se_2+\frac{3a_s}{\delta^2}-\frac{3a_t}{\delta^2}-\frac{3b_t}\delta-3e_3\Bigr],\qquad b_s=\frac{a_s(H_s-a_s)}3,
\end{aligned}\label{eq:nulaurent}
\end{equation}

the last being the first no-log condition (\ref{eq:nolog1}) with $c_2=3$. We verified (\ref{eq:nulaurent}) against $N_i/M$ built from the Q polynomials to $10^{-174}$--$10^{-193}$ on all six regime branches, both sites, at $\epsilon=10^{-6}/8$, and the no-log relation to $10^{-199}$. Thus $Q_{\rm num}$ is a function of $(H_1,a_1,H_2,a_2)$, $\zeta$, $\theta_{12}$, $\omega_1,\omega_1'$ only.\footnote{Agent U2, \texttt{RESULTS-\allowbreak{}untwist-\allowbreak{}limit.\allowbreak{}md} L3 and \texttt{untwist\_\allowbreak{}limit\_\allowbreak{}frob.\allowbreak{}py}; independently re-verified by U1 (\texttt{RESULTS-\allowbreak{}untwisted-\allowbreak{}Q.\allowbreak{}md} T2: the $\nu$-formulas hold to $10^{-185}$ for all six branches, both sites).
The result needs no use of the Q polynomials; it was reached by ``the Frobenius mechanism'', which is representation-generic.} The reduction also improves the limit identity by 14 or more digits (the residuals of Section~\ref{sec:untwist} drop from $10^{-43}$--$10^{-75}$ to $10^{-55}$--$10^{-157}$) and requires $\epsilon$-orders $\le5$ instead of $8$.

\subsection{Why this must hold for every \texorpdfstring{$\Lambda$}{}}\label{sec:nugen}

The argument is local.
For exponents $\{0,\lambda_2+1,\lambda_1+2\}$ the jet of every solution up to order $\lambda_1+1$ lies in the 2-plane spanned by the exponent-0 and the exponent-$(\lambda_2+1)$ Frobenius vectors; the coefficients of these vectors are rational functions of $(H_s,a_s,b_s)$, the other site's data through the regular parts, and $e_k(\zeta)$. Hence every non-forced $p_{kl}=D_s\,\nu_{kl}$ with $D_s$ the $2\times2$ determinant of the leading components of the two non-trivial solutions (the analogue of $[A,B]_{23}$) and $\nu_{kl}$ rational, and $M=P^3D_s\nu_{0,\lambda_2+1}$, so that $P^3D_s$ cancels between the sides.
We have \emph{not} derived the $\nu_{kl}$ for $\lambda_1\ge3$, and the untwisting limit for $[3,1,0]$, $[3,2,0]$ is not done; the general statement is an expectation based on the adjoint and on the verified exponent and jet-rank structure (Section~\ref{sec:frob}).
This is work package B of round 5.

\section{The general kernel}\label{sec:kernel}

\textbf{Punchline.} In the companion normalisation $\Theta=K/(\omega_1\omega_1')^{2\lambda_1-\lambda_2}$ the kernel is a polynomial in $(a,a',t)=(1/\omega_1,1/\omega_1',1/\theta_{12})$. At $t=0$ it factorises into one-site kernels $\theta_i(a)\theta_j(a')$, and the one-site kernel is a dressed Wronskian with a closed form valid for all $\Lambda$. Every two-row kernel is a block of the one-row master kernel $R^{\lambda_1}$ times $\theta\otimes\theta'$ (block theorem). $R^{\lambda_1}$ is known to first order in $t$ for all pairs and to all orders for $[l,l,0]$; the general four-index form is open.

\subsection{Conventions for this section}\label{sec:kconv}

Put $n=\lambda_1+1$ and let pairs $(kl)$ run over $0\le k<l\le n$ with $l\ge\lambda_2+1$. Let $a=1/\omega_1$, $a'=1/\omega_1'$, $t=1/\theta_{12}$, $x=at$, $y=a't$, and let the monomial weight be $e_{kl}=k+l-\lambda_2-1$ (relative to $M$).
The jet formula in canonical normalisation is

\begin{equation}
c\,Q_{\rm num}\,M_1M_2=\sum K_{ij}\,m_i(\theta_1)m_j(\theta_2),\qquad \Theta:=\frac K{(\omega_1\omega_1')^{2\lambda_1-\lambda_2}},\label{eq:Thetadef}
\end{equation}

with $M=P^3p_{0,\lambda_2+1}$ exactly, no factor $4$. (D2's $G=4K$ for the adjoint, from its vacuum monomial $2h^2P^3p_{02}$; D3's identity (\ref{eq:Theta}) reads $\Theta=G/(4\omega_1^3\omega_1'^3)$.) The common exponential factor $e^{(\zeta_2+2\zeta_3)\theta_s}$ of all monomials of a site is dropped, since the formula is homogeneous of degree one in the monomials, so every quantity is an exact rational.
We use the closed forms below only after checking them at points not used in their discovery.

\subsection{The one-site kernel}\label{sec:onesite}

At $t=0$ the sites decouple and $\Theta\to\theta(a)\otimes\theta(a')$ exactly (symbolic identity entry by entry, for all computed two-site representations).
The one-site kernel $\theta_{kl}(a)$ is exactly computable for any $\Lambda$: for one site the eigenstates are GT weight vectors, the Q's are exact rationals, and the same fit gives, with rank equal to the number of pairs and residual exactly zero (eight twist points, all diagonal states, $a$-dependence from four values of $\omega$),

\begin{equation}
\theta_{kl}(a)=c_{kl}\,a^{k+l-\lambda_2-1},\qquad c_{kl}=(-1)^{k+l-\lambda_2-1}\binom{\lambda_1-\lambda_2}{l-\lambda_2-1}\binom{\lambda_1+1}{k}\frac{l-k}{(\lambda_2+1)\binom{l}{\lambda_2+1}},\label{eq:ckl}
\end{equation}

independent of $\omega$ and $\zeta$, with $c_{0,\lambda_2+1}=1$. The formula was guessed from $\Lambda\le[4,2,0]$ and then correctly \emph{predicted} the other eight representations $[3,0,0],[3,3,0],[4,0,0],[4,3,0],[4,4,0],[5,1,0],[5,2,0],[5,3,0]$ (17 in all).
Examples (pairs in lexicographic order): $[2,1,0]$: $(02)\,1,(03)\,{-\tfrac12},(12)\,{-\tfrac32},(13)\,1,(23)\,{-\tfrac12}$; $[1,1,0]$: $(02)\,1,(12)\,{-1}$; $[3,1,0]$: $1,-1,\tfrac13,-2,\tfrac83,-1,-2,1,-\tfrac13$. Appendix~\ref{app:C} has the table up to $\lambda_1=4$.

The structure is a decomposable bivector.
For $\lambda_2=0$ the generating symbol is

\begin{equation}
\sum_{k<l}\theta^{[\lambda_1,0,0]}_{kl}\bigl(x_1^kx_2^l-x_1^lx_2^k\bigr)=(x_2-x_1)\bigl[(1-ax_1)(1-ax_2)\bigr]^{\lambda_1},\label{eq:symbol}
\end{equation}

exact for $\lambda_1=1,\dots,5$, and in jet language, for general $\lambda_2$,

\begin{equation}
\sum_{k<l}\theta_{kl}(a)\,m_{kl}=\epsilon_\Lambda\,a^{-\lambda_2}\,P^3\;W\bigl[(1-a\partial)^{\lambda_1}P^2,\ (1-a\partial)^{\lambda_1}P^3\bigr](\theta),\qquad \epsilon_\Lambda=(-1)^{\lambda_2}\frac{\lambda_2!(\lambda_1-\lambda_2)!}{\lambda_1!},\label{eq:dressedW}
\end{equation}

a Wronskian of the two dual Q's after dressing by $(1-a\partial)^{\lambda_1}$. In Taylor-normalised jets $\tilde p_{kl}=p_{kl}/(k!\,l!)$ the kernel is $\hat g\wedge\hat h$ with $\hat g_k=(-a)^k/(n-1-k)!$, $\hat h_k=(-a)^k/(n-k)!$: it is a Plücker vector of a 2-plane, the image of the lowest-weight vector $(\xi_1-\xi_2)$ of the spin-$\lambda_1$ irrep of an auxiliary $\mathfrak{sl}_2$ under $\exp(-aS_+)$. A one-site top-state sum rule fixes the normalisation (Section~\ref{sec:sumrule}).\footnote{Agent G1, \texttt{RESULTS-\allowbreak{}general-\allowbreak{}kernel.\allowbreak{}md} headline 1, K1 (exact Fractions; GT states with multiplicity one only: 6/8, 9/15, 12/24, \dots of the states; the others are not GT eigenvectors and were not used) and section K3(d).
The symbol identity holds for $\lambda_1=1,\dots,5$ and $c^{(\lambda_2)}_{kl}/c^{(0)}_{kl}=\epsilon_\Lambda$ for all pairs.}

\subsection{The block theorem}\label{sec:block}

The second structural fact: the $\lambda_2$-dependence of the two-site kernel is trivial.
For all non-forced pairs $i,j$,

\begin{equation}
\Theta^{[\lambda_1,\lambda_2,0]}_{ij}(a,a',t)=\Bigl[\frac{\lambda_2!(\lambda_1-\lambda_2)!}{\lambda_1!}\Bigr]^2(aa')^{-\lambda_2}\;\Theta^{[\lambda_1,0,0]}_{ij}(a,a',t),\label{eq:block}
\end{equation}

so that, equivalently, $\Theta_{ij}=\theta_i(a)\theta_j(a')R^{\lambda_1}_{ij}(at,a't)$ with a universal, $\lambda_2$-independent Laurent polynomial $R^{\lambda_1}_{(kl),(k'l')}(x,y)$, $R(0,0)=1$, homogeneous of degree $0$ in $(a,a',\theta_{12})$ (a function of $x,y$ only; the argument is $(x,y)=(at,a't)$).
All the $\lambda_2$ content is in the one-site kernel $\theta$. In the $K$ normalisation this has a particularly clean form (a rewriting of (\ref{eq:block}) that we also checked on the stored exact kernels, see below):

\begin{equation}
K^{[\lambda_1,\lambda_2,0]}_{ij}=\Bigl[\frac{\lambda_2!(\lambda_1-\lambda_2)!}{\lambda_1!}\Bigr]^2K^{[\lambda_1,0,0]}_{ij}\quad(\text{non-forced pairs}).\label{eq:blockK}
\end{equation}

For $\lambda_1=3$ the prefactor is $1/9$ for both $\lambda_2=1,2$, so the $[3,1,0]$ and $[3,2,0]$ kernels in the $K$ normalisation coincide on the seven shared pairs.
We confirmed this at two parameter points, $(\omega_1,\omega_1',\theta_{12})=(5/2,7/4,31/35)$ and $(-2/5,3,29/45)$, entry by entry on the exact stored polynomial laws (0 mismatches in $49$ entries each), and confirmed (\ref{eq:block}) for $[3,1,0]$ against $[3,0,0]$ at $(a,a',t)=(2/5,4/7,35/31)$ (0 mismatches in 81 entries).
G1's own verification: symbolic for $\lambda_1\le3$ (all $\lambda_2$), symbolic for $[4,4,0]$ against $[4,0,0]$, direct independent fits for $[4,1,0],[4,2,0],[4,3,0]$ at two $\theta_{12}$ and two $(\omega,\omega')$ (deviation $\le3\cdot10^{-58}$).
The corner block (pairs $(k,\lambda_1+1)$) equals $N(aa')^{\lambda_1-\lambda_2}$ times the single-index kernel $\Theta^{[\lambda_1,\lambda_1,0]}$, $N=[\lambda_2!(\lambda_1-\lambda_2)!/\lambda_1!]^2$.

The block theorem is useful as a predictor: $[4,2,0]$ has 12 pairs and $\Theta=\tfrac1{36}(aa')^{-2}\Theta^{[4,0,0]}|$ (verified directly, deviation $\le8\cdot10^{-62}$), $[5,1,0]$ has 20 pairs, $\Theta=\tfrac1{25}(aa')^{-1}\Theta^{[5,0,0]}|$, $a$-degree $9$ and $t$-degree $16$. The $[5,1,0]$ prediction was not tested directly: a corner-block check was started but the $1225$-state build was too slow and was stopped.

\subsection{Structure: polynomiality, weights, degrees}\label{sec:structure}

Every kernel computed is a polynomial in $(a,a',t)$ with rational coefficients; each monomial $a^pa'^qt^n$ of entry $(i,j)$ obeys $p+q-n=e_i+e_j$ with $0\le p,q\le2\lambda_1-\lambda_2$ (maximum attained) and the exchange rule

\begin{equation}
\Theta_{ij}(a,a',t)=\Theta_{ji}(a',a,-t)\quad\text{exactly}.\label{eq:exch}
\end{equation}

Equivalently, writing $\Theta_{ij}=\sum_nt^n\Theta_{ij,n}(a,a')$, each $\Theta_{ij,n}$ is homogeneous of degree $e_i+e_j+n$ in $(a,a')$ (the ``torus'' law: $\omega_1,\omega_1'$ scale like $\zeta$ and $\theta_{12}$ like $1/\zeta$).
The $t$-degree is $D_t=4(\lambda_1-1)-\max(0,2\lambda_2-3)$ for $\lambda_1\ge2$, $D_t=1$ for $\lambda_1=1$, and $2l-1$ for $[l,l,0]$; the one-row sequence $1,4,8,12$ is $4(l-1)$ from $l=2$ on.
The earlier guess $\binom{l+2}2-2=13$ for $[4,0,0]$ is refuted (12).
Predicted: $[5,0,0]$, $[5,1,0]$: $16$; $[5,2,0]$: $15$; $[5,3,0]$: $13$. Table~\ref{tab:kstruct} summarises.

{\footnotesize
\begin{longtable}{lccccl}
\caption{Structure of the two-site kernels (G1 \texttt{RESULTS-\allowbreak{}general-\allowbreak{}kernel.\allowbreak{}md} K0, K2, K3): number of pairs, $a$-degree ($=2\lambda_1-\lambda_2$), $t$-degree, generic rank of $\Theta$, and ranks of the $t$-expansion coefficients $\Theta_n$. All kernels satisfy $\Theta(t=0)=\theta\otimes\theta'$, the exchange rule, and the top-state sum rule exactly.
Rank-deficient at a generic point: $[1,0,0],[3,0,0],[3,1,0],[4,0,0],[4,1,0]$ (cause not understood).}\label{tab:kstruct}\\
\toprule
$\Lambda$ & pairs & $a$-deg & $t$-deg & generic rank & ranks of $\Theta_n$, $n=0,1,\dots$ \\
\midrule
\endfirsthead
\toprule
$\Lambda$ & pairs & $a$-deg & $t$-deg & generic rank & ranks of $\Theta_n$, $n=0,1,\dots$ \\
\midrule
\endhead
\bottomrule
\endlastfoot
$[1,0,0]$ & 3 & 2 & 1 & 2 & 1,2 \\
$[1,1,0]$ & 2 & 1 & 1 & 2 & 1,2 \\
$[2,0,0]$ & 6 & 4 & 4 & 6 & 1,2,4,6,6 \\
$[2,1,0]$ & 5 & 3 & 4 & 5 & 1,2,4,5,4 \\
$[2,2,0]$ & 3 & 2 & 3 & 3 & 1,2,3,2 \\
$[3,0,0]$ & 10 & 6 & 8 & 8 & 1,2,4,8,10,10,10,8,6 \\
$[3,1,0]$ & 9 & 5 & 8 & 8 & 1,2,4,8,9,9,8,6,4 \\
$[3,2,0]$ & 7 & 4 & 7 & 7 & 1,2,4,7,7,6,4,2 \\
$[3,3,0]$ & 4 & 3 & 5 & 4 & 1,2,3,4,3,2 \\
$[4,0,0]$ & 15 & 8 & 12 & 12 & 1,2,4,8,13,15,15,15,15,12,10,8,6 \\
$[4,4,0]$ & 5 & 4 & 7 & 5 & 1,2,3,4,5,4,3,2 \\
$[4,1,0]$ & 14 & 7 & 12 & 12 & 1,2,4,8,13,14,14,14,13,10,8,6,4 (block thm) \\
$[4,2,0]$ & 12 & 6 & 11 & 12 & 1,2,4,8,12,12,12,11,9,6,4,2 (block thm) \\
$[4,3,0]$ & 9 & 5 & 9 & 9 & 1,2,4,7,9,9,8,6,4,2 (block thm) \\
\end{longtable}}


The supports shrink with $n$: the nonzero entries of $\Theta_n$ per order are $25,23,19,12,4$ of $25$ for the adjoint and $49,47,43,37,29,19,9,2$ of $49$ for $[3,2,0]$; the top order has only two to four non-zero entries.
The ``two channels per order'' doubling of the rank ends at $n=3$ (for $[4,0,0]$: $13$ at $n=4$, not $16$).

\subsection{The first-order term}\label{sec:first}

The coefficient of $t^1$ is known in closed form for all pairs and all representations.
With $\rho(p)=(n-k)(n-l)$, $\sigma(p)=k+l-1$, $f_2(p)=2(n-1)-\sigma(p)$, $N_0=n(n-1)$, and $p=(kl)$, $q=(k'l')$,

\begin{equation}
\Theta_{1,pq}=t\,\theta_p(a)\theta_q(a')\Bigl[\alpha_{pq}a+\beta_{pq}a'+\gamma_{pq}\frac{a^2}{a'}+\varepsilon_{pq}\frac{a'^2}{a}\Bigr],\label{eq:first}
\end{equation}

\begin{equation}
\alpha_{pq}=-\frac3{n^2}\rho(p)\bigl(\rho(q)-N_0\bigr)-\frac2n\rho(p)\sigma(q)+\frac1nf_2(p)\bigl(\rho(q)-N_0\bigr)+f_2(p)\sigma(q),\qquad\beta_{pq}=-\alpha_{qp},
\end{equation}

\begin{equation}
\gamma_{pq}=\frac{k'l'\rho(p)}{n^2},\qquad\varepsilon_{pq}=-\frac{kl\,\rho(q)}{n^2}.
\end{equation}

This formula was fitted on $\lambda_1=1,2,3$ and then reproduced the $\lambda_1=4$ single-index kernel and all $\lambda_2$-blocks exactly: 0 mismatches over 1062 random-point entry tests ($\lambda_1\le4$) and 675 entries of the held-out $[4,4,0]$ and $[4,0,0]$. As a matrix on the pair space $\alpha,\beta$ have rank 2 and $\gamma,\varepsilon$ rank 1; the factors $(n-k)(n-l)$ and $k'l'$ are the raising and lowering amplitudes of the spin-$n/2$ representation of the auxiliary $\mathfrak{sl}_2$ acting on the jet index (double hops: $\gamma$ raises both indices of site 1 and lowers both of site 2).
On a single index ($l=l'=n$): $\alpha=(\lambda_1-k)k'$, $\beta=-(\lambda_1-k')k$, $\gamma=\varepsilon=0$. $\Theta_1$ has rank exactly 2, with $\theta(a)$ in its column space and $\theta(a')$ in its row space.
Beyond first order we know the extreme terms.
At order $2$ the exponent pairs $x^uy^v$ ($u+v=2$) are $(u,v)=(-2,4),(-1,3),(0,2),(1,1),(2,0),(3,-1),(4,-2)$ with coefficient ranks $1,2,4,5,4,2,1$; the extreme ones are $\tfrac12(n-k)_2(n-l)_2k'_2l'_2/(n(n-1))^2$ and its mirror.
General $m$ (exact on all 225 entries of $[4,0,0]$ for $m=1,\dots,4$): the coefficient of $x^{2m}y^{-m}$ in $R_{pq}$ is

\begin{equation}
\frac1{m!}\frac{(n-k)_m(n-l)_m(k')_m(l')_m}{(n)_m^2}=(m!)^3\frac{\binom{n-k}m\binom{n-l}m\binom{k'}m\binom{l'}m}{\binom nm^2},\label{eq:hops}
\end{equation}

with falling factorials $(z)_m$, and a sign $(-1)^m$ and exchanged roles for the mirror $y^{2m}x^{-m}$.

For $\lambda_1=1$ the full kernel is an R-matrix: $\Theta^{[1,\lambda_2,0]}=(1+z)\,\theta(a)\otimes\theta(a')-z\,\theta(a')\otimes\theta(a)$ with $z=t/(\omega-\omega')=aa't/(a'-a)$, i.e. a permutation exchange of the two coherent states with the rational R-matrix $R(u)=u+tP$.

\subsection{The single-index family \texorpdfstring{$[l,l,0]$}{[l,l,0]} to all orders}\label{sec:single}

For $[l,l,0]$ (one Plücker index $k\in\{0,\dots,l\}$) the kernel is

\begin{equation}
\Theta_{kk'}=(-1)^{k+k'}\binom lk\binom l{k'}a^ka'^{k'}\sum_{m=0}^{l-k}\sum_{m'=0}^{l-k'}\binom{l-k}m\binom{l-k'}{m'}F_{mm'}(k,k')\,(at)^m(-a't)^{m'},\label{eq:single}
\end{equation}

\begin{equation}
F_{00}=1,\quad F_{m0}=(k')_m,\quad F_{0m'}=(k)_{m'},\quad F_{mm'}=(k+m+1)_{m'-1}(k'+m'+1)_{m-1}\,(kk'+m'k+mk'),
\end{equation}

here with rising Pochhammer symbols; uniformly $F=\Gamma(k+m+m')\Gamma(k'+m+m')\,[(k+m)(k'+m')-mm']/(\Gamma(k+m+1)\Gamma(k'+m'+1))$. It is exact for $l=1,\dots,4$ (all $4+9+16+25$ entries) and \emph{predicted correctly for $l=5$}: direct fits of $[5,5,0]$ (441 states, 36 unknowns, rank 36, residual $0.0$) agree to $1.5\cdot10^{-59}$--$9.2\cdot10^{-58}$ at three parameter points.
Its symbol is nice.
With $A=1-a\xi$, $B=1-a'\eta$ and $\hat\Theta(\xi,\eta)=\sum\Theta_{kk'}\xi^k\eta^{k'}$,

\begin{equation}
\hat\Theta^{[l,l,0]}=A^lB^l+t\,l^2aa'(\xi-\eta)A^{l-1}B^{l-1}+\dots,\label{eq:symbol2}
\end{equation}

every $t^n$ term with $n\ge1$ is divisible by $(\xi-\eta)$ (so on $\xi=\eta$, plane waves like the top state's $P^3=e^{\zeta_3u}$, there is \emph{no coupling}), and the top coefficient is $t^{2l-1}$ times $(-1)^{l+1}\frac{(2l-2)!(2l)!}{2(l-1)!\,l!}(aa')^l(\xi-\eta)$, with values $1,12,720,100800,25401600,10059033600$ for $l=1,\dots,6$ ($l\le4$ from fitted kernels, $l=5,6$ from (\ref{eq:single})).
In the variables $\Xi=\omega-\xi$, $H=\omega'-\eta$, $\nu=\omega-\omega'$, $\delta=\xi-\eta=\nu-\Xi+H$, the rescaled $\hat K=(\omega\omega')^l\hat\Theta$ is $\omega$-free: $K_1=\Xi H+t\delta$, $K_2=\Xi^2H^2+4t\delta\Xi H+2t^2\delta(3\delta-2\nu)-12t^3\delta$, $K_3=\Xi^3H^3+9t\delta\Xi^2H^2+18t^2\delta\Xi H(2\delta-\nu)+\dots$. In a Laplace--Borel form, $\hat\Theta=\sum_{K,K'}\binom lK\binom l{K'}\{\mathbb E[(u+vs)^K(u'+v's')^{K'}]-KK'vv'\,\mathbb E[(u+vs_2)^{K-1}(u'+v's_2')^{K'-1}]\}$ with $u=-a\xi$, $v=at$, $u'=-a'\eta$, $v'=-a't$ and independent Gamma variables $s\sim\Gamma(K')$, $s'\sim\Gamma(K)$, $s_2\sim\Gamma(K'+1)$, $s_2'\sim\Gamma(K+1)$ (the Pochhammer symbols are the moments).

Its determinants have a pattern.
With $\nu=\omega-\omega'=(a'-a)/(aa')$, $T=l(l+1)/2$ and $d_l=\lceil l/2\rceil$, $\det\Theta^{[l,l,0]}=c_l(aa')^{s_l}t^T\nu^{T-d_l}p_l(\nu,t)$ with $p_l$ homogeneous of degree $d_l$:

\begin{equation}
p_1=\nu+t,\quad p_2=\nu+3t,\quad p_3=\nu^2+6t\nu+3t^2,\quad p_4=\nu^2+10t\nu+15t^2,\quad p_5=\nu^3+15t\nu^2+45t^2\nu+15t^3.\label{eq:pl}
\end{equation}

The kernel loses rank at $\nu=0$ ($\omega=\omega'$, where CG-completeness fails) and on the roots of $p_l$. We note, as a curiosity, that $\Theta$ is rank-deficient at a generic point for $[1,0,0],[3,0,0],[3,1,0],[4,0,0]$ (ranks $2/3$, $8/10$, $8/9$, $12/15$): lost directions of unknown origin; all other computed representations have full rank.

\subsection{The top-state sum rule}\label{sec:sumrule}

A strong exact constraint comes from the top state $e_1\otimes e_1$, whose Q-functions are explicit: $P^3=e^{\zeta_3u}$, $P^2=e^{\zeta_2u}p_2$ with $p_2=\sum_{k\ge0}(\zeta_3-\zeta_2)^{-k}(\Pi^{\lambda_2})^{(k)}$ (from the Wronskian relation with $Q_1=e^{-\zeta_1u}$; monic, $p_2'-(\zeta_3-\zeta_2)p_2=-(\zeta_3-\zeta_2)\Pi^{\lambda_2}$).
The sum rule is

\begin{equation}
\sum_{ij}\Theta_{ij}\,m_i^{\rm top}(\theta_1)\,m_j^{\rm top}(\theta_2)=M^{\rm top}(\theta_1)M^{\rm top}(\theta_2)\,R(a)R(a'),\qquad R(a)=(1-\zeta_2a)^{\lambda_1-\lambda_2}(1-\zeta_3a)^{\lambda_1},\label{eq:sumrule}
\end{equation}

exact for ten representations (two-site; fully symbolic in $\zeta_2,\zeta_3,\theta_1,\theta_2,a,a'$ for $[1,0,0],[1,1,0],[2,2,0],[2,1,0]$, exact rational at 3--4 random points for the rest) and ten one-site ones.
The right-hand side is not $M^{\rm top}M^{\rm top}$, as the plan had it: the companion-frame top state is not $e_1\otimes e_1$ because $P_{\rm phys}^{-1}e_{\rm hw}$ is not parallel to $e_{\rm hw}$, and $c\,a^{2\lambda_1-\lambda_2}/x_{\rm hw}=R(a)$ with $x_{\rm hw}=(\zeta_2-\zeta_1)^{\lambda_2}$ exactly.
At $a=0$ the rule gives $R(0)=1$ and each further power of $a$ one linear relation between the top-state jets and $\Theta$: every candidate kernel must satisfy it.

\subsection{One-row representations are the master kernels}\label{sec:master}

The $\lambda_2=0$ representations work, with the same monomials, and are the master objects of the block theorem.
The oper has exponents $\{0,1,\lambda_1+2\}$, $c_2=0$ (a scalar, see Section~\ref{sec:laurent}), $M=P^3p_{01}$ with nothing forced, and $\binom{\lambda_1+2}2$ monomials ($3,6,10$ for $[1,0,0],[2,0,0],[3,0,0]$).
Nothing breaks, and in particular the plan's worry that $S_2$'s double pole would be non-scalar did not materialise.

\subsection{The auxiliary \texorpdfstring{$\mathfrak{sl}_2$}{sl} and what is open}\label{sec:auxsl2}

The jet index space (jets of order $\le\lambda_1+1$, dimension $\lambda_1+2$) carries an $\mathfrak{sl}_2$ of spin $(\lambda_1+1)/2$ acting on symbols by the M\"obius map $(U_sf)(\xi_1,\xi_2)=\prod_i(1-s\xi_i)^{\lambda_1+1}f(\xi_i/(1-s\xi_i))$; the one-site kernel obeys $\theta^{(a+s)}=U_s\theta^{(a)}$ exactly (checked symbolically, $\lambda_1\le4$), a coherent state of the spin-$\lambda_1$ irrep of $\Lambda^2V$. The two-site kernel is \emph{not} covariant under the diagonal shift $a\to a+s$, $a'\to a'+s$ (tested, $\lambda_1=1,2$).
Further tests that failed: $\Theta$ is not a compound matrix $\Lambda^2M$ ($\Theta_0$, $\Theta_1$ are Plücker, $\Theta_2,\Theta_3,\dots$ and $\Theta$ are not); the ``box exchange'' hypothesis (the kernel as a sum over permuted parameters $u=(\omega+ct,\omega'+c't)$ with shifts up to $\pm3$ in $t/2$ units) is refuted for $\lambda_1=2$ with the fused-R picture not literally a sum over parameter permutations; and the XXX R-matrix poles $(\omega-\omega'+mt)$ appear only in determinants of $\Theta$, not in its entries.
The number of $t$-orders is only linear in $\lambda_1$, not quadratic as a fusion count of $\lambda_1^2$ R-matrix factors would suggest.

\textbf{Open:} the four-index master kernel $R^{\lambda_1}_{(kl),(k'l')}(x,y)$ beyond first order.
We know the single-index family (all orders), the first-order term (all pairs), the extreme multi-hops, universality in $\lambda_2$, the rank pattern, and the sum rule; the full-$\lambda_1=5$ grids and the full $(a',t)$ polynomial grids for $[4,1,0],[4,2,0],[4,3,0]$ (replaced by the block theorem) were not done.

\section{The untwisting limit and the CG generating function}\label{sec:untwist}

\textbf{Punchline.} In the regime $\zeta=\epsilon\phi$, $\omega=\epsilon\hat\omega$ the jet formula is a statement about the $\epsilon$-jet of the Q-functions, with cancellations of depth up to six orders.
Reduced to Laurent data it becomes a polynomial identity in $(\hat\omega_1,\hat\omega_1')$ whose solution is the CG generating function.
For all six HW branches of the adjoint the identity holds to $10^{-55}$--$10^{-157}$ and returns every CG coefficient, $\phi$-independent to $2\cdot10^{-77}$.

\subsection{The regime and the difficulty}\label{sec:regime}

Let $\zeta=\epsilon\phi$, $\omega_1=\epsilon\hat\omega_1$, $\omega_1'=\epsilon\hat\omega_1'$ and $\epsilon\to0$. Chain side (Section~\ref{sec:chain}): $\langle\Phi_\infty|\Psi_B\rangle=\langle\varphi_u\otimes e_1^T|\Psi_0\rangle(1+O(\epsilon))$ for every HW-limit branch $\Psi_B=\Psi_0+\epsilon\Psi_1+\cdots$ with $\Psi_0$ an exact HW vector of weight $\mu$, G-norm 1; so

\begin{equation}
\langle\varphi_u(\hat\omega_1,\hat\omega_1')\otimes e_1|\Psi_0\rangle=\sum_\nu\bigl[e_1^T\rho(n')\bigr]_\nu\,C_\nu,\qquad C_\nu=\langle e_\nu\otimes e_1|\Psi_0\rangle,\label{eq:CGgf}
\end{equation}

is the CG generating function: a polynomial in the rotation parameters whose coefficients are the CG coefficients of the GT basis.
We want to write the left-hand side as a Q-expression using (\ref{eq:jet}).

The difficulty is that the limit is delicate on the Q side.
At $\zeta=0$ the three dual Q-functions $e^{\zeta_ju}p_j$ lose their exponentials and their Wronskians vanish (they carry factors $\zeta_b-\zeta_a$), so the untwisted limits of the three branch Q's are linearly dependent: the limit of the solution space is a $\zeta$-jet ($p^{(0)}$ and its $\epsilon$-derivatives), not just a triple of polynomials.
Moreover the jet formula's total is $\epsilon^{N_B}$ smaller than its individual terms, with $N_B=8,7,7,6,6,4$ for $(4,2,0),(4,1,1),(3,3,0),(3,2,1)\times2,(2,2,2)$, while single terms are $O(\epsilon^2\ldots\epsilon^3)$: a cancellation of up to six orders.

\subsection{\texorpdfstring{$\epsilon$}{}-order bookkeeping}\label{sec:bookkeeping}

The ingredients scale as follows (for the adjoint, at generic $\phi$; $\phi'$ gives identical exponents).
The frame scalar is $c=\epsilon^8\hat c$ with $\hat c=(\phi_2-\hat\omega_1)(\phi_3-\hat\omega_1)^2(\phi_2-\hat\omega_1')(\phi_3-\hat\omega_1')^2(\phi_1-\phi_2)^2$; the kernel splits by degree in $(\omega_1,\omega_1')$, $G^s(\epsilon\hat\omega)=\sum_d\epsilon^dG^{s[d]}(\hat\omega)$ (at fixed $\theta_{12}$); $M_1M_2\sim\epsilon^2$ on all branches; $\langle\Phi_\infty|\Psi_B\rangle\sim\epsilon^0$ at all 16 grid points; and

\begin{equation}
\langle x_0^G|\Psi_B\rangle\sim\epsilon^{2+k_B},\qquad k_B=0,1,1,2,2,4\quad\text{for}\quad(4,2,0),\,(4,1,1),\,(3,3,0),\,(3,2,1)\times2,\,(2,2,2).\label{eq:kB}
\end{equation}

Table~\ref{tab:order} gives the exponents.
The interpretation of $k_B$: on the exact HW kets of every non-top $\mu$, $\langle x(\phi)\otimes x(\phi)|v\rangle=0$ identically ($\le4\cdot10^{-202}$ at 200 digits, both $\phi$, the whole HW space and both octet copies); for the top it is $(\phi_2-\phi_1)^2$. So $\langle x_0^G|\Psi_B\rangle$ gets its first contribution from the $k_B$-th order vector $\Psi_{k_B}$, and we define

\begin{equation}
X_B:=\langle x(\phi)\otimes x(\phi)|\Psi_{k_B}\rangle,\label{eq:XB}
\end{equation}

which turns out to be a product of differences $\phi_i-\phi_j$ (below).
The consistency check is $N_B=8+0-(2+k_B)+2=8-k_B$ on all six branches.

{\small
\begin{longtable}{lcccccc}
\caption{$\epsilon$-orders for the six adjoint branches (U2 L1; the same at $\phi'$). $N_i$ are the exponents of the five jet monomials $(P^3p_{02},P^{3\prime}p_{02},P^3p_{03},P^{3\prime}p_{03},P^3p_{23})$ (the same on both sites); $N_B$ is the exponent of $\sum GNN$; ``first $n$'' is the first order at which $\langle x(\phi)\otimes x(\phi)|\Psi_n\rangle\ne0$.}\label{tab:order}\\
\toprule
branch & $k_B$ & $\langle x_0^G|\Psi_B\rangle$ & $M_1M_2$ & $N_i$ exponents & $N_B$ & first $n$ \\
\midrule
\endfirsthead
\toprule
branch & $k_B$ & $\langle x_0^G|\Psi_B\rangle$ & $M_1M_2$ & $N_i$ exponents & $N_B$ & first $n$ \\
\midrule
\endhead
\bottomrule
\endlastfoot
$(4,2,0)$ & 0 & $\epsilon^2$ & $\epsilon^2$ & 1,2,1,2,3 & 8 & 0 \\
$(4,1,1)$ & 1 & $\epsilon^3$ & $\epsilon^2$ & 1,1,2,2,3 & 7 & 1 (odd $n$ only) \\
$(3,3,0)$ & 1 & $\epsilon^3$ & $\epsilon^2$ & 1,2,1,2,3 & 7 & 1 (odd $n$ only) \\
$(3,2,1)$ a, b & 2 & $\epsilon^4$ & $\epsilon^2$ & 1,1,1,1,2 & 6 & 2 \\
$(2,2,2)$ & 4 & $\epsilon^6$ & $\epsilon^2$ & 1,1,1,1,1 & 4 & 4 ($n=0..3$ vanish to $10^{-74}$) \\
\end{longtable}}


\subsection{The leading polynomial identity}\label{sec:leadpoly}

All orders $t<N_B$ of $\sum_{ij}G^s_{ij}(\epsilon\hat\omega)N_iN_j$ cancel, with residuals $4.5\cdot10^{-51}$ ($(4,2,0)$), $1.2\cdot10^{-55}$, $1.1\cdot10^{-55}$, $6.6\cdot10^{-66}$, $2.9\cdot10^{-84}$ at $t=N_B-1$, against $0.2$--$0.4$ at $t=N_B$. The leading polynomial is $L_B(\hat\omega)=\sum_{d+n+m=N_B}G^{s[d]}_{ij}N^{(n)}_iN^{(m)}_j$, where $N_i^{(n)}$ is the $\epsilon^n$ coefficient of $N_i$, and the identity is

\begin{equation}
L_B(\hat\omega_1,\hat\omega_1')=\hat c\;\langle\varphi_u\otimes e_1|\Psi_0\rangle\;\frac{m_B}{X_B},\qquad m_B=\lim\epsilon^{-2}M_1M_2,\ \ \kappa_B=1,\label{eq:leadpoly}
\end{equation}

as a polynomial in $(\hat\omega_1,\hat\omega_1')$, coefficient by coefficient (relative coefficient residual $10^{-39}$ to $10^{-75}$ over six branches and two directions; the constant $\kappa_B$ between the two sides is exactly 1).
Each octet copy pairs with its own HW ket (different generating functions, same constant $R_B=m_B/X_B$).
For the top state $m_B=-(\phi_3-\phi_2)^2\theta_{12}^2$ ($Q_1=1$; checked to 13 digits), so $R_B(\mathrm{top})=-(\phi_3-\phi_2)^2\theta_{12}^2/(\phi_2-\phi_1)^2$; in general $M_s=P^3A_sQ_1$ with $A_s=(\zeta_3-\zeta_2)e^{\sigma\theta_s}(\theta_s-\theta_t)$ and $m_B=-(\phi_3-\phi_2)^2\theta_{12}^2Q_1^0(\theta_1)Q_1^0(\theta_2)$ (not separately checked beyond the top).
Which $(d,n,m)$ contribute at order $N_B$ (terms above $10^{-30}$): $(4,2,0)$: $d=2..6$, $n,m=1..4$ (largest $(4,2,2),(3,2,3),(5,1,2),(4,1,3)$); $(4,1,1)$ and $(3,3,0)$: $d=1..5$, $n,m\le5$ (largest $(3,2,2),(4,1,2),(2,2,3)$); octets $d=0..4$, $n,m\le4$; singlet only six triples $(1,1,2),(2,1,1),(1,2,1),(0,2,2),(0,1,3),(0,3,1)$.

{\small
\begin{longtable}{lrrrr}
\caption{The constant $R_B=m_B/X_B$ of (\ref{eq:leadpoly}) at the two directions (U2 table L2(c)): best-fit value, and relative residual of the polynomial identity.}\label{tab:RB}\\
\toprule
branch & $R_B(\phi)$ & residual $(\phi)$ & $R_B(\phi')$ & residual $(\phi')$ \\
\midrule
\endfirsthead
\toprule
branch & $R_B(\phi)$ & residual $(\phi)$ & $R_B(\phi')$ & residual $(\phi')$ \\
\midrule
\endhead
\bottomrule
\endlastfoot
$(4,2,0)$ & $-1.1217793595$ & $2.4\cdot10^{-43}$ & $-0.0892224658$ & $7.7\cdot10^{-39}$ \\
$(4,1,1)$ & $-4.8470775935$ & $2.7\cdot10^{-47}$ & $0.1598811435$ & $2.8\cdot10^{-46}$ \\
$(3,3,0)$ & $-1.1349116102$ & $1.8\cdot10^{-47}$ & $0.0721026407$ & $1.6\cdot10^{-45}$ \\
$(3,2,1)$ a, b & $-7.5969660751$ & $2.9\cdot10^{-57}$, $7.3\cdot10^{-58}$ & $-0.2001611630$ & $2.1\cdot10^{-54}$, $3.4\cdot10^{-55}$ \\
$(2,2,2)$ & $-90.949002471$ & $7.0\cdot10^{-75}$ & $-0.7937987684$ & $3.8\cdot10^{-72}$ \\
\end{longtable}}


\subsection{The reduced form: only Laurent data}\label{sec:redform}

Using Section~\ref{sec:nu} the leading identity can be written without any Q-function or normalisation:

\begin{equation}
\sum_{ij}G^s_{ij}(\epsilon\hat\omega_1,\epsilon\hat\omega_1')\,\nu_i(\theta_1)\nu_j(\theta_2)=\epsilon^{6-k_B}\,\hat c\,\frac{\langle\varphi_u\otimes e_1|\Psi_0\rangle}{X_B}+O(\epsilon^{7-k_B}),\label{eq:limitnu}
\end{equation}

verified coefficientwise to $10^{-55}\ldots10^{-157}$ with lower orders vanishing to $10^{-62}\ldots10^{-165}$; no $M$, no $\kappa_B$ and no Q-normalisation enters.
The $\nu$-series are needed only to order $3$ (top), $5$ ($(4,1,1)$, $(3,3,0)$), $4$ (octets) and $2$ (singlet), i.e. the Laurent data to order $\epsilon^3,\epsilon^5,\epsilon^4,\epsilon^2$. Table~\ref{tab:nudepth} lists the $(i,n)$ pairs that enter.

{\small
\begin{longtable}{lcl}
\caption{Depth needed in the $\nu$-form (U2 L3): $i=0,\dots,4$ labels $(1,\kappa,\pi,\kappa\pi,\varepsilon_s)$, $n$ the $\epsilon$-order of $\nu_i$.}\label{tab:nudepth}\\
\toprule
branch & $\nu$ exponents & $n$ needed per site \\
\midrule
\endfirsthead
\toprule
branch & $\nu$ exponents & $n$ needed per site \\
\midrule
\endhead
\bottomrule
\endlastfoot
$(4,2,0)$ & 0,1,0,1,2 & only the leading coefficients: $(0,0),(1,1),(2,1),(3,2),(4,3)$ \\
$(4,1,1)$ & 0,0,1,1,2 & $i{=}1$: 0,1,2,4; $i{=}2$: 1,2,4; $i{=}3$: 1,2,3,5; $i{=}4$: 2,3,4 \\
$(3,3,0)$ & 0,1,0,1,2 & $i{=}1$: 1; $i{=}2$: 0,1,2,4; $i{=}3$: 1,2,3,5; $i{=}4$: 2,3,4 \\
$(3,2,1)$ a, b & 0,0,0,0,1 & $i{=}1..3$: 0..4; $i{=}4$: 1..4 \\
$(2,2,2)$ & 0,0,0,0,0 & $i{=}1..4$: 0..2 \\
\end{longtable}}


\subsection{The CG dictionary and \texorpdfstring{$\phi$}{}-independence}\label{sec:dictionary}

Solving the polynomial identity (\ref{eq:leadpoly}) or (\ref{eq:limitnu}) for the HW ket (least squares over the polynomial coefficients, unknown in the HW space $H_\mu$, then G-normalise and fix the sign by the first non-zero GT component) returns the exact CG slices: $C_7=1$ for the top ($\mathtt{gl3}$ index; notebook $\nu=1$), $-0.70710678119$ for $(4,1,1)$ and $(3,3,0)$, the octet pairs $(C_2,C_4)=(0.208887289734,0.049311600013)$ and $(-0.0797878448606,0.337986734608)$, and $C_0=0.117851130198$ for the singlet; the full 64-component vector equals the chain's $\Psi_0$ to $3\cdot10^{-77}$ (octets), $6\cdot10^{-100}$ ($(4,1,1)$), $3\cdot10^{-102}$ (singlet).
The result is \emph{$\phi$-independent}: repeating the whole procedure with $\phi'=(-2/5,3/8,7/11)$ changes $N^{(n)}$, $\hat c$, the Gauss parameters, $R_B$, $X_B$ and $m_B$ (Table~\ref{tab:RB}) but the recovered $C_\nu$ differ by $\le2\cdot10^{-77}$ (0 for the one-dimensional sectors, $10^{-202}$ for the singlet).\footnote{Agent U2, \texttt{RESULTS-\allowbreak{}untwist-\allowbreak{}limit.\allowbreak{}md} headlines 1--4, L2(d).
Normalisation-free check: $\rho_B(\hat\omega;\hat\omega^0)=[\hat c(\hat\omega^0)/\hat c(\hat\omega)]L(\hat\omega)/L(\hat\omega^0)$ against the exact ratio from $\langle\varphi_u\otimes e_1|\Psi_0\rangle$ at all 16 grid points: maximal deviation $2.5\cdot10^{-153}$, $3.1\cdot10^{-52}$, $4.4\cdot10^{-51}$, $3.0\cdot10^{-64}$, $1.2\cdot10^{-64}$, $3.4\cdot10^{-82}$ ($\phi$).}

The recipe for the CG generating function in ``CG from Q'' form is thus: write $\mathrm{GF}(\hat\omega)=\sum_{d+n+m=6-k_B}G^{s[d]}_{ij}\nu_i^{(n)}(\theta_1)\nu_j^{(m)}(\theta_2)/\hat c$ and expand in the basis $\hat c\,[e_1^T\rho(n'(\hat\omega))]_\nu$ (polynomials; the weight-$(0,2,1)$ component $-2a'g'$ carries $(\phi_2-\hat\omega_1')^{-2}$ that $\hat c$ does not clear, harmless because $C_3=\langle e_3\otimes e_1|\Psi_0\rangle=0$ on every HW ket of the sectors involved); the overall constant is fixed by the G-norm and the sign convention.
The final formula needs exactly: (i) the kernel $G$ split by degree; (ii) the one-site numbers $\nu(\theta_s;\epsilon)=(1,a_s/3,\pi,a_s\pi/3,\varepsilon_s)$, that is $(H_s,a_s)$ per site and branch as $\epsilon$-series to the depths above; (iii) the Gauss parameters (\ref{eq:gauss}) and $\hat c$. The Q-polynomials are not needed beyond these numbers.

\section{The \texorpdfstring{$\epsilon$}{}-series of the Laurent data}\label{sec:series}

\textbf{Punchline.} The Laurent data of the HW-limit branches have exact $\epsilon$-series, polynomial in the direction $\phi$ (degree $n$ at order $\epsilon^n$): linear for the top state, algebraic with a square root for $(4,1,1)$ and $(3,3,0)$, given to $n=6$ for the singlet and to $n=5$ over $\mathbb{Q}(\sqrt5)$ for the octets.
They obey $H_1+H_2=\epsilon\sum_{i\ne j}\phi_i\mu_j$ and $b_s=a_s(H_s-a_s)/3$ order by order.

\subsection{The series}\label{sec:seriesdata}

In units $\theta_{12}=1$ (at general $\theta_{12}$, $H^{(n)},a^{(n)}\sim\theta_{12}^{n-1}$ and $b^{(n)}\sim\theta_{12}^{n-2}$, verified at two $\theta$), with $\epsilon$ written $e$ where terse:

\begin{itemize}
\item \emph{Top, $(4,2,0)$, exact to all orders:} $H_1=4+\epsilon(\phi_1+2\phi_2+3\phi_3)$, $a_1=a_2=3\epsilon\phi_3$ ($\kappa=\zeta_3$, $P^3=e^{\zeta_3u}$), and $\Psi_B=e_1\otimes e_1$ exactly for all $\epsilon$. With $\theta_{12}$: $H_1^{(0)}=4/\theta_{12}$.
\item \emph{$(4,1,1)$}, with $x=\epsilon(\phi_2-\phi_3)$: $H_1=6+\tfrac\epsilon2(2\phi_1+5\phi_2+5\phi_3)+\bigl(\sqrt{1+x^2/4}-1\bigr)$, $a_1=6+\tfrac32\epsilon(\phi_2+\phi_3)+3\bigl(\sqrt{1+x^2/4}-1\bigr)$ (series to $n=6$: $x^2/8-x^4/128+x^6/1024$, resp. $3x^2/8,\dots$; odd orders $n\ge3$ vanish).
\item \emph{$(3,3,0)$}, with $y=\epsilon(\phi_1-\phi_2)$: $H_1=6+\tfrac32\epsilon(\phi_1+\phi_2+2\phi_3)+\bigl(\sqrt{1+y^2/4}-1\bigr)$, $a_1=a_2=3\epsilon\phi_3$ exactly.
\item \emph{Singlet $(2,2,2)$}, with $s_1=\phi_1+\phi_2+\phi_3$, $Q=\sum\phi_i^2-\sum_{i<j}\phi_i\phi_j$, $D=(\phi_1-2\phi_2+\phi_3)(\phi_1+\phi_2-2\phi_3)(2\phi_1-\phi_2-\phi_3)$:
\end{itemize}

\begin{equation}
\begin{aligned}
H_1&=12+2s_1\epsilon+\tfrac{Q}{6}\epsilon^2+0\cdot\epsilon^3-\tfrac{Q^2}{288}\epsilon^4+0\cdot\epsilon^5+(\deg6)\,\epsilon^6,\\
a_1&=12+s_1\epsilon+\tfrac Q4\epsilon^2+\tfrac D{72}\epsilon^3-\tfrac{Q^2}{192}\epsilon^4-\tfrac{11}{10368}DQ\,\epsilon^5+(\deg6)\,\epsilon^6.
\end{aligned}\label{eq:singlet}
\end{equation}

\begin{itemize}
\item \emph{Octets $(3,2,1)$ copy 0} (the other copy has $\sqrt5\to-\sqrt5$; even orders are rational, odd orders carry $\sqrt5$), in terms of $p_i=\phi_i$: $H_1^{(0)}=9$,
\end{itemize}

\begin{equation}
\begin{aligned}
H_1^{(1)}&=\frac{(15+\sqrt5)p_1+(20-2\sqrt5)p_2+(25+\sqrt5)p_3}{10},\\
H_1^{(2)}&=\frac{2(2p_1^2-3p_1p_2-p_1p_3+3p_2^2-3p_2p_3+2p_3^2)}{75},
\end{aligned}
\end{equation}

$H_1^{(3)}=-\sqrt5(p_1-p_3)^2(p_1-2p_2+p_3)/375$, $a_1^{(0)}=(15-3\sqrt5)/2$, $a_1^{(1)}=[(3-\sqrt5)p_1+(9-3\sqrt5)p_2+(18+4\sqrt5)p_3]/10$, $b_1^{(0)}=3\sqrt5$, $b_1^{(1)}=(12p_1+11p_2+22p_3)/5$; orders up to $n=5$ in the data (Appendix~\ref{app:D}).

The structure is exact, not guessed: (i) the \emph{global rules} $H_1+H_2=\epsilon\sum_{i\ne j}\phi_i\mu_j$ (it holds at order 1 only) and $b_1+b_2=\epsilon^2(\phi_1\phi_2\mu_3+\phi_1\phi_3\mu_2+\phi_2\phi_3\mu_1)$; (ii) $b_s=a_s(H_s-a_s)/3$ order by order; (iii) the six no-log conditions hold to $10^{-148}$. Algebraic degree: on $(4,1,1)$ and $(3,3,0)$ the $\sqrt{1+x^2/4}$ is read off the series to $n=6$ (not proved) and suggests a $2\times2$ mixing, the HW-limit branch mixing with one neighbour; on the singlet and octets we did not identify closed algebraic forms.

\subsection{How they were obtained}\label{sec:seriesmethod}

Two independent routes agree exactly.
Agent W4 did exact Rayleigh--Schrödinger degenerate perturbation theory on the weight-$\mu$ space, with isolating operator $H_1+a_1/7$ (all six Laurent operators commute, and every series eigen-equation holds exactly): series to order $7$ at $\phi$ and $\phi'$ (data file record \texttt{E4\_\allowbreak{}SERIES}), symbolic $\phi$ to order $3$, and $\Psi_n$ to order $4$ symbolically.
Agent U1 solved the six no-log conditions plus the two global rules order by order: that works (exact symbolic $\phi$ to $n=6$) when the Jacobian of the eight equations in six unknowns has rank 6, which is the case for $(4,1,1)$, $(3,3,0)$, $(2,2,2)$. For the top state and the octets the Jacobian has rank 5 (kernel dimension 1), the implemented degenerate-solvability step failed (a cubic with no rational root), and U1 used instead the closed form from the Wronskian structure (top) and a polynomial fit in $\phi$ over 29 generic $\phi$ at 200--300 digits followed by exact $\mathbb{Q}(\sqrt5)$ reconstruction by PSLQ (octets, $n\le5$; $n=6$ would need $\ge30$ directions).
All series were validated against 200 and 400-digit numerics at $\phi=(1/3,-2/7,5/11)$ to $10^{-60}$ ($n\le5$).\footnote{Agent U1, \texttt{RESULTS-\allowbreak{}untwisted-\allowbreak{}Q.\allowbreak{}md} T2(b)/T3 (\texttt{untwisted\_\allowbreak{}Q\_\allowbreak{}laurent\_\allowbreak{}exact.\allowbreak{}pkl}, $n\le5$; \texttt{untwisted\_\allowbreak{}Q\_\allowbreak{}pert\_\allowbreak{}sym.\allowbreak{}pkl}, $n\le6$; \texttt{untwisted\_\allowbreak{}Q\_\allowbreak{}octet\_\allowbreak{}fit.\allowbreak{}pkl}); agent W4, \texttt{Gaudin-\allowbreak{}untwisted-\allowbreak{}exact.\allowbreak{}wb} E4, data file \texttt{ref/\allowbreak{}untwisted-\allowbreak{}exact-\allowbreak{}210.\allowbreak{}txt}. U1 and W4 agree exactly to the depth where both have data; the degenerate solver is the missing piece (work package C1 of round 5).}

Tables~\ref{tab:series-ex} and~\ref{tab:series-num} (Appendix D) list the series at $\phi$. A few exact entries as a sanity check: top $H_1=140/31+(260/231)\epsilon$, $a_1=(15/11)\epsilon$, $b_{1,2}=\pm(700/341)\epsilon-(25/231)\epsilon^2$; $(4,1,1)$ at $\phi$: $H_1^{(1)}=349/462$, $a_1^{(1)}=39/154$, $H_1^{(2)}=100719/1660120$; singlet $H_1^{(1)}=232/231$ and $a_1^{(1)}=116/231$ ($2s_1$ and $s_1$ at $\phi$, in units $\theta_{12}=31/35$, since $2(1/3-2/7+5/11)=232/231$).

\subsection{The perturbative content: \texorpdfstring{$\Psi_1$}{} and \texorpdfstring{$X_B$}{XB}}\label{sec:Psi1}

The exact perturbation theory also gives the irreducible content of the first-order vector.
For $(4,1,1)$: $\Psi_1=\tfrac{31}{140}(\phi_3-\phi_2)(1,1)$ lies entirely in the $(4,2,0)$ irrep (its G-norm$^2$ is $(\theta_{12}^2/8)(\phi_2-\phi_3)^2$ while $\Psi_0=(1,-1)$ has G-norm$^2$ 2); $(3,3,0)$ is the same with $\phi_2-\phi_1$, entirely in $(4,2,0)$; the singlet $\Psi_1$ lies entirely in the $(3,2,1)$ octet irrep (G-norm$^2$ $(15376/1225)Q$, no $(3,3,0),(4,1,1),(4,2,0)$ component); the octets have components in all four irreps, proportional to $(\phi_2-\phi_3)$ and $(\phi_1-\phi_2)$ combinations.
The leading coefficient of the reference overlap, closed-form $X_B=\langle x(\phi)\otimes x(\phi)|\Psi_{k_B}\rangle$, with $\Psi_0$ of G-norm 1:

\begin{equation}
\begin{aligned}
X_{(4,2,0)}&=(\phi_1-\phi_2)^2,\qquad X_{(4,1,1)}=-\tfrac{\theta_{12}}2(\phi_1-\phi_2)(\phi_1-\phi_3)(\phi_2-\phi_3),\\
X_{(3,3,0)}&=-\tfrac{\theta_{12}}2(\phi_1-\phi_2)^3,\qquad X_{(2,2,2)}=\Bigl(\tfrac{\theta_{12}}2\Bigr)^4\bigl[(\phi_1-\phi_2)(\phi_1-\phi_3)(\phi_2-\phi_3)\bigr]^2,\\
X_{\rm octet}&=-\frac{5766(\sqrt5-3)}{6125}(\phi_1-\phi_2)^2(\phi_1-\phi_3)(\phi_2-\phi_3)\quad\text{and its conjugate,}
\end{aligned}\label{eq:XBforms}
\end{equation}

(for $(4,1,1)$, $(3,3,0)$ the formulas are for $\Psi_0$ of G-norm$^2$ 2, divided by $\sqrt{g_0}$ in $\langle x_0^G|\Psi_B\rangle=\epsilon^{2+k_B}X_B/\sqrt{g_0}$; for the octet $g_0=1080(7-3\sqrt5)$, for the singlet $g_0=288$; the data file has $\texttt{E5\_KB}$ with the exact strings).
These are the closed forms of U2's numerical $X_B$ up to normalisation: e.g. at $\phi$ the numbers are $0.383220,-0.017394,-0.074288,0.0013930,6.9931\cdot10^{-6}$.

The generating functions in the data file are $\hat c\langle\varphi_u\otimes e_1|\Psi_B\rangle=\sum_\nu W_\nu s[\nu]$ with $W_\nu=\hat c[e_1\rho(n')]_\nu$ polynomial in $(\hat\omega_1,\hat\omega_1')$ ($\nu=1..6,8$; $W_7$ has a simple pole $1/(\hat\omega_1'-\phi_2)$, which never carries a slice) and $s[\nu]=\Psi(\nu,1)$, exactly the CG coefficients at order $0$ and exact linear or quadratic forms in $\phi$ at orders $1,2$ (record \texttt{E5\_\allowbreak{}GF} of \texttt{untwisted-\allowbreak{}exact-\allowbreak{}210.\allowbreak{}txt}).

\subsection{Worked numbers for \texorpdfstring{$\mathrm{adj}\otimes\mathrm{adj}$}{adjadj}}\label{sec:worked}

The complete exact CG table of $\mathrm{adj}\otimes\mathrm{adj}$ is Table~\ref{tab:cg} (rows with $\Lambda=(2,1,0)$): $C_1=1$ for $(4,2,0)$; $C_2=-1/\sqrt2$ for $(4,1,1)$; $C_3=-1/\sqrt2$ for $(3,3,0)$; $C_8=2/\sqrt{288}=\sqrt2/12$ for the singlet $(2,2,2)$; and for the octet copy 1,

\begin{equation}
C_4=\frac{21-9\sqrt5}{\sqrt{7560-3240\sqrt5}},\qquad C_5=\frac{-3+3\sqrt5}{\sqrt{7560-3240\sqrt5}},
\end{equation}

copy 2 with $\sqrt5\to-\sqrt5$ (indices in notebook order: $\nu=4,5$ are the two weight-$(1,1,1)$ states).
Their Laurent data at $\zeta=0$ and the first corrections are in Tables~\ref{tab:untw} and~\ref{tab:series-ex}.
What the Q-side formula needs for each branch is: $G^s$, the series of $(H_s,a_s)$ to the depths of Table~\ref{tab:nudepth}, the Gauss parameters and $\hat c$; the output is the polynomial in $(\hat\omega_1,\hat\omega_1')$ of which $C_\nu$ are the coefficients up to the overall constant $m_B/X_B$ (Table~\ref{tab:RB}), fixed by the G-norm.
A single end-to-end script reproducing the whole table from the pieces ($G$, the series, the Gauss parameters) was not written; it is item F3 of the round-5 plan.

\section{The untwisted model}\label{sec:untwisted}

\textbf{Punchline.} At $\zeta=0$ the Laurent data of a HW state are Casimir eigenvalues: $H_1$ from $C_2$, $a_1-a_2$ from the cubic Casimir, and $b_s=0$ except in multiplicity-2 sectors.
The untwisted oper is Fuchsian at infinity with three polynomial solutions of degrees $(\mu_3,\mu_2+1,\mu_1+2)$, the Weyl-shifted weight; it is a ${}_3F_2$ equation only on branches with a consecutive exponent pair.
The monic twisted Q-functions have poles in $\epsilon$; the finite objects are the Plücker monomials and the Laurent data.

\subsection{Laurent data from Casimirs}\label{sec:casimir}

On the HW kets of weight $\mu$ in $V_\Lambda\otimes V_\Lambda$ at $\zeta=0$ all six Laurent operators act within the HW space, and we can compute their eigenvalues in closed form.
With $T=|\Lambda|$ and the Casimirs of the GT normalisation $E_{ik}E_{ki}$ and $E_{ij}E_{jk}E_{ki}$,

\begin{equation}
C_2(l)=\sum_il_i(l_i+4-2i),\qquad C_3(l)=l_1^3+l_2^3+l_3^3+4l_1^2+l_2^2-2l_3^2-l_1l_2-l_1l_3-l_2l_3+4l_1-2l_2-2l_3,\label{eq:Cas}
\end{equation}

one has $H_1=-H_2=\eta/\theta_{12}$, $a_1=-a_2=\alpha/\theta_{12}$, $b_1=-b_2=\beta/\theta_{12}^2$ with $\eta,\alpha,\beta$ independent of $\theta$ (verified at three $\theta$ points), and

\begin{equation}
\eta=T^2-\tfrac12\bigl[C_2(\mu)-2C_2(\Lambda)\bigr],\qquad
\alpha=\frac{4T(1+T)^2-3(1+T)C_2(\mu)+C_3(\mu)}6=\frac{(\mu_1+2)(\mu_2+1)\mu_3}2,\label{eq:etaalpha}
\end{equation}

(the first from $\sum_{ik}E^1_{ik}E^2_{ki}=\tfrac12(C_2^{\rm tot}-C_2^{(1)}-C_2^{(2)})$, as in (\ref{eq:H1}); the second found by an exact fit over 30 multiplicity-one sectors and then verified on every sector, and agreeing with U1's independently found product form over 23 sectors).
The product is half the product of the entries of the Weyl-shifted weight $(\mu_1+2,\mu_2+1,\mu_3)$: $\alpha=0$ iff $\mu_3=0$, $\alpha=\eta$ for $\mu_3>0$ except when $\lambda_2\ge2$ where $\alpha-\eta=4,6,6,10$ for $[2,2,0]\,(4,2,2)$, $[3,2,0]\,(6,2,2)$, $(4,4,2)$, $(4,3,3)$. For the adjoint $\eta=4,6,6,9,12$ and $\alpha=0,6,0,15/2,12$ on $(4,2,0),(4,1,1),(3,3,0),(3,2,1),(2,2,2)$. In terms of the pair ($\mu_3$, cubic Casimir) this is a ``determinant'' of the shifted weight: a cubic Casimir quantity.

In \emph{multiplicity-1} sectors $\beta=0$. In the multiplicity-2 sectors (the only ones that occur) the two copies are the eigenvectors of the $2\times2$ matrix of $S3c[1,2]=a_1$ (Table~\ref{tab:univ}): $\alpha_\pm=(\mathrm{tr}\mp\sqrt{\mathrm{disc}})/2$, $\beta=\pm\sqrt{\mathrm{disc}}$ with

\begin{equation}
\mathrm{tr}=\eta+2c_2,\qquad\mathrm{disc}=\eta^2-4c_2^2\qquad(7\text{ of the 8 multiplicity-2 sectors}),\label{eq:multlaw}
\end{equation}

giving $\mathrm{disc}=45$ for $[2,1,0]$ $(3,2,1)$; $105$ and $192$ for $[3,1,0]$ $(5,2,1)$, $(4,3,1)$; $105$ and $192$ for $[3,2,0]$ $(5,4,1)$, $(5,3,2)$; and $\beta_{\rm copy}=\alpha_{\rm copy}(\eta-\alpha_{\rm copy})/c_2$ (the no-log relation).
\textbf{There is one exception: $[3,2,0]$ $(5,3,2)$}, where $\mathrm{tr}=56=\eta+2c_2+18$ and $\mathrm{disc}=192$ (the prediction would have been $228$); there the no-log relation in this form fails.
No closed form for $\beta(\Lambda,\mu)$ is known: it is not a function of $C_2,C_3$ alone but needs the multiplicity-space operator.
For the adjoint octet $\alpha_\pm=(15\mp3\sqrt5)/2$, $\beta=\pm3\sqrt5$, with eigenvalue equation $\lambda^2-15\lambda+45=0$. At the standard point this gives $a_1=4.68085\ldots$ and $12.25463\ldots$ for the two copies.

Two sum rules hold as operator identities on all of $V\otimes V$ at $\zeta=0$ (exactly, five or six representations, 37 sectors):

\begin{equation}
b_1+b_2=0,\qquad a_1+a_2+\theta_1b_1+\theta_2b_2=0,\qquad (H_1+H_2=0).\label{eq:sumrules0}
\end{equation}

Hence $S_2=O(u^{-2})$ and $S_3=O(u^{-3})$ as $u\to\infty$ (no constants at $\zeta=0$): the oper is regular singular at infinity.\footnote{Agents U1 (\texttt{RESULTS-\allowbreak{}untwisted-\allowbreak{}Q.\allowbreak{}md} T0, Fractions on exact HW kets, five representations, two $\theta$ points) and W4 (\texttt{Gaudin-\allowbreak{}untwisted-\allowbreak{}exact.\allowbreak{}wb} E2, all 37 sectors of six representations; the 45 branches in Table~\ref{tab:untw}).
W4's exact tables: $\alpha_{\rm copy}$ copy 1 is the smaller, ``(3,2,1)'': $\alpha_\pm=(15\mp3\sqrt5)/2$.} The exact-data file lists all of it (Appendix~\ref{app:E}).

\subsection{Fuchsian at infinity: three polynomial solutions}\label{sec:fuchs}

The untwisted dual oper $\partial^3-S_1\partial^2+S_2\partial-S_3$ has, at $\theta_s$, exponents $\{0,\lambda_2+1,\lambda_1+2\}$ (Section~\ref{sec:frob}), and at infinity solutions growing as $u^k$, where for the adjoint $k$ solves $k^3-9k^2+(14+\eta)k-2\alpha=0$. The indicial roots at infinity are, on all 45 branches,

\begin{equation}
k=(\mu_3,\ \mu_2+1,\ \mu_1+2),\label{eq:expinf}
\end{equation}

the Weyl-shifted weight (e.g. top $(0,3,6)$, $(4,1,1)$: $(1,2,6)$, singlet $(2,3,4)$).
The Fuchs relation reads: twice the sum of the three exponents at one $\theta_s$, minus $\sum k$, equals $3$ (for the adjoint $12-\sum k=12-(|\mu|+3)=3$).
No logarithms occur (all exponents differ by integers, the local monodromy at $\theta_s$ is trivial by the no-log conditions, and at infinity the exponents are integers), so \emph{all three} solutions are polynomials, of degrees $(\mu_3,\mu_2+1,\mu_1+2)$, with Wronskian proportional to $\Pi^3$ and apparent singularities at $\theta_1,\theta_2$. The exact dimension of polynomial solutions of degree $\le d$ is, for the top, $1,1,1,2,2,2,3$ ($d=0,\dots,6$), for $(4,1,1)$ $0,1,2,2,2,2,3$, and always reaches $3$. This refutes the expectation recorded in the round-4 plan that fewer than three polynomial solutions exist at $\zeta=0$ (the Wronskians of the twisted Q's degenerate as $\zeta\to0$, which suggested a degenerate solution space; Section~\ref{sec:poles} explains where the degeneracy really sits).

The exact echelon bases in $x=(u-\theta_2)/\theta_{12}$ (site 2 at $x=0$, site 1 at $x=1$), all verified to solve the oper in sympy:

\begin{equation}
\begin{aligned}
(4,2,0):&\ \bigl\{1,\ \tfrac12x^2(2x-3),\ \tfrac14x^4(2x-3)^2\bigr\}=\{1,F,F^2\},\ \ F=\textstyle\int\Pi\ (\text{a }\mathrm{Sym}^2),\\
(4,1,1):&\ \bigl\{\tfrac12(2x-1),\ x^2,\ \tfrac12x^4(2x^2-6x+5)\bigr\},\\
(3,3,0):&\ \bigl\{1,\ x^2(x-1)^2,\ \tfrac14x^2(4x^3-10x+5)\bigr\},\\
(2,2,2):&\ \bigl\{\tfrac14(2x-1)^2,\ \tfrac18(2x-1)(4x^2+2x-1),\ x^4\bigr\},\\
(3,2,1)_{\rm copy\,0}:&\ \bigl\{\tfrac1{10}(10x-5+\sqrt5),\ \tfrac12x^2(2x-3-\sqrt5),\ \tfrac12x^4(2x-5+\sqrt5)\bigr\},
\end{aligned}\label{eq:v0bases}
\end{equation}

(copy 1: $\sqrt5\to-\sqrt5$; the copies are Galois conjugate).
For the top state the basis is $\{1,F,F^2\}$ with $F=\int\Pi$ (note the twisted top-state Q-functions are $e^{\zeta_3u}$ and $e^{\zeta_2u}(\Pi+\Pi'/d+\Pi''/d^2)$ with $d=\zeta_3-\zeta_2$).\footnote{Agent U1, \texttt{RESULTS-\allowbreak{}untwisted-\allowbreak{}Q.\allowbreak{}md} T1 and \texttt{v0\_\allowbreak{}bases()} in \texttt{untwisted\_\allowbreak{}Q.\allowbreak{}py}; W4 E3 (exponents and degrees on all 45 branches, \texttt{Gaudin-\allowbreak{}untwisted-\allowbreak{}exact.\allowbreak{}wb}).
The explicit statement for the top Q's was checked numerically at 200 digits by U1.}

\subsection{The \texorpdfstring{${}_3F_2$}{F} criterion}\label{sec:3F2}

A natural hope was that the untwisted oper is a hypergeometric ${}_3F_2$ equation, so that its polynomial solutions would be hypergeometric polynomials.
This is true only sometimes.
W4 compared the oper after Möbius maps and gauge factors with the ${}_3F_2$ equation for all role assignments of the three singular points $\theta_2,\theta_1,\infty$ and found: \textbf{the oper is a ${}_3F_2$ equation if and only if some singular point carries a consecutive pair of exponents}, which happens on exactly 20 of the 45 branches.
At $\theta_s$ the exponents $\{0,\lambda_2+1,\lambda_1+2\}$ contain a consecutive pair iff $\lambda_1=\lambda_2$ (pair $\{\lambda+1,\lambda+2\}$: $[1,1,0],[2,2,0]$, all five branches); at infinity iff $\mu_2=\mu_3$ or $\mu_1=\mu_2$. The parameters are $b=(-\lambda_2,-\lambda_1-1)$ with witness roles $s=0$ at $\theta_2$, $s=1$ at infinity, $s=\infty$ at $\theta_1$ and gauge $(1-x)^{\ldots}$ (for $\lambda_1=\lambda_2$: $s=x$, gauge $(1-x)^{\lambda+1}$), the parameters $a_j$ being tabulated by witness.
These are degenerate integer parameters, reducible, so ``hypergeometric'' is not the natural frame: the polynomial bases above are the closed forms.
On the other 25 branches (all with $\mu_1>\mu_2>\mu_3$ and $\lambda_1\ne\lambda_2$, including every multiplicity-2 sector) there is no consecutive pair at any of the three points, hence no ${}_3F_2$ form.
U1's adjoint verdict (only $(4,1,1)$, $(3,3,0)$, $(2,2,2)$, with $X=(u-\theta_2)/(u-\theta_1)$ and gauge $(u-\theta_1)^c$, $c=2,5,3$ or $4$) agrees with W4's.\footnote{Agent W4, \texttt{Gaudin-\allowbreak{}untwisted-\allowbreak{}exact.\allowbreak{}wb} E3; Table~\ref{tab:untw} last column.
We re-counted the criterion from the exponents in the 45 branches of the data files and obtained 20.}

The untwisted \emph{direct} oper (our $\Pi^p\circ$adjoint$\circ\Pi^{-p}$) has polynomial solutions of degrees $\{2\lambda_1+2-\mu_3,\,2\lambda_1+1-\mu_2,\,2\lambda_1-\mu_1\}$ on 43 of the 45 branches; the exceptions are $[1,1,0]$ $(2,2,0)$ (degrees $\{0,2,4\}$) and $[2,2,0]$ $(4,4,0)$ ($\{0,2,6\}$).

\subsection{Why the twisted monic Q's have poles in \texorpdfstring{$\epsilon$}{}}\label{sec:poles}

We expected the untwisting limit of the twisted Q's to be the three untwisted polynomials.
The monic $p_j(u;\epsilon)$ have \emph{no} coefficientwise limit: the coefficient $c_m$ of $p_j$ behaves as $\epsilon^{-\min(\mu_j-\mu_3,\ \mu_j-m)}$ (top: $P^1$ coefficients $\epsilon^{-4..0}$, $P^2$ $\epsilon^{-2..0}$, $P^3$ finite; $(4,1,1)$: $P^1$ at $\epsilon^{-3}$; $(3,3,0)$: $P^1,P^2$ at $\epsilon^{-3}$; $(3,2,1)$: $P^1$ at $\epsilon^{-2}$, $P^2$ at $\epsilon^{-1}$; singlet finite).
Expanding in the exact untwisted basis $v_a(\epsilon)$ (solutions of the twisted oper with the same jets at $u=0$), $P^j=\sum_aM_{ja}(\epsilon)v_a(\epsilon)$, the matrix $M$ has pole order $\mu_j-\mu_3$ on $v_3$ and $O(\epsilon^{-(N_j-1)})$ on $v_2,v_1$, with $\det M=c_W\,\Pi(0)^3\sim\epsilon^3$ exactly; the normalised rows $(1,\rho_2,\rho_1)_j$ have $\rho=\epsilon r^{(1)}+\cdots$ with the \emph{first-order direction common to all $j$}, differing only at $\epsilon^2$ (octets, singlet) to $\epsilon^4$ (top: $R_1-R_3=O(\epsilon^4)$, $R_2-R_3=O(\epsilon^3)$).
So the degeneracy sits in the monic normalisation, not in the solution space: the limiting 3-plane is the span of the untwisted polynomials and the $\zeta$-jet basis is the Smith (saturation) basis of $M$; its elementary-divisor valuations were not computed, and no closed form for $r^{(1)}$ was found.
Limits of $\epsilon$-differences and derivatives: the leading Laurent coefficients of $P^j$ solve the untwisted oper, the lower ones solve $D_0p^{(n+1)}=-D_1p^{(n)}-\cdots$ (inhomogeneous), consistent with the numerics.
The finite objects are the Plücker monomials and the Laurent data, which is why the formulation of Section~\ref{sec:nu} is the right one.

\section{Surprises and corrections along the way}\label{sec:surprises}

\textbf{Punchline.} About two dozen of our own expectations were wrong or had to be sharpened.
We list them candidly, with where each was caught, because the pattern (conflating $|\Lambda|$ with $c_2$, expecting degeneracy where there was none, expecting too many monomials) tells us where to be careful with $\lambda_1\ge3$ and $\lambda_3>0$.

{\footnotesize
\begin{longtable}{p{0.3cm}p{5.2cm}p{5.8cm}p{2.6cm}}
\caption{Corrections to the plans, in the order they were found.
``Round'' is the plan round that carried the wrong expectation; the correcting agent is in the last column.}\label{tab:surprises}\\
\toprule
\# & expected & found & by \\
\midrule
\endfirsthead
\toprule
\# & expected & found & by \\
\midrule
\endhead
\bottomrule
\endlastfoot
1 & $H_1=+[\ldots]$ (plan, round 1) & sign is $-$: $H_1=-[\sum\zeta_iE_{ii}+\sum E^1_{ik}E^2_{ki}/\theta_{12}]+|\Lambda|\sigma+|\Lambda|^2/\theta_{12}$, $H_1+H_2=-W+2|\Lambda|\sigma$ & W1 \\
2 & $x_0$ (XXX limit) not parallel to the Gaudin $x_0^G$ & exactly parallel ($10^{-31}$), $x_0\sim h$ per site, $x$ linear in $z$ at all orders & P2, D2 \\
3 & twelve $P^{3(c)}p_{kl}$ monomials per site & five (adjoint); in general $\binom{\lambda_1+2}2-\binom{\lambda_2+1}2$ with the lone factor always $P^3$ & D1, D2, P4 \\
4 & $M_{\rm vac}=\tfrac52h^2W'P^3$ & $2h^2W'P^3$ (the weight-2 minor contains an $n=1$ piece fixed by $\mathrm{minor}(-1,0)=0$) & D2 \\
5 & $S_2$ double pole $=|\Lambda|$ (or Casimir$/2$) & $c_2=\lambda_2(\lambda_1+1)$ & P3, then P4 \\
6 & $a_s^{\rm vac}=|\Lambda|\omega_3$, $f(S3c)=|\Lambda|\zeta_3$ & $c_2\omega_3$, $c_2\zeta_3$ (equal on $[2,1,0]$, $[3,1,0]$) & P4, W3 \\
7 & the direct oper is the formal adjoint & the adjoint has no $e^{-\zeta u}\times$polynomial solutions; direct $=\Pi^p\circ$adjoint$\circ\Pi^{-p}$, $p=\lambda_1-\lambda_3$ & P1, P2, P3 \\
8 & dual Wronskian $F'=\Pi$ for $[3,2,0]$ & $F'=\Pi^{\lambda_2}=\Pi^2$; direct exponents $\{0,2,5\}$, not $\{0,3,5\}$ & P4 \\
9 & $c_H$ has $1/\theta_{12}$ coefficient $2\lambda_1$ & $2e_2(\Lambda)$ ($2\lambda_1\lambda_2$ for $\lambda_3=0$) & W2 \\
10 & the $[3,1,0]$ fourth localiser is a single product, $S3c[s,1]=c_\Lambda S3c(H-S3c)$ & two quadratic monomials plus a $W$ term; for $[3,2,0]$ no $c_\Lambda$ at all, $\gamma$ changed $2\to1$ & W2, W3 \\
11 & $\ell_1=1$ risk: no localiser inside $D_2$ & $\ell_1=3$ already inside $D_1$; every Laurent coefficient localises & W1 \\
12 & $\omega_1'$ fixed control fails in Gaudin as in XXX & passes at finite $\omega_2$, fails at $\omega_2=\infty$ & W1, W2 \\
13 & fixed-coefficient orthogonality among brackets & the relation has state-dependent $\Delta$ (Lagrange identity); six residue columns & P1 \\
14 & $\zeta$-independent kernel with vacuum monomial $M=Q_1P^3$ & $M=P^3p_{0,\lambda_2+1}$ and the scalar $c$ are needed; direct-$Q_1$ bases carry a hidden $\sigma$ & D1 \\
15 & naive $t\to\infty$ at fixed $a$ gives the Gaudin kernel & gives a rank-one state-independent $(1-a)^3(1-a')^3$; wrong limit & D2 \\
16 & the SoV covectors degenerate at $h\to0$ in the companion frame & they are independent there; the degeneracy lives in the diagonal frame ($P^{-1}_{\rm phys}=O(h)$) & D3 \\
17 & one polynomial Gaudin $B_G(u)$ & none; solvable hierarchy, no polynomial of low degree & D3 \\
18 & cleaning needed for $(5,3,2)$ etc. & no cleaning anywhere ($K/S=O(\epsilon)$ on all branches) & W1--W3 \\
19 & top-state sum rule $=M^{\rm top}M^{\rm top}$ & extra factor $R(a)R(a')$, $R=(1-\zeta_2a)^{\lambda_1-\lambda_2}(1-\zeta_3a)^{\lambda_1}$ & G1 \\
20 & $\lambda_2=0$: $S_2$ double pole non-scalar & scalar $0$; one-row reps are the master kernels & G1 \\
21 & $\le2$ polynomial solutions of the untwisted oper & all three are polynomials, degrees $(\mu_3,\mu_2+1,\mu_1+2)$ & U1, W4 \\
22 & the untwisted oper is $_3F_2$ & only with a consecutive exponent pair (20 of 45 branches) & U1, W4 \\
23 & the monic twisted Q's have untwisted limits & poles in $\epsilon$; finite objects are Plücker monomials and Laurent data & U1 \\
24 & $\hat c$ clears all denominators of $e_1\rho(n')$ & not the weight-$(0,2,1)$ component $-2a'g'$ with $(\phi_2-\hat\omega_1')^{-2}$ (harmless) & U2 \\
25 & $\det N$ rank 5 for all $A\ne B$ & rank 4 iff $\Delta\mu$ has one positive and two negative components & P3 \\
26 & gauge power fixed by $\lambda_3=0$ data & $p=\lambda_1-\lambda_3$ vs $\lambda_1$ not separated (all data $\lambda_3=0$) & P3, P4 \\
27 & octet copies distinguished only at higher order & they split already at $\epsilon^0$ by $a_s^{(0)}$ (while $H_s^{(0)}$ coincide) & U2 \\
\end{longtable}}


A few of these deserve a sentence more than the table gives.
The first plan said that with the XXX-like ansatz the Gaudin formula would involve twelve monomials and $M_{\rm vac}=\tfrac52h^2W'P^3$ (items 3, 4); both numbers came from the heuristic that the XXX ladder points $\theta_s+mh$ become the order-3 jet, and in fact the true count follows from the Frobenius structure of the oper.
The \emph{cleaning question} (item 18) was Paul's own: whether the per-sector, per-copy cleaning of the single-vacuum approach can be avoided; the plan's careful list of expected failures ($(5,3,2)$, multiplicity 3) is a good example of a plan that was too pessimistic.
The conflation of $|\Lambda|$ and $c_2$ (items 5, 6) is the one to remember: on $[2,1,0]$ and $[3,1,0]$ one has $\lambda_2=1$ and hence $c_2=\lambda_1+1=|\Lambda|$, so the two quantities coincide and only $[3,2,0]$ could tell them apart.
Finally, item 21 is a good reminder that exponents integer-spaced with trivial monodromy mean polynomial solutions: the ``degeneracy'' of the Wronskians at $\zeta=0$ does not reduce the solution space.

\section{Outlook: the round-5 plan}\label{sec:outlook}

\textbf{Punchline.} The adjoint works end to end.
Round 5 aims to turn ``the adjoint works'' into ``the formula'': a closed form for the master kernel $R^{\lambda_1}$, the Laurent-data form and untwisting limit for $[3,1,0]$ and $[3,2,0]$, a derivation of the $\epsilon$-series, the norm of the untwisted eigenvectors from the Q side, and the write-up.
Nothing below is implemented; Paul asked for the plan, not for work.\footnote{\texttt{Paul/\allowbreak{}Mathematica/\allowbreak{}Gaudin/\allowbreak{}Gaudin-\allowbreak{}sl3-\allowbreak{}round5-\allowbreak{}plan.\allowbreak{}md} (written 2026-10-16 about 18:00, not dispatched).
Suggested allocation is sonnet agents only, $\le4$ Python agents at once, the wolfbook builder alone on the kernel K1.}

\subsection{WP-A: the master kernel in closed form}\label{sec:WPA}

The goal is a four-index closed form of $R^{\lambda_1}_{(kl),(k'l')}(x,y)$, $x=at$, $y=a't$, valid for all $\lambda_1$; then every $\Lambda$ follows by the block theorem and $\theta\otimes\theta'$. Data in hand: exact $R$ for $\lambda_1\le4$ ($[4,0,0]$: 225 entries), the first-order formula, the single-index $\Gamma$-form, the extreme-hop coefficients, the exchange rule, homogeneity, the rank pattern $1,2,4,8$ and the top-state sum rule.
Three attempts are planned: (A1) a fit-guided Pochhammer ansatz generalising the $\Gamma$-form, validated on a new representation, $[5,0,0]$ (1225 states), done at 40 digits through the $\nu$-form, which needs only Laurent data and no Q-solving, so that Section~\ref{sec:nu} pays off computationally; (A2) a structural derivation of the jet basis for one-row representations $V_{(\lambda_1,0,0)}=\mathrm{Sym}^{\lambda_1}(\mathbb{C}^3)$ in the polynomial realisation $E_{ij}=x_i\partial_j$, where the SoV wavefunction rule becomes explicit differential or integral operators (start with $\lambda_1=1,2$); (A3) the auxiliary $\mathfrak{sl}_2$: is $R$ covariant under a twisted shift $(a,a',t)\to(a+s,a'+s,t')$? If so $R$ is fixed by its value at $a=0$, that is by the one-site data plus finitely many coefficients; also whether $R^{\lambda_1}$ is the spin-$(\lambda_1+1)/2$ fusion of the $\lambda_1=1$ R-matrix form.
Acceptance: $R^{\lambda_1}$ reproduces all exact kernels for $\lambda_1\le4$ and predicts $[5,0,0]$ (and $[5,1,0]$) on $\ge200$ states to $\ge30$ digits.

\subsection{WP-B: Laurent-data form and the limit for \texorpdfstring{$[3,1,0]$}{[3,1,0]}, \texorpdfstring{$[3,2,0]$}{[3,2,0]}}\label{sec:WPB}

(B1) Derive $\nu_{kl}$ for general $(\lambda_1,\lambda_2)$ from the Frobenius recursion at $\theta_s$ (exponents $\{0,\lambda_2+1,\lambda_1+2\}$; the order-$(\lambda_1+1)$ jet lies in a 2-plane with coefficients rational in $(H_s,a_s,b_s)$), prove that $P^3D_s$ cancels between the sides, and verify on the 225 states of $[3,1,0]$ and $[3,2,0]$ to $10^{-50}$. (B2) Run U2's pipeline for both representations: $k_B$ per branch (10 each) and the depth; test whether $k_B$ follows a weight rule (adjoint $0,1,1,2,2,4$) and the exact CG tables of the data files. (B3) The exact $\epsilon$-series of the Laurent data for the $10+10$ branches (duality check: $[3,2,0]$'s series are the dual images of $[3,1,0]$'s).
Acceptance: all 20 CG tables from the Q side to $\ge30$ digits, $\phi$-independent.

\subsection{WP-C: the mechanism of the \texorpdfstring{$\epsilon$}{}-series}\label{sec:WPC}

(C1) Implement the degenerate solver: Fredholm alternative at each order for the rank-5 cases (top state, octets), symbolic in $\phi$, to $n=8$, from the Bethe side alone. (C2) Identify the $2\times2$ mixings behind the algebraic $(4,1,1)$, $(3,3,0)$ forms and the algebraic degrees on the others; conjecture for $H_s(\zeta)$ on a HW-limit branch as eigenvalues of a finite matrix built from the twisted $\mathfrak{gl}_3$ Casimirs. (C3) Closed form via Casimirs on the chain side: the Laurent operators on the $\mu$-isotypic component are expressible through $C_2$, $C_3$, $\sum\zeta_kE_{kk}$ and products; write them in that algebra and diagonalise symbolically on each isotypic component (small matrices since the twist mixes only within weight spaces).
Acceptance: an algorithm or formula producing the series of every branch of $[2,1,0]$, $[3,1,0]$, $[3,2,0]$ exactly.

\subsection{WP-D: the norm from the Q side}\label{sec:WPD}

The CG formula fixes the generating function up to one constant per branch, the G-norm.
Classical results say the norm of a Gaudin Bethe vector is the Hessian of the Yang--Yang function at the Bethe roots (Reshetikhin--Varchenko, Mukhin--Varchenko); here the Bethe equations are the no-log conditions.
Hypothesis: $\langle\Psi|\Psi\rangle_G/\langle x_0|\Psi\rangle^2$ is proportional to the determinant of the Jacobian of the six no-log conditions with respect to the Laurent data, or to the equal-twist $\det N$ suitably normalised (D3's observation that $Q_{\rm num}^{\rm eq}M_1M_2/\det N_B$ is state-independent is norm-like).
Test numerically on the 64 states, then take the untwisting limit; if none of the candidates is state-independent the true Yang--Yang Hessian in terms of Q-zeros is needed, a different formalism, and the plan says to flag it and stop.

\subsection{WP-E, WP-F: the write-up and housekeeping}\label{sec:WPE}

The notes you are reading are E1.
E2 is the paper: \texttt{paper/\allowbreak{}gaudinCG.\allowbreak{}tex} is a gl(4)-rectangular stub, to be replaced by the sl(3) Gaudin CG paper (recommended) or a new file; proposed structure: introduction; the two-site model, cdet, Laurent data, opers; vacua, frames and the two-vacua family (no cleaning); the jet formula (basis, K tensor, block theorem, one-site kernel, master kernel); the overlap as a function of Laurent data; the untwisting limit; worked example adj$\otimes$adj; outlook; appendices on conventions, kernels, series.
Housekeeping: rerun the four Gaudin notebooks in VS Code to store native outputs (the agents stored none), gitignore \texttt{*.\allowbreak{}pkl} ($\approx230$ MB in \texttt{Paul/\allowbreak{}Python/\allowbreak{}gaudin\_\allowbreak{}sl3/\allowbreak{}}) before committing, promote the notebooks to \texttt{Clean/\allowbreak{}} and run the docs agent, update \texttt{CLAUDE.\allowbreak{}md}'s outdated goal statement, a one-file demonstration \texttt{cg\_\allowbreak{}from\_\allowbreak{}Q.\allowbreak{}py}, and the backlog: $[4,2,0]$ (dim 27, multiplicity 3) as the last cleaning check, the $(5,3,2)$ exception of the multiplicity-2 law and $\beta(\Lambda,\mu)$, the second null vector of the rank-4 pairs, the $\lambda_3>0$ dependence of the gauge rule.

The decisions that are Paul's: the paper file; the priority between WP-A (the more striking result) and WP-C/D (the more complete one, no chain input); whether to spend a round on $[4,2,0]$ before writing; and the update of \texttt{CLAUDE.\allowbreak{}md}.

\appendix

\section{Conventions and normalisations}\label{app:conv}

This appendix reconciles the notations of the different agents and fixes every factor.
The main text uses one convention throughout (the canonical one below); each source file is mapped to it.

\subsection{The kernel normalisations}\label{app:norm}

{\small
\begin{longtable}{p{2.3cm}p{4.6cm}p{7cm}}
\caption{Kernel normalisations in the different sources and their relations. $c$ is the frame scalar (\ref{eq:cfactor}); $M=P^3p_{0,\lambda_2+1}$.}\label{tab:norms}\\
\toprule
symbol & where & definition and relation \\
\midrule
\endfirsthead
\toprule
symbol & where & definition and relation \\
\midrule
\endhead
\bottomrule
\endlastfoot
$K$ & G1, P4, main text & $c\,Q_{\rm num}M_1M_2=\sum K\,m\,m$, $m_{kl}=P^3p_{kl}$, $M$ exactly (no extra factor) \\
$G$ & D1, D2, D3 & $4c\,Q_{\rm num}M_1M_2=\sum G\,N\,N$ (D2's ``$4c$''); for the adjoint $G=4K$, because $M_{\rm vac}=2h^2P^3p_{02}$ per site in the XXX limit and $2\times2=4$ \\
$G^s$ & U2 & $G^s:=G/4=K$ ($c\,Q_{\rm num}M_1M_2=\sum G^sNN$; U2 states ``G of the spec = D2's printed G divided by 4'') \\
$\Theta$ & G1, D3 & $\Theta=K/(\omega_1\omega_1')^{2\lambda_1-\lambda_2}$, a polynomial in $(a,a',t)=(1/\omega_1,1/\omega_1',1/\theta_{12})$; adjoint: $\Theta=G/(4\omega_1^3\omega_1'^3)$ (D3, exact in 25/25 entries) \\
$Q_{\rm num}$ & all & $\langle\Phi_\infty|\Psi\rangle/\langle x_0^G|\Psi\rangle$ in the diagonal frame; companion-frame version $Q_{\rm num}^{\rm comp}=Q_{\rm num}/(\kappa^G\kappa^{G\prime})=c\,Q_{\rm num}/(\omega_1\omega_1')^{2\lambda_1-\lambda_2}$ (D1, D3) \\
$T(B)$ & D1 & $\langle\Phi_\infty|\Psi\rangle D_B(\omega_0)/\langle\Omega_0\otimes\Omega_0|\Psi\rangle=c_T(\zeta)\,Q_{\rm num}M_1M_2$ \\
$\nu$ & U2 & $N_i/M=(1,\kappa,\pi,\kappa\pi,\varepsilon_s)$; the main text writes $\nu$ for this vector, and $\nu(u)=\Pi(u)$ appears only in the XXX $B$-operator discussion (Section~\ref{sec:Bxxx}) \\
$D_2$ & W2, W3 & the quadratic localising algebra; not to be confused with D2 (agent) or $D_B$ ($\det N$) \\
$\gamma_2,\gamma_3,d_3$ & this text & Frobenius coefficients of $f_0,f_2$ (U2 calls $\gamma_2,\gamma_3$ ``$c_2,c_3$''; renamed to avoid the clash with the double pole $c_2$) \\
\end{longtable}}


The relation $\Theta^{\rm adj}=G/(4\omega^3\omega'^3)$ was found by D3 to be exact; G1 verified symbolically in 25/25 entries that $\Theta$ defined from $K$ equals it, so no further factor hides in the conventions.
The ``4 of D2/D3'' is a convention, not physics: with $M=P^3p_{0,\lambda_2+1}$ exactly one uses $K=G/4$ and no 4 appears anywhere.

\subsection{Frames and bases}\label{app:basis}

\begin{itemize}
\item \emph{Frames.} Diagonal frame: states of the diagonal-twist Hamiltonians; companion frame: $\Psi_c=(P_{\rm phys}^{-1}\otimes P_{\rm phys}^{-1})|_{z=\zeta}\Psi_d$, $x_0^c=e_1\otimes e_1$ (D3), $\Phi^\infty_c=\Omega_c^\infty[\omega_1]\otimes\Omega_c^\infty[\omega_1']$. $x_0^G$ in the diagonal frame is $x(\zeta)\otimes x(\zeta)$ and $\langle x_0^c|\Psi_c\rangle=\langle x_0^d|\psi_d\rangle$.
\item \emph{Basis orders.} Notebook (W-agents, data files): GT patterns descending, index 1 or 0 is the HW ket; \texttt{gl3\_\allowbreak{}reps} (P, D, G, U agents): lowest weight first, HW ket last.
Adjoint map $\mathrm{NB}\to\mathtt{gl3}$: $(7,5,6,4,2,1,3,0)$ (0-based).
Weights in notebook order for the adjoint: $(2,1,0),(2,0,1),(1,2,0),(1,1,1),(1,1,1),(1,0,2),(0,2,1),(0,1,2)$; the metric diagonal $G=(1,1,1,2,6,9,4,9)$. U2's slices (\texttt{gl3} index $\nu$): top $C_7$, $(4,1,1)\ C_5$, $(3,3,0)\ C_6$, octets $(C_2,C_4)$, singlet $C_0$; W4's (notebook, 1-based): $C_1,C_2,C_3,(C_5,C_4)$ resp. $(C_4,C_5)$, $C_8$; Table~\ref{tab:cg} uses W4's.
\item \emph{Signs.} HW kets are G-normalised with the first non-zero component, in descending GT order, positive (D1, U2); W4's canonical HW kets are primitive integer vectors (first non-zero component positive).
For $(4,4,2)$ of $[3,2,0]$ the exact slice has an exactly vanishing component and the chain branch's sign convention differs; the CG cells align signs and use norm-based errors (W3).
\item \emph{Copies.} In a multiplicity-2 sector copy 1 is the eigenvector of $S3c[1,2]$ with the smaller eigenvalue ($\alpha_1^{(0)}$); U1's ``copy 0'' is W4's copy 1 and $\sqrt5\to-\sqrt5$ exchanges the copies.
\item \emph{$\theta_{12}$ units.} U1's series are in units $\theta_{12}=1$, with $H^{(n)},a^{(n)}\sim\theta_{12}^{n-1}$, $b^{(n)}\sim\theta_{12}^{n-2}$; W4's data are $\theta$-scaled, $H=\eta/\theta_{12}$, $a=\alpha/\theta_{12}$, $b=\beta/\theta_{12}^2$, with $\theta=(2/7,-3/5)$, $\theta_{12}=31/35$ for the numerical tables and series (so ``$140/31$'' is $4/\theta_{12}$).
\item \emph{Kernel pair orders.} G1: lexicographic pairs $(kl)$, $0\le k<l\le\lambda_1+1$, $l\ge\lambda_2+1$. D2/D1: $N=(P^3p_{02},P^{3\prime}p_{02},P^3p_{03},P^{3\prime}p_{03},P^3p_{23})$; the map to G1's order $(02),(03),(12),(13),(23)$ is $1\to(02)$, $2\to(12)$, $3\to(03)$, $4\to(13)$, $5\to(23)$ because $P^{3\prime}p_{0l}=P^3p_{1l}$.
\item \emph{$c$ for general $\Lambda$:} (\ref{eq:cfactor}); wrong exponents tested and excluded (Section~\ref{sec:frames}).
\end{itemize}

\subsection{Numerical points and grids}\label{app:grids}

\begin{itemize}
\item $\theta=(2/7,-3/5)$, $\theta_{12}=31/35$; $\theta$ variants for kernel laws: $(1/5,-4/9)$, $(3/7,-1/4)$, $(1/2,-1/3)$, $(-1/7,-6/5)$, $(4/5,-2/3)$ ($\theta_{12}=29/45,19/28,5/6,37/35,22/15$); further $\theta_{12}$ values $7/9,13/8,-9/7,-17/10,11/6,3/4,-5/3,19/13,-7/11,9/5,-13/9,25/14,-19/10,15/11,-11/7$ (G1, 12--20 values); pairs with equal $\theta_{12}$: $(2/7,-3/5),(-1/3,-128/105),(5/4,51/140)$.
\item Twist points: $P=(3/2,5/4,2/3)$, $P'=(2,5/3,-1/2)$, $X_1=(7/5,-3/4,1/6)$, $X_2=(-4/3,2/9,5/7)$, $X_3=(1/3,-2/7,5/11)$; further $X_4..X_6$ and 26 real-spectrum random rationals (D1); some triples give a complex spectrum of $S_2+\pi S_3$ and are skipped.
\item Vacuum twists: $\omega=(5/2,7/3,3/5)$, $\omega'=(7/4,-2/5,9/7)$; $\omega_0=(5/2,7/3,3/5)$ for $T(B)$; $(\omega_1,\omega_1')$: base $(5/2,7/4)$, second $(-2/5,3)$, a $3\times3$ grid $\omega_1\in\{5/2,-2/5,9/7\}\times\omega_1'\in\{7/4,3,-5/3\}$, a 9-value grid $\{5/2,-2/5,9/7,3,-5/3,7/4,11/3,-7/2,13/5\}$ (72 ordered pairs), and G1's $w'$ list up to $23/9,-3/2,10/3$.
\item Regime: $\phi=(1/3,-2/7,5/11)$, control $\phi'=(-2/5,3/8,7/11)$; $\hat\omega_1\in\{3/4,-5/3,1/2,-7/3\}$, $\hat\omega_1'\in\{-12/5,7/2,1/8,5/4\}$ (the 24-point Stage-D grid is $\{3/4,-5/3\}\times\{-1/5,9/7,5/3\}$ for $\hat\omega$ and $\{-12/5,7/2\}\times\{1/8,-3/4\}$ for $\hat\omega'$); $\epsilon\in\{1,2,4\}\cdot10^{-6}$ (chain-side, three values) and $\epsilon_k=10^{-6}2^{-k}$, $k=0..11$ (series fits, $K=12$; the extended set $K=24$, 400 digits, reliable to $n=18$).
\item Generic points for operator identities: $u_0=3/7$, $u_1=-2/9$ (note: Laurent fits at $\theta_1=3/7=u_0$ hit a pole).
\item Precision: chain 200 digits; Q side 60 (finite twist), 80 (fits, kernels), 100--120 (regime), 200 (series), 400 (extended series); fixed-point 320-bit integer QR in D1; exact Fractions in D3, G1 (one-site).
\end{itemize}

\subsection{Vacuum eigenvalue functions}\label{app:vac}

The exact eigenvalue functions of the vacuum $\Omega_{\zeta\omega}\otimes\Omega_{\zeta\omega}$ (exact rationals; $x_s=u-\theta_s$), with $S_1=e_1(\omega)+|\Lambda|(1/x_1+1/x_2)$, $e_1(\omega)=163/30$ for $\omega$ and $369/140$ for $\omega'$, and the constants $e_2(\omega)=131/15$, $e_3(\omega)=7/2$ ($e_2(\omega')=29/28$, $e_3(\omega')=-9/10$).
The table lists the Laurent data of $S_2,S_3$ (the double poles are $c_2$ and $a_s=c_2\omega_3$).

{\small
\begin{longtable}{llccccc}
\caption{Vacuum Laurent data (exact). $H_s$ is the residue of $S_2$ at $\theta_s$, $b_s$ that of $S_3$, $a$ the common double pole of $S_3$. Sources: P1, P3, P4 section 0.}\label{tab:vac}\\
\toprule
$\Lambda$ & twist & $H_1$ & $H_2$ & $b_1$ & $b_2$ & $a$ \\
\midrule
\endfirsthead
\toprule
$\Lambda$ & twist & $H_1$ & $H_2$ & $b_1$ & $b_2$ & $a$ \\
\midrule
\endhead
\bottomrule
\endlastfoot
$[2,1,0]$ & $\omega$ & $12539/930$ & $4139/930$ & $2173/310$ & $493/310$ & $9/5$ \\
$[2,1,0]$ & $\omega'$ & $40463/4340$ & $1263/4340$ & $30501/4340$ & $-19899/4340$ & $27/7$ \\
$[3,1,0]$ & $\omega$ & $5789/310$ & $1589/310$ & $3027/310$ & $507/310$ & $12/5$ \\
$[3,1,0]$ & $\omega'$ & $54107/4340$ & $-4693/4340$ & $40869/4340$ & $-34731/4340$ & $36/7$ \\
$[3,2,0]$ & $\omega$ & $885/31$ & $45/31$ & $2376/155$ & $-144/155$ & $24/5$ \\
$[3,2,0]$ & $\omega'$ & $48341/2170$ & $-10459/2170$ & $44217/2170$ & $-31383/2170$ & $72/7$ \\
\end{longtable}}


Check against (\ref{eq:vacdata}) for $[3,2,0]$, $\omega$: $H_1=2\omega_1+3\omega_2+5\omega_3+12/\theta_{12}=5+7+3+420/31=15+420/31=885/31$.

\section{The kernels}\label{app:K}

\subsection{The adjoint kernel}\label{app:Kadj}

{\scriptsize
\begin{longtable}{cp{13.5cm}}
\caption{The adjoint kernel $G_{ij}(\omega_1,\omega_1',t)$, $t=1/\theta_{12}$, in the basis $N=(P^3p_{02},\,P^{3\prime}p_{02},\,P^3p_{03},\,P^{3\prime}p_{03},\,P^3p_{23})$ (weights $2,3,3,4,5$); $G=4K$. Source: \texttt{RESULTS-\allowbreak{}twovac-\allowbreak{}limit.\allowbreak{}md} section L2(d), parsed from the agent's printout (D2); equal to the independent fit of D1 in 25/25 entries.}\label{tab:G210}\\
\toprule
$(i,j)$ & $G_{ij}$ \\
\midrule
\endfirsthead
\toprule
$(i,j)$ & $G_{ij}$ \\
\midrule
\endhead
\bottomrule
\endlastfoot
1,1 & $-4 \omega_1^{2} \omega_1'^{2} (3 t^{2} + t \omega_1 - t \omega_1' - \omega_1 \omega_1')$ \\
1,2 & $2 (6 t^{3} \omega_1^{2} + 24 t^{3} \omega_1 \omega_1' + 6 t^{3} \omega_1'^{2} + 2 t^{2} \omega_1^{3} + 12 t^{2} \omega_1^{2} \omega_1' - 2 t^{2} \omega_1'^{3} + 2 t \omega_1^{3} \omega_1' - 6 t \omega_1^{2} \omega_1'^{2} - 2 t \omega_1 \omega_1'^{3} - 3 \omega_1^{3} \omega_1'^{2})$ \\
1,3 & $-2 \omega_1^{2} \omega_1'^{2} (2 t + \omega_1)$ \\
1,4 & $4 (-3 t^{4} - 4 t^{3} \omega_1 - 2 t^{3} \omega_1' - t^{2} \omega_1^{2} + 3 t^{2} \omega_1 \omega_1' + t^{2} \omega_1'^{2} + 3 t \omega_1^{2} \omega_1' + t \omega_1 \omega_1'^{2} + \omega_1^{3} \omega_1')$ \\
1,5 & $-2 (6 t^{3} + 8 t^{2} \omega_1 + 4 t^{2} \omega_1' + 4 t \omega_1^{2} + 2 t \omega_1 \omega_1' + \omega_1^{3})$ \\
2,1 & $-2 (6 t^{3} \omega_1^{2} + 24 t^{3} \omega_1 \omega_1' + 6 t^{3} \omega_1'^{2} + 2 t^{2} \omega_1^{3} - 12 t^{2} \omega_1 \omega_1'^{2} - 2 t^{2} \omega_1'^{3} - 2 t \omega_1^{3} \omega_1' - 6 t \omega_1^{2} \omega_1'^{2} + 2 t \omega_1 \omega_1'^{3} + 3 \omega_1^{2} \omega_1'^{3})$ \\
2,2 & $-12 t^{2} \omega_1^{2} - 48 t^{2} \omega_1 \omega_1' - 12 t^{2} \omega_1'^{2} - 4 t \omega_1^{3} - 12 t \omega_1^{2} \omega_1' + 12 t \omega_1 \omega_1'^{2} + 4 t \omega_1'^{3} + 9 \omega_1^{2} \omega_1'^{2}$ \\
2,3 & $12 t^{4} + 16 t^{3} \omega_1 + 8 t^{3} \omega_1' + 4 t^{2} \omega_1^{2} - 12 t^{2} \omega_1 \omega_1' - 4 t^{2} \omega_1'^{2} - 4 t \omega_1^{2} \omega_1' + 4 t \omega_1 \omega_1'^{2} + 3 \omega_1^{2} \omega_1'^{2}$ \\
2,4 & $2 (6 t^{3} + 8 t^{2} \omega_1 + 4 t^{2} \omega_1' + 2 t \omega_1^{2} - 6 t \omega_1 \omega_1' - 2 t \omega_1'^{2} - 3 \omega_1^{2} \omega_1')$ \\
2,5 & $6 t^{2} + 8 t \omega_1 + 4 t \omega_1' + 3 \omega_1^{2}$ \\
3,1 & $2 \omega_1^{2} \omega_1'^{2} (2 t - \omega_1')$ \\
3,2 & $12 t^{4} - 8 t^{3} \omega_1 - 16 t^{3} \omega_1' - 4 t^{2} \omega_1^{2} - 12 t^{2} \omega_1 \omega_1' + 4 t^{2} \omega_1'^{2} - 4 t \omega_1^{2} \omega_1' + 4 t \omega_1 \omega_1'^{2} + 3 \omega_1^{2} \omega_1'^{2}$ \\
3,3 & $\omega_1^{2} \omega_1'^{2}$ \\
3,4 & $-2 (-6 t^{3} - 2 t^{2} \omega_1 + 2 t^{2} \omega_1' + 2 t \omega_1 \omega_1' + \omega_1^{2} \omega_1')$ \\
3,5 & $6 t^{2} + 4 t \omega_1 + \omega_1^{2}$ \\
4,1 & $4 (-3 t^{4} + 2 t^{3} \omega_1 + 4 t^{3} \omega_1' + t^{2} \omega_1^{2} + 3 t^{2} \omega_1 \omega_1' - t^{2} \omega_1'^{2} - t \omega_1^{2} \omega_1' - 3 t \omega_1 \omega_1'^{2} + \omega_1 \omega_1'^{3})$ \\
4,2 & $2 (-6 t^{3} + 4 t^{2} \omega_1 + 8 t^{2} \omega_1' + 2 t \omega_1^{2} + 6 t \omega_1 \omega_1' - 2 t \omega_1'^{2} - 3 \omega_1 \omega_1'^{2})$ \\
4,3 & $-2 (6 t^{3} + 2 t^{2} \omega_1 - 2 t^{2} \omega_1' - 2 t \omega_1 \omega_1' + \omega_1 \omega_1'^{2})$ \\
4,4 & $-4 (3 t^{2} + t \omega_1 - t \omega_1' - \omega_1 \omega_1')$ \\
4,5 & $-2 (2 t + \omega_1)$ \\
5,1 & $-2 (-6 t^{3} + 4 t^{2} \omega_1 + 8 t^{2} \omega_1' - 2 t \omega_1 \omega_1' - 4 t \omega_1'^{2} + \omega_1'^{3})$ \\
5,2 & $6 t^{2} - 4 t \omega_1 - 8 t \omega_1' + 3 \omega_1'^{2}$ \\
5,3 & $6 t^{2} - 4 t \omega_1' + \omega_1'^{2}$ \\
5,4 & $2 (2 t - \omega_1')$ \\
5,5 & $1$ \\
\end{longtable}}


{\scriptsize
\begin{longtable}{cp{13.5cm}}
\caption{The adjoint kernel in twisted jets $\tilde N$ built from $e^{-\omega_su}P^j$ ($\omega_s=\omega_1$ on site 1, $\omega_1'$ on site 2): it depends only on $\delta=\omega_1'-\omega_1$ and $t=1/\theta_{12}$. Same basis order as Table~\ref{tab:G210}.
Source: \texttt{RESULTS-\allowbreak{}twovac-\allowbreak{}limit.\allowbreak{}md}.}\label{tab:Gtw}\\
\toprule
$(i,j)$ & $\tilde G_{ij}$ \\
\midrule
\endfirsthead
\toprule
$(i,j)$ & $\tilde G_{ij}$ \\
\midrule
\endhead
\bottomrule
\endlastfoot
1,1 & $-4 \delta^{2} t^{2} (\delta + 3 t)(2 \delta + 3 t)$ \\
1,2 & $-4 \delta t^{2} (2 \delta + 3 t)^{2}$ \\
1,3 & $-4 \delta t^{2} (\delta^{2} + 5 \delta t + 3 t^{2})$ \\
1,4 & $-4 t^{2} (3 \delta^{2} + 8 \delta t + 3 t^{2})$ \\
1,5 & $-4 t^{2} (2 \delta + 3 t)$ \\
2,1 & $4 \delta t^{2} (2 \delta + 3 t)^{2}$ \\
2,2 & $2 \delta t (\delta + 6 t)(2 \delta + 3 t)$ \\
2,3 & $2 t^{2} (5 \delta^{2} + 10 \delta t + 6 t^{2})$ \\
2,4 & $4 t (\delta^{2} + 5 \delta t + 3 t^{2})$ \\
2,5 & $2 t (2 \delta + 3 t)$ \\
3,1 & $4 \delta t^{2} (\delta^{2} + 5 \delta t + 3 t^{2})$ \\
3,2 & $2 t^{2} (5 \delta^{2} + 10 \delta t + 6 t^{2})$ \\
3,3 & $2 \delta t^{2} (\delta + 6 t)$ \\
3,4 & $4 t^{2} (2 \delta + 3 t)$ \\
3,5 & $6 t^{2}$ \\
4,1 & $-4 t^{2} (3 \delta^{2} + 8 \delta t + 3 t^{2})$ \\
4,2 & $-4 t (\delta^{2} + 5 \delta t + 3 t^{2})$ \\
4,3 & $-4 t^{2} (2 \delta + 3 t)$ \\
4,4 & $-4 t (\delta + 3 t)$ \\
4,5 & $-4 t$ \\
5,1 & $4 t^{2} (2 \delta + 3 t)$ \\
5,2 & $2 t (2 \delta + 3 t)$ \\
5,3 & $6 t^{2}$ \\
5,4 & $4 t$ \\
5,5 & $1$ \\
\end{longtable}}


For the adjoint the canonical $\Theta$ is $G/(4\omega_1^3\omega_1'^3)$ with $\omega=1/a$. A few entries in this normalisation (G1's K0 table, pairs $(02),(03),(12),(13),(23)$, $a'=a$ primed, $t=1/\theta_{12}$): $\Theta_{(02)(02)}=1+(a-a')t-3aa't^2$, $\Theta_{(02)(03)}=-a'(2at+1)/2$, $\Theta_{(03)(02)}=a(2a't-1)/2$, $\Theta_{(03)(03)}=aa'/4$, $\Theta_{(23)(23)}=a^3a'^3/4$, $\Theta_{(13)(13)}=-a^2a'^2(3aa't^2-at+a't-1)$, $\Theta_{(03)(23)}=aa'^3(6a^2t^2+4at+1)/4$. The full table follows from $G$ above or from the exchange rule $\Theta_{ij}(a,a',t)=\Theta_{ji}(a',a,-t)$; all 25 entries are stored exactly in \texttt{general\_\allowbreak{}kernel\_\allowbreak{}rep\_\allowbreak{}210.\allowbreak{}pkl}.

\subsection{\texorpdfstring{$[3,1,0]$}{[3,1,0]} and \texorpdfstring{$[3,2,0]$}{[3,2,0]}: base-point kernels}\label{app:K310}

The exact kernel $K$ at the base point $(\omega_1,\omega_1',\theta_{12})=(5/2,7/4,31/35)$ was evaluated from G1's stored exact polynomial laws (\texttt{general\_\allowbreak{}kernel\_\allowbreak{}rep\_\allowbreak{}310.\allowbreak{}pkl}) in the basis of pairs $(02),(03),(04),(12),(13),(14),(23),(24),(34)$, rows $i$ and columns $j$; $K_{99}=1/9$. By the block theorem (\ref{eq:blockK}) the kernel of $[3,2,0]$ is the sub-block with rows and columns $(03),(04),(13),(14),(23),(24),(34)$ (we verified exact equality of all 49 entries, at this point and at $(-2/5,3,29/45)$); it reproduces the base-point values printed by P4 (for example $K_{11}=-2477531875/3936256$, $K_{12}=-15649375/95232$, $K_{17}=1241317375/2770563$, $K_{71}=-3124391459/709264128$, $K_{77}=1/9$ in the $[3,2,0]$ order).

\begin{itemize}
\item Row $(02)$: $\frac{315367650959375}{30261936128}$, $\frac{718176615625}{244047872}$, $\frac{6342284375}{5904384}$, $-\frac{6585110076593125}{117265002496}$, $\frac{4028315502843156875}{655020316754688}$, $-\frac{92562651185431875}{14086458424832}$, $\frac{37273401868703875}{1760807303104}$, $-\frac{84056493864875}{113600471168}$, $-\frac{6592621876625}{2748398496}$.
\item Row $(03)$: $-\frac{1163517096875}{488095744}$, $-\frac{2477531875}{3936256}$, $-\frac{15649375}{95232}$, $\frac{56471003621623654875}{3493441689358336}$, $-\frac{102534604765838875}{42259375274496}$, $\frac{201835815824125}{113600471168}$, $-\frac{1066229095827875}{227200942336}$, $\frac{841817515875}{1832265664}$, $\frac{1241317375}{2770563}$.
\item Row $(04)$: $\frac{4287885875}{23617536}$, $\frac{8703625}{190464}$, $\frac{42875}{4608}$, $-\frac{59815795579096375}{42259375274496}$, $\frac{213468283144375}{1022404240512}$, $-\frac{1720011432125}{10993593984}$, $\frac{29908228000}{85887453}$, $-\frac{4528239625}{88658016}$, $-\frac{66781625}{2144952}$.
\item Row $(12)$: $\frac{2958159365546875}{117265002496}$, $-\frac{26499034046453294375}{1746720844679168}$, $\frac{38085651450599375}{10564843818624}$, $\frac{142414335777875}{3782742016}$, $-\frac{409392543660729125}{42259375274496}$, $\frac{917570019491125}{227200942336}$, $-\frac{1924121671270925}{227200942336}$, $\frac{399981133475}{458066416}$, $\frac{30381041525}{44329008}$.
\item Row $(13)$: $-\frac{45956030935735311125}{5240162534037504}$, $\frac{92531938123379125}{21129687637248}$, $-\frac{1336404322859375}{1022404240512}$, $-\frac{1017838029413828875}{84518750548992}$, $\frac{2973626853294125}{1022404240512}$, $-\frac{7269823449875}{5496796992}$, $\frac{12484966188275}{5496796992}$, $-\frac{1902394025}{5541126}$, $-\frac{179699825}{1072476}$.
\item Row $(14)$: $\frac{94733549284435075}{112691667398656}$, $-\frac{319234542990475}{908803769344}$, $\frac{2843820794675}{21987187968}$, $\frac{496844069229025}{454401884672}$, $-\frac{651463216075}{2748398496}$, $\frac{7355513725}{59105344}$, $-\frac{5263492675}{29552672}$, $\frac{16750475}{476656}$, $\frac{147925}{11532}$.
\item Row $(23)$: $\frac{4226124769018075}{1760807303104}$, $-\frac{347834314964675}{454401884672}$, $\frac{1716815095975}{5496796992}$, $\frac{1917549022254845}{908803769344}$, $-\frac{5119851097655}{10993593984}$, $\frac{13720239855}{59105344}$, $-\frac{17517770585}{59105344}$, $\frac{14281435}{238328}$, $\frac{227285}{11532}$.
\item Row $(24)$: $-\frac{584809663796365}{1817607538688}$, $\frac{1123221350545}{14658125312}$, $-\frac{14416082485}{354632064}$, $-\frac{462993234235}{1832265664}$, $\frac{4168591945}{88658016}$, $-\frac{27137425}{953312}$, $\frac{6967205}{238328}$, $-\frac{57785}{7688}$, $-\frac{365}{186}$.
\item Row $(34)$: $\frac{2805654394403}{87948751872}$, $-\frac{3124391459}{709264128}$, $\frac{60638627}{17159616}$, $\frac{6925335571}{354632064}$, $-\frac{2906729}{1072476}$, $\frac{100009}{46128}$, $-\frac{37121}{23064}$, $\frac{203}{372}$, $\frac{1}{9}$.
\end{itemize}

Laws (all verified, Table~\ref{tab:laws} and Section~\ref{sec:structure}): each entry is a polynomial in $(\omega_1,\omega_1')$. For $[3,2,0]$ the table of total degrees is rows $8,7,7,6,6,5,4$ / $7,6,6,5,5,4,3$ (twice) / $6,5,5,4,4,3,2$ (twice) / $5,4,4,3,3,2,1$ / $4,3,3,2,2,1,0$, i.e.\ $d_{ij}=8-(e_i+e_j)$ with $e=(0,1,1,2,2,3,4)$, at most degree 4 in each of $\omega_1,\omega_1'$; for $[3,1,0]$, $e=(0,1,2,1,2,3,3,4,5)$, bidegree $(5,10)$ in $(1/\omega_1',t)$ at fixed $a$, 998 monomials; the dependence on $\theta$ is through $\theta_{12}$ only (equal-$\theta_{12}$ pairs agree to $4.8\cdot10^{-72}$, $8\cdot10^{-70}$, $2.4\cdot10^{-68}$ for $[2,1,0],[2,0,0],[3,2,0]$).

\section{The one-site kernel}\label{app:C}

Table~\ref{tab:onesite} lists the coefficients of $\theta_{kl}(a)=c_{kl}a^{k+l-\lambda_2-1}$ for all two-row and one-row representations up to $\lambda_1=4$. They are computed from the closed form (\ref{eq:ckl}); we cross-checked the entries printed by G1 for $[1,0,0],[1,1,0],[2,0,0],[2,1,0],[2,2,0],[3,1,0],[3,2,0],[4,2,0]$ (they agree).

{\scriptsize
\begin{longtable}{lcp{11cm}}
\caption{The one-site kernel $\theta_{kl}(a)=c_{kl}\,a^{k+l-\lambda_2-1}$: coefficients $c_{kl}$ from the closed form (eq.~\ref{eq:ckl}) for all two-row (and one-row) representations up to $\lambda_1=4$. Pairs $(kl)$ in lexicographic order, $0\le k<l\le\lambda_1+1$, $l\ge\lambda_2+1$. The entries for $\Lambda\le[4,2,0]$ coincide with the exact one-site fits of G1; the rest were predicted by the formula and then confirmed (17 representations in all).}\label{tab:onesite}\\
\toprule
$\Lambda$ & \# pairs & $c_{kl}$ \\
\midrule
\endfirsthead
\toprule
$\Lambda$ & \# pairs & $c_{kl}$ \\
\midrule
\endhead
\bottomrule
\endlastfoot
$[1,0,0]$ & 3 & (01): $1$, (02): $-1$, (12): $1$ \\
$[1,1,0]$ & 2 & (02): $1$, (12): $-1$ \\
$[2,0,0]$ & 6 & (01): $1$, (02): $-2$, (03): $1$, (12): $3$, (13): $-2$, (23): $1$ \\
$[2,1,0]$ & 5 & (02): $1$, (03): $-\frac{1}{2}$, (12): $-\frac{3}{2}$, (13): $1$, (23): $-\frac{1}{2}$ \\
$[2,2,0]$ & 3 & (03): $1$, (13): $-2$, (23): $1$ \\
$[3,0,0]$ & 10 & (01): $1$, (02): $-3$, (03): $3$, (04): $-1$, (12): $6$, (13): $-8$, (14): $3$, (23): $6$, (24): $-3$, (34): $1$ \\
$[3,1,0]$ & 9 & (02): $1$, (03): $-1$, (04): $\frac{1}{3}$, (12): $-2$, (13): $\frac{8}{3}$, (14): $-1$, (23): $-2$, (24): $1$, (34): $-\frac{1}{3}$ \\
$[3,2,0]$ & 7 & (03): $1$, (04): $-\frac{1}{3}$, (13): $-\frac{8}{3}$, (14): $1$, (23): $2$, (24): $-1$, (34): $\frac{1}{3}$ \\
$[3,3,0]$ & 4 & (04): $1$, (14): $-3$, (24): $3$, (34): $-1$ \\
$[4,0,0]$ & 15 & (01): $1$, (02): $-4$, (03): $6$, (04): $-4$, (05): $1$, (12): $10$, (13): $-20$, (14): $15$, (15): $-4$, (23): $20$, (24): $-20$, (25): $6$, (34): $10$, (35): $-4$, (45): $1$ \\
$[4,1,0]$ & 14 & (02): $1$, (03): $-\frac{3}{2}$, (04): $1$, (05): $-\frac{1}{4}$, (12): $-\frac{5}{2}$, (13): $5$, (14): $-\frac{15}{4}$, (15): $1$, (23): $-5$, (24): $5$, (25): $-\frac{3}{2}$, (34): $-\frac{5}{2}$, (35): $1$, (45): $-\frac{1}{4}$ \\
$[4,2,0]$ & 12 & (03): $1$, (04): $-\frac{2}{3}$, (05): $\frac{1}{6}$, (13): $-\frac{10}{3}$, (14): $\frac{5}{2}$, (15): $-\frac{2}{3}$, (23): $\frac{10}{3}$, (24): $-\frac{10}{3}$, (25): $1$, (34): $\frac{5}{3}$, (35): $-\frac{2}{3}$, (45): $\frac{1}{6}$ \\
$[4,3,0]$ & 9 & (04): $1$, (05): $-\frac{1}{4}$, (14): $-\frac{15}{4}$, (15): $1$, (24): $5$, (25): $-\frac{3}{2}$, (34): $-\frac{5}{2}$, (35): $1$, (45): $-\frac{1}{4}$ \\
$[4,4,0]$ & 5 & (05): $1$, (15): $-4$, (25): $6$, (35): $-4$, (45): $1$ \\
\end{longtable}}


\section{The \texorpdfstring{$\epsilon$}{}-series and the CG tables}\label{app:D}

{\scriptsize
\begin{longtable}{llccc}
\caption{Exact $\epsilon$-series coefficients of the Laurent data of the six adjoint branches at $\phi=(1/3,-2/7,5/11)$, $\theta_{12}=31/35$ (so $H_1^{(0)}=4/\theta_{12}$ etc.), orders $n=0,1,2$ ($H_s=\sum_n\epsilon^nH_s^{(n)}$).
Source: \texttt{ref/\allowbreak{}untwisted-\allowbreak{}exact-\allowbreak{}210.\allowbreak{}txt}, record \texttt{E4\_\allowbreak{}SERIES} (agent W4; exact Rayleigh--Schr\"odinger).}\label{tab:series-ex}\\
\toprule
branch & series & $n=0$ & $n=1$ & $n=2$ \\
\midrule
\endfirsthead
\toprule
branch & series & $n=0$ & $n=1$ & $n=2$ \\
\midrule
\endhead
\bottomrule
\endlastfoot
$(4,2,0)$ & $H_1$ & $\frac{140}{31}$ & $\frac{260}{231}$ & 0 \\
$(4,2,0)$ & $a_1$ & 0 & $\frac{15}{11}$ & 0 \\
$(4,2,0)$ & $b_1$ & 0 & $\frac{700}{341}$ & $-\frac{25}{231}$ \\
$(4,1,1)$ & $H_1$ & $\frac{210}{31}$ & $\frac{349}{462}$ & $\frac{100719}{1660120}$ \\
$(4,1,1)$ & $a_1$ & $\frac{210}{31}$ & $\frac{39}{154}$ & $\frac{302157}{1660120}$ \\
$(4,1,1)$ & $b_1$ & 0 & $\frac{1160}{1023}$ & $-\frac{107}{462}$ \\
$(3,3,0)$ & $H_1$ & $\frac{210}{31}$ & $\frac{221}{154}$ & $\frac{5239}{123480}$ \\
$(3,3,0)$ & $a_1$ & 0 & $\frac{15}{11}$ & 0 \\
$(3,3,0)$ & $b_1$ & 0 & $\frac{1050}{341}$ & $\frac{5}{154}$ \\
$(3,2,1) copy 1$ & $H_1$ & $\frac{315}{31}$ & $\frac{82}{77}+\frac{157}{1155}\sqrt{5}$ & $\frac{4645474}{140072625}$ \\
$(3,2,1) copy 1$ & $a_1$ & $\frac{525}{62}-\frac{105}{62}\sqrt{5}$ & $\frac{509}{770}+\frac{541}{2310}\sqrt{5}$ & $\frac{7649591}{140072625}-\frac{13875662}{700363125}\sqrt{5}$ \\
$(3,2,1) copy 1$ & $b_1$ & $\frac{3675}{961}\sqrt{5}$ & $\frac{76}{31}$ & $-\frac{1}{11}+\frac{1952029}{20010375}\sqrt{5}$ \\
$(3,2,1) copy 2$ & $H_1$ & $\frac{315}{31}$ & $\frac{82}{77}-\frac{157}{1155}\sqrt{5}$ & $\frac{4645474}{140072625}$ \\
$(3,2,1) copy 2$ & $a_1$ & $\frac{525}{62}+\frac{105}{62}\sqrt{5}$ & $\frac{509}{770}-\frac{541}{2310}\sqrt{5}$ & $\frac{7649591}{140072625}+\frac{13875662}{700363125}\sqrt{5}$ \\
$(3,2,1) copy 2$ & $b_1$ & $-\frac{3675}{961}\sqrt{5}$ & $\frac{76}{31}$ & $-\frac{1}{11}-\frac{1952029}{20010375}\sqrt{5}$ \\
$(2,2,2)$ & $H_1$ & $\frac{420}{31}$ & $\frac{232}{231}$ & $\frac{782347}{11205810}$ \\
$(2,2,2)$ & $a_1$ & $\frac{420}{31}$ & $\frac{116}{231}$ & $\frac{782347}{7470540}$ \\
$(2,2,2)$ & $b_1$ & 0 & $\frac{2320}{1023}$ & $-\frac{17}{231}$ \\
\end{longtable}}


{\tiny
\begin{longtable}{llrrrrrrrr}
\caption{The same series, numerically (6 significant digits), orders $n=0,\dots,7$; exact values to order 7 are in the data file.
Odd orders vanish identically on the algebraic branches $(4,1,1)$ and $(3,3,0)$ beyond $n=1$ and on the singlet for $H_1$ (see text).}\label{tab:series-num}\\
\toprule
branch & series & $n=0$ & $n=1$ & $n=2$ & $n=3$ & $n=4$ & $n=5$ & $n=6$ & $n=7$ \\
\midrule
\endfirsthead
\toprule
branch & series & $n=0$ & $n=1$ & $n=2$ & $n=3$ & $n=4$ & $n=5$ & $n=6$ & $n=7$ \\
\midrule
\endhead
\bottomrule
\endlastfoot
$(4,2,0)$ & $H_1$ & 4.51613 & 1.12554 & 0 & 0 & 0 & 0 & 0 & 0 \\
$(4,2,0)$ & $a_1$ & 0 & 1.36364 & 0 & 0 & 0 & 0 & 0 & 0 \\
$(4,2,0)$ & $b_1$ & 0 & 2.05279 & -0.108225 & 0 & 0 & 0 & 0 & 0 \\
$(4,1,1)$ & $H_1$ & 6.77419 & 0.755411 & 0.0606697 & 0 & -0.00163007 & 0 & 8.75937e-05 & 0 \\
$(4,1,1)$ & $a_1$ & 6.77419 & 0.253247 & 0.182009 & 0 & -0.00489022 & 0 & 0.000262781 & 0 \\
$(4,1,1)$ & $b_1$ & 0 & 1.13392 & -0.231602 & 0.0202232 & 0 & -0.000543358 & 0 & 2.91979e-05 \\
$(3,3,0)$ & $H_1$ & 6.77419 & 1.43506 & 0.0424279 & 0 & -0.0007972 & 0 & 2.9958e-05 & 0 \\
$(3,3,0)$ & $a_1$ & 0 & 1.36364 & 0 & 0 & 0 & 0 & 0 & 0 \\
$(3,3,0)$ & $b_1$ & 0 & 3.07918 & 0.0324675 & 0.0192854 & 0 & -0.000362364 & 0 & 1.36173e-05 \\
$(3,2,1) copy 1$ & $H_1$ & 10.1613 & 1.36889 & 0.0331648 & -9.34224e-05 & -0.00028439 & 7.10857e-06 & 4.61811e-06 & -2.62555e-07 \\
$(3,2,1) copy 1$ & $a_1$ & 4.68085 & 1.18472 & 0.0103104 & 0.000820064 & -0.000148242 & -8.918e-06 & 3.12722e-06 & 7.11436e-08 \\
$(3,2,1) copy 1$ & $b_1$ & 8.55104 & 2.45161 & 0.127221 & 0.00973109 & -0.000715096 & -5.10438e-05 & 1.19737e-05 & 3.85103e-07 \\
$(3,2,1) copy 2$ & $H_1$ & 10.1613 & 0.760985 & 0.0331648 & 9.34224e-05 & -0.00028439 & -7.10857e-06 & 4.61811e-06 & 2.62555e-07 \\
$(3,2,1) copy 2$ & $a_1$ & 12.2546 & 0.137354 & 0.0989128 & 0.00171494 & -0.000786658 & -4.38194e-05 & 1.17683e-05 & 1.1896e-06 \\
$(3,2,1) copy 2$ & $b_1$ & -8.55104 & 2.45161 & -0.309039 & 0.00973109 & 0.000715096 & -5.10438e-05 & -1.19737e-05 & 3.85103e-07 \\
$(2,2,2)$ & $H_1$ & 13.5484 & 1.00433 & 0.0698162 & 0 & -0.000539655 & 0 & 2.38687e-06 & 0 \\
$(2,2,2)$ & $a_1$ & 13.5484 & 0.502165 & 0.104724 & -0.00635184 & -0.000809482 & 0.000180024 & 3.5803e-06 & -4.80708e-06 \\
$(2,2,2)$ & $b_1$ & 0 & 2.26784 & -0.0735931 & 0.0403721 & 0 & -0.000607703 & 0 & 1.14447e-05 \\
\end{longtable}}


Symbolic data to higher order are in \texttt{untwisted\_\allowbreak{}Q\_\allowbreak{}pert\_\allowbreak{}sym.\allowbreak{}pkl} (singlet to $n=6$, algebraic branches to $n=6$), \texttt{untwisted\_\allowbreak{}Q\_\allowbreak{}laurent\_\allowbreak{}exact.\allowbreak{}pkl} (all branches $n\le5$), \texttt{untwisted\_\allowbreak{}Q\_\allowbreak{}octet\_\allowbreak{}fit.\allowbreak{}pkl} and the data-file records \texttt{E4\_\allowbreak{}SYMBOLIC} (orders $0..3$), \texttt{E4\_\allowbreak{}PSI} (to order 4), \texttt{E4\_\allowbreak{}GNORM}, \texttt{E5\_\allowbreak{}KB}, \texttt{E5\_\allowbreak{}SLICE}, \texttt{E5\_\allowbreak{}GF}. The numerical $\epsilon$-series of the monomials $N_i,M$, the Plücker coordinates and the polynomial coefficients are in \texttt{untwisted\_\allowbreak{}Q\_\allowbreak{}series\_\allowbreak{}210.\allowbreak{}pkl} (12 points, 200 digits, $n\le10$), \texttt{\_\allowbreak{}ext.\allowbreak{}pkl} (24 points, 400 digits, $n\le20$, reliable to 18) and \texttt{\_\allowbreak{}phi2.\allowbreak{}pkl} ($\phi'$).

\subsection{CG tables}\label{app:CG}

{\scriptsize
\begin{longtable}{llcccc>{\raggedleft\arraybackslash}p{2.6cm}}
\caption{Exact CG slices $C_\nu=\langle e_\nu\otimes e_1|\Psi_0\rangle/\|\Psi_0\|_G$ of the highest-weight kets of $V_\Lambda\otimes V_\Lambda$, six representations (45 HW branches); $\nu$ is the site-1 Gelfand--Tsetlin index in notebook order (index 1 = highest weight) with its weight; the two copies in a multiplicity-2 sector are the eigenvectors of $\theta_{12}\,a_1$, copy 1 having the smaller eigenvalue; $C_\nu=(c_A+c_B\sqrt{D_0})/\sqrt{g_A+g_B\sqrt{D_0}}$. Source: \texttt{ref/\allowbreak{}untwisted-\allowbreak{}exact-\allowbreak{}<Lambda>.\allowbreak{}txt}, records \texttt{BRANCH}.}\label{tab:cg}\\
\toprule
$\Lambda$ & $\mu$ & copy & $\nu$, weight & $c_A+c_B\sqrt{D_0}$ & $g_A+g_B\sqrt{D_0}$ & $C_\nu$ \\
\midrule
\endfirsthead
\toprule
$\Lambda$ & $\mu$ & copy & $\nu$, weight & $c_A+c_B\sqrt{D_0}$ & $g_A+g_B\sqrt{D_0}$ & $C_\nu$ \\
\midrule
\endhead
\bottomrule
\endlastfoot
$(1,1,0)$ & $(2,2,0)$ &  & 1, (1,1,0) & $1$ & $1$ & 1.0000000000 \\
$(1,1,0)$ & $(2,1,1)$ &  & 2, (1,0,1) & $-1$ & $2$ & -0.7071067812 \\
$(2,1,0)$ & $(4,2,0)$ &  & 1, (2,1,0) & $1$ & $1$ & 1.0000000000 \\
$(2,1,0)$ & $(4,1,1)$ &  & 2, (2,0,1) & $-1$ & $2$ & -0.7071067812 \\
$(2,1,0)$ & $(3,3,0)$ &  & 3, (1,2,0) & $-1$ & $2$ & -0.7071067812 \\
$(2,1,0)$ & $(3,2,1)$ & 1 & 4, (1,1,1) & $21-9\sqrt{5}$ & $7560-3240\sqrt{5}$ & 0.0493116000 \\
 &  &  & 5, (1,1,1) & $-3+3\sqrt{5}$ &  & 0.2088872897 \\
$(2,1,0)$ & $(3,2,1)$ & 2 & 4, (1,1,1) & $21+9\sqrt{5}$ & $7560+3240\sqrt{5}$ & 0.3379867346 \\
 &  &  & 5, (1,1,1) & $-3-3\sqrt{5}$ &  & -0.0797878449 \\
$(2,1,0)$ & $(2,2,2)$ &  & 8, (0,1,2) & $2$ & $288$ & 0.1178511302 \\
$(2,2,0)$ & $(4,4,0)$ &  & 1, (2,2,0) & $1$ & $1$ & 1.0000000000 \\
$(2,2,0)$ & $(4,3,1)$ &  & 2, (2,1,1) & $-1$ & $4$ & -0.5000000000 \\
$(2,2,0)$ & $(4,2,2)$ &  & 3, (2,0,2) & $1$ & $12$ & 0.2886751346 \\
$(3,1,0)$ & $(6,2,0)$ &  & 1, (3,1,0) & $1$ & $1$ & 1.0000000000 \\
$(3,1,0)$ & $(6,1,1)$ &  & 2, (3,0,1) & $-1$ & $2$ & -0.7071067812 \\
$(3,1,0)$ & $(5,3,0)$ &  & 3, (2,2,0) & $-1$ & $4$ & -0.5000000000 \\
$(3,1,0)$ & $(5,2,1)$ & 1 & 4, (2,1,1) & $264-24\sqrt{105}$ & $1496880-136080\sqrt{105}$ & 0.0564580374 \\
 &  &  & 5, (2,1,1) & $-30+6\sqrt{105}$ &  & 0.0983443476 \\
$(3,1,0)$ & $(5,2,1)$ & 2 & 4, (2,1,1) & $264+24\sqrt{105}$ & $1496880+136080\sqrt{105}$ & 0.2998902852 \\
 &  &  & 5, (2,1,1) & $-30-6\sqrt{105}$ &  & -0.0538008073 \\
$(3,1,0)$ & $(4,4,0)$ &  & 7, (1,3,0) & $1$ & $12$ & 0.2886751346 \\
$(3,1,0)$ & $(4,3,1)$ & 1 & 8, (1,2,1) & $-168+96\sqrt{3}$ & $5225472-2985984\sqrt{3}$ & -0.0074430331 \\
 &  &  & 9, (1,2,1) & $24-24\sqrt{3}$ &  & -0.0758903002 \\
$(3,1,0)$ & $(4,3,1)$ & 2 & 8, (1,2,1) & $-168-96\sqrt{3}$ & $5225472+2985984\sqrt{3}$ & -0.1036680780 \\
 &  &  & 9, (1,2,1) & $24+24\sqrt{3}$ &  & 0.0203347447 \\
$(3,1,0)$ & $(4,2,2)$ &  & 10, (1,1,2) & $18$ & $233280$ & 0.0372677996 \\
 &  &  & 11, (1,1,2) & $-6$ &  & -0.0124225999 \\
$(3,1,0)$ & $(3,3,2)$ &  & 14, (0,2,2) & $-3$ & $17280$ & -0.0228217732 \\
$(3,2,0)$ & $(6,4,0)$ &  & 1, (3,2,0) & $1$ & $1$ & 1.0000000000 \\
$(3,2,0)$ & $(6,3,1)$ &  & 2, (3,1,1) & $-1$ & $4$ & -0.5000000000 \\
$(3,2,0)$ & $(6,2,2)$ &  & 3, (3,0,2) & $1$ & $12$ & 0.2886751346 \\
$(3,2,0)$ & $(5,5,0)$ &  & 4, (2,3,0) & $-1$ & $2$ & -0.7071067812 \\
$(3,2,0)$ & $(5,4,1)$ & 1 & 5, (2,2,1) & $22-2\sqrt{105}$ & $18480-1680\sqrt{105}$ & 0.0423435280 \\
 &  &  & 7, (2,2,1) & $-3+\sqrt{105}$ &  & 0.2037459499 \\
$(3,2,0)$ & $(5,4,1)$ & 2 & 5, (2,2,1) & $22+2\sqrt{105}$ & $18480+1680\sqrt{105}$ & 0.2249177139 \\
 &  &  & 7, (2,2,1) & $-3-\sqrt{105}$ &  & -0.0701153289 \\
$(3,2,0)$ & $(5,3,2)$ & 1 & 6, (2,1,2) & $-168+96\sqrt{3}$ & $5225472-2985984\sqrt{3}$ & -0.0074430331 \\
 &  &  & 8, (2,1,2) & $24-24\sqrt{3}$ &  & -0.0758903002 \\
$(3,2,0)$ & $(5,3,2)$ & 2 & 6, (2,1,2) & $-168-96\sqrt{3}$ & $5225472+2985984\sqrt{3}$ & -0.1036680780 \\
 &  &  & 8, (2,1,2) & $24+24\sqrt{3}$ &  & 0.0203347447 \\
$(3,2,0)$ & $(4,4,2)$ &  & 12, (1,2,2) & $4$ & $2880$ & 0.0745355992 \\
$(3,2,0)$ & $(4,3,3)$ &  & 13, (1,1,3) & $-3$ & $17280$ & -0.0228217732 \\
$(4,1,0)$ & $(8,2,0)$ &  & 1, (4,1,0) & $1$ & $1$ & 1.0000000000 \\
$(4,1,0)$ & $(8,1,1)$ &  & 2, (4,0,1) & $-1$ & $2$ & -0.7071067812 \\
$(4,1,0)$ & $(7,3,0)$ &  & 3, (3,2,0) & $-1$ & $6$ & -0.4082482905 \\
$(4,1,0)$ & $(7,2,1)$ & 1 & 4, (3,1,1) & $450-90\sqrt{21}$ & $5443200-1088640\sqrt{21}$ & 0.0557300243 \\
 &  &  & 5, (3,1,1) & $-42+18\sqrt{21}$ &  & 0.0600589519 \\
$(4,1,0)$ & $(7,2,1)$ & 2 & 4, (3,1,1) & $450+90\sqrt{21}$ & $5443200+1088640\sqrt{21}$ & 0.2670185879 \\
 &  &  & 5, (3,1,1) & $-42-18\sqrt{21}$ &  & -0.0385423778 \\
$(4,1,0)$ & $(6,4,0)$ &  & 7, (2,3,0) & $3$ & $360$ & 0.1581138830 \\
$(4,1,0)$ & $(6,3,1)$ & 1 & 8, (2,2,1) & $-150+60\sqrt{6}$ & $6912000-2764800\sqrt{6}$ & -0.0081097820 \\
 &  &  & 9, (2,2,1) & $18-12\sqrt{6}$ &  & -0.0304894699 \\
$(4,1,0)$ & $(6,3,1)$ & 2 & 8, (2,2,1) & $-150-60\sqrt{6}$ & $6912000+2764800\sqrt{6}$ & -0.0802785656 \\
 &  &  & 9, (2,2,1) & $18+12\sqrt{6}$ &  & 0.0128118003 \\
$(4,1,0)$ & $(6,2,2)$ &  & 10, (2,1,2) & $40$ & $4838400$ & 0.0181848242 \\
 &  &  & 11, (2,1,2) & $-10$ &  & -0.0045462060 \\
$(4,1,0)$ & $(5,5,0)$ &  & 13, (1,4,0) & $-1$ & $144$ & -0.0833333333 \\
$(4,1,0)$ & $(5,4,1)$ & 1 & 14, (1,3,1) & $690-150\sqrt{21}$ & $1669248000-362880000\sqrt{21}$ & 0.0010394114 \\
 &  &  & 15, (1,3,1) & $-90+30\sqrt{21}$ &  & 0.0188810654 \\
$(4,1,0)$ & $(5,4,1)$ & 2 & 14, (1,3,1) & $690+150\sqrt{21}$ & $1669248000+362880000\sqrt{21}$ & 0.0238611846 \\
 &  &  & 15, (1,3,1) & $-90-30\sqrt{21}$ &  & -0.0039407078 \\
$(4,1,0)$ & $(5,3,2)$ &  & 16, (1,2,2) & $-120$ & $261273600$ & -0.0074239234 \\
 &  &  & 17, (1,2,2) & $30$ &  & 0.0018559808 \\
$(4,1,0)$ & $(4,4,2)$ &  & 22, (0,3,2) & $4$ & $1036800$ & 0.0039283710 \\
\end{longtable}}


\section{The untwisted data}\label{app:E}

{\scriptsize
\begin{longtable}{lllrccccc}
\caption{Untwisted Laurent data and oper data on all 45 highest-weight branches: $H_1=-H_2=\eta/\theta_{12}$, $a_1=-a_2=\alpha/\theta_{12}$, $b_1=-b_2=\beta/\theta_{12}^2$; exponents of the dual oper at infinity $(\mu_1+2,\mu_2+1,\mu_3)$ (the three polynomial solutions have exactly these degrees); degrees of the three polynomial solutions of the untwisted direct oper (see text for the two exceptions); last column: whether the oper is a ${}_3F_2$ equation (consecutive exponent pair at a singular point).
Source: \texttt{ref/\allowbreak{}untwisted-\allowbreak{}exact-\allowbreak{}<Lambda>.\allowbreak{}txt}.}\label{tab:untw}\\
\toprule
$\Lambda$ & $\mu$ & copy & $\eta$ & $\alpha$ & $\beta$ & at $\infty$ & direct deg. & ${}_3F_2$ \\
\midrule
\endfirsthead
\toprule
$\Lambda$ & $\mu$ & copy & $\eta$ & $\alpha$ & $\beta$ & at $\infty$ & direct deg. & ${}_3F_2$ \\
\midrule
\endhead
\bottomrule
\endlastfoot
$(1,1,0)$ & $(2,2,0)$ &  & $2$ & $0$ & $0$ & $(4,3,0)$ & $(0,2,4)$ & yes \\
$(1,1,0)$ & $(2,1,1)$ &  & $4$ & $4$ & $0$ & $(4,2,1)$ & $(3,2,0)$ & yes \\
$(2,1,0)$ & $(4,2,0)$ &  & $4$ & $0$ & $0$ & $(6,3,0)$ & $(6,3,0)$ & no \\
$(2,1,0)$ & $(4,1,1)$ &  & $6$ & $6$ & $0$ & $(6,2,1)$ & $(5,4,0)$ & yes \\
$(2,1,0)$ & $(3,3,0)$ &  & $6$ & $0$ & $0$ & $(5,4,0)$ & $(6,2,1)$ & yes \\
$(2,1,0)$ & $(3,2,1)$ & 1 & $9$ & $\frac{15}{2}-\frac{3}{2}\sqrt{5}$ & $3\sqrt{5}$ & $(5,3,1)$ & $(5,3,1)$ & no \\
$(2,1,0)$ & $(3,2,1)$ & 2 & $9$ & $\frac{15}{2}+\frac{3}{2}\sqrt{5}$ & $-3\sqrt{5}$ & $(5,3,1)$ & $(5,3,1)$ & no \\
$(2,1,0)$ & $(2,2,2)$ &  & $12$ & $12$ & $0$ & $(4,3,2)$ & $(4,3,2)$ & yes \\
$(2,2,0)$ & $(4,4,0)$ &  & $8$ & $0$ & $0$ & $(6,5,0)$ & $(0,2,6)$ & yes \\
$(2,2,0)$ & $(4,3,1)$ &  & $12$ & $12$ & $0$ & $(6,4,1)$ & $(5,2,0)$ & yes \\
$(2,2,0)$ & $(4,2,2)$ &  & $14$ & $18$ & $0$ & $(6,3,2)$ & $(4,3,0)$ & yes \\
$(3,1,0)$ & $(6,2,0)$ &  & $6$ & $0$ & $0$ & $(8,3,0)$ & $(8,5,0)$ & no \\
$(3,1,0)$ & $(6,1,1)$ &  & $8$ & $8$ & $0$ & $(8,2,1)$ & $(7,6,0)$ & yes \\
$(3,1,0)$ & $(5,3,0)$ &  & $10$ & $0$ & $0$ & $(7,4,0)$ & $(8,4,1)$ & no \\
$(3,1,0)$ & $(5,2,1)$ & 1 & $13$ & $\frac{21}{2}-\frac{1}{2}\sqrt{105}$ & $\sqrt{105}$ & $(7,3,1)$ & $(7,5,1)$ & no \\
$(3,1,0)$ & $(5,2,1)$ & 2 & $13$ & $\frac{21}{2}+\frac{1}{2}\sqrt{105}$ & $-\sqrt{105}$ & $(7,3,1)$ & $(7,5,1)$ & no \\
$(3,1,0)$ & $(4,4,0)$ &  & $12$ & $0$ & $0$ & $(6,5,0)$ & $(8,3,2)$ & yes \\
$(3,1,0)$ & $(4,3,1)$ & 1 & $16$ & $12-4\sqrt{3}$ & $8\sqrt{3}$ & $(6,4,1)$ & $(7,4,2)$ & no \\
$(3,1,0)$ & $(4,3,1)$ & 2 & $16$ & $12+4\sqrt{3}$ & $-8\sqrt{3}$ & $(6,4,1)$ & $(7,4,2)$ & no \\
$(3,1,0)$ & $(4,2,2)$ &  & $18$ & $18$ & $0$ & $(6,3,2)$ & $(6,5,2)$ & yes \\
$(3,1,0)$ & $(3,3,2)$ &  & $20$ & $20$ & $0$ & $(5,4,2)$ & $(6,4,3)$ & yes \\
$(3,2,0)$ & $(6,4,0)$ &  & $12$ & $0$ & $0$ & $(8,5,0)$ & $(8,3,0)$ & no \\
$(3,2,0)$ & $(6,3,1)$ &  & $16$ & $16$ & $0$ & $(8,4,1)$ & $(7,4,0)$ & no \\
$(3,2,0)$ & $(6,2,2)$ &  & $18$ & $24$ & $0$ & $(8,3,2)$ & $(6,5,0)$ & yes \\
$(3,2,0)$ & $(5,5,0)$ &  & $14$ & $0$ & $0$ & $(7,6,0)$ & $(8,2,1)$ & yes \\
$(3,2,0)$ & $(5,4,1)$ & 1 & $19$ & $\frac{35}{2}-\frac{1}{2}\sqrt{105}$ & $\sqrt{105}$ & $(7,5,1)$ & $(7,3,1)$ & no \\
$(3,2,0)$ & $(5,4,1)$ & 2 & $19$ & $\frac{35}{2}+\frac{1}{2}\sqrt{105}$ & $-\sqrt{105}$ & $(7,5,1)$ & $(7,3,1)$ & no \\
$(3,2,0)$ & $(5,3,2)$ & 1 & $22$ & $28-4\sqrt{3}$ & $8\sqrt{3}$ & $(7,4,2)$ & $(6,4,1)$ & no \\
$(3,2,0)$ & $(5,3,2)$ & 2 & $22$ & $28+4\sqrt{3}$ & $-8\sqrt{3}$ & $(7,4,2)$ & $(6,4,1)$ & no \\
$(3,2,0)$ & $(4,4,2)$ &  & $24$ & $30$ & $0$ & $(6,5,2)$ & $(6,3,2)$ & yes \\
$(3,2,0)$ & $(4,3,3)$ &  & $26$ & $36$ & $0$ & $(6,4,3)$ & $(5,4,2)$ & yes \\
$(4,1,0)$ & $(8,2,0)$ &  & $8$ & $0$ & $0$ & $(10,3,0)$ & $(10,7,0)$ & no \\
$(4,1,0)$ & $(8,1,1)$ &  & $10$ & $10$ & $0$ & $(10,2,1)$ & $(9,8,0)$ & yes \\
$(4,1,0)$ & $(7,3,0)$ &  & $14$ & $0$ & $0$ & $(9,4,0)$ & $(10,6,1)$ & no \\
$(4,1,0)$ & $(7,2,1)$ & 1 & $17$ & $\frac{27}{2}-\frac{3}{2}\sqrt{21}$ & $3\sqrt{21}$ & $(9,3,1)$ & $(9,7,1)$ & no \\
$(4,1,0)$ & $(7,2,1)$ & 2 & $17$ & $\frac{27}{2}+\frac{3}{2}\sqrt{21}$ & $-3\sqrt{21}$ & $(9,3,1)$ & $(9,7,1)$ & no \\
$(4,1,0)$ & $(6,4,0)$ &  & $18$ & $0$ & $0$ & $(8,5,0)$ & $(10,5,2)$ & no \\
$(4,1,0)$ & $(6,3,1)$ & 1 & $22$ & $16-4\sqrt{6}$ & $8\sqrt{6}$ & $(8,4,1)$ & $(9,6,2)$ & no \\
$(4,1,0)$ & $(6,3,1)$ & 2 & $22$ & $16+4\sqrt{6}$ & $-8\sqrt{6}$ & $(8,4,1)$ & $(9,6,2)$ & no \\
$(4,1,0)$ & $(6,2,2)$ &  & $24$ & $24$ & $0$ & $(8,3,2)$ & $(8,7,2)$ & yes \\
$(4,1,0)$ & $(5,5,0)$ &  & $20$ & $0$ & $0$ & $(7,6,0)$ & $(10,4,3)$ & yes \\
$(4,1,0)$ & $(5,4,1)$ & 1 & $25$ & $\frac{35}{2}-\frac{5}{2}\sqrt{21}$ & $5\sqrt{21}$ & $(7,5,1)$ & $(9,5,3)$ & no \\
$(4,1,0)$ & $(5,4,1)$ & 2 & $25$ & $\frac{35}{2}+\frac{5}{2}\sqrt{21}$ & $-5\sqrt{21}$ & $(7,5,1)$ & $(9,5,3)$ & no \\
$(4,1,0)$ & $(5,3,2)$ &  & $28$ & $28$ & $0$ & $(7,4,2)$ & $(8,6,3)$ & no \\
$(4,1,0)$ & $(4,4,2)$ &  & $30$ & $30$ & $0$ & $(6,5,2)$ & $(8,5,4)$ & yes \\
\end{longtable}}


The sectors of each $V_\Lambda\otimes V_\Lambda$ and their multiplicities: $[1,1,0]$: $(2,2,0),(2,1,1)$. $[2,1,0]$: $(4,2,0),(4,1,1),(3,3,0),(3,2,1)\times2,(2,2,2)$. $[2,2,0]$: $(4,4,0),(4,3,1),(4,2,2)$. $[3,1,0]$: $(6,2,0),(6,1,1),(5,3,0),(5,2,1)\times2,(4,4,0),(4,3,1)\times2,(4,2,2),(3,3,2)$. $[3,2,0]$: $(6,4,0),(6,3,1),(6,2,2),(5,5,0),(5,4,1)\times2,(5,3,2)\times2,(4,4,2),(4,3,3)$. $[4,1,0]$: $(8,2,0),(8,1,1),(7,3,0),(7,2,1)\times2,(6,4,0),(6,3,1)\times2,(6,2,2),(5,5,0),(5,4,1)\times2,(5,3,2),(4,4,2)$. Multiplicity $\le2$ everywhere (45 branches in all).
One-site Casimirs $(C_2,C_3)(\Lambda)$ from the data files: $[1,1,0]$: $(4,8)$; $[2,1,0]$: $(9,30)$; $[2,2,0]$: $(12,36)$; $[3,1,0]$: $(16,72)$; $[3,2,0]$: $(19,77)$; $[4,1,0]$: $(25,140)$. The Laurent operators in the canonical HW basis, the Gram matrices and the HW kets themselves are in the files (records \texttt{LAURMAT}, \texttt{GRAM}, \texttt{HWBASIS}).

\section{Reproducibility}\label{app:repro}

All Python work lives in \texttt{Paul/\allowbreak{}Python/\allowbreak{}gaudin\_\allowbreak{}sl3/\allowbreak{}} (run with \texttt{python -\allowbreak{}X utf8}, Python 3.13, sympy 1.14, mpmath 1.3, numpy; no \texttt{uv}), all notebooks in \texttt{Paul/\allowbreak{}Mathematica/\allowbreak{}Gaudin/\allowbreak{}Experiments/\allowbreak{}}, reference data in \texttt{Paul/\allowbreak{}Mathematica/\allowbreak{}Gaudin/\allowbreak{}ref/\allowbreak{}}. Nothing was committed by the agents; about $230$ MB of \texttt{*.\allowbreak{}pkl} caches in \texttt{Paul/\allowbreak{}Python/\allowbreak{}gaudin\_\allowbreak{}sl3/\allowbreak{}} should be gitignored before committing.
The four notebooks store no native outputs (results are in markdown cells below each code cell); Paul reruns them in VS Code.
The common pitfall: run heavy stages separately (an all-stage run stalled under CPU contention from parallel agents).

{\footnotesize
\begin{longtable}{p{1.0cm}p{4.2cm}p{4.6cm}p{3.6cm}}
\caption{Which file produces what.
Times are those reported by the agents (warm = with caches, cold = from scratch).
Specs are \texttt{SPEC-\allowbreak{}*.\allowbreak{}md}/\texttt{*-\allowbreak{}spec.\allowbreak{}md} next to each file.}\label{tab:repro}\\
\toprule
agent & files & produces & time, precision \\
\midrule
\endfirsthead
\toprule
agent & files & produces & time, precision \\
\midrule
\endhead
\bottomrule
\endlastfoot
W1 & \texttt{Experiments/\allowbreak{}Gaudin-\allowbreak{}L2-\allowbreak{}chain.\allowbreak{}wb} (48 cells, 16 code) & Hamiltonians, vacuum, localisers, two-vacua untwist and CG for $[2,1,0]$, $[3,1,0]$ (switch $\Lambda$ in the first code cell) & one clean run, 200 digits \\
W2 & \texttt{Experiments/\allowbreak{}Gaudin-\allowbreak{}L2-\allowbreak{}chain-\allowbreak{}2.\allowbreak{}wb} (80 cells, 21 code) & $(\infty,\infty)$ family, $[3,1,0]$ fourth localiser, closed forms, A5 reference tables & every code cell run twice, 200 digits \\
W3 & \texttt{Experiments/\allowbreak{}Gaudin-\allowbreak{}L2-\allowbreak{}chain-\allowbreak{}320.\allowbreak{}wb} (75 cells, 32 code) & $[3,2,0]$ chain side, cleaning verdict, A7 tables & one clean run, 8 min; scan cell 59 about 6 min cold \\
W4 & \texttt{Experiments/\allowbreak{}Gaudin-\allowbreak{}untwisted-\allowbreak{}exact.\allowbreak{}wb} (39 cells, 14 code); \texttt{ref/\allowbreak{}untwisted-\allowbreak{}exact-\allowbreak{}\{110,210,220,310,320,410\}.\allowbreak{}txt} & exact untwisted reference data, perturbation theory, CG tables & one clean run, exact \\
P1 & \texttt{gaudin\_\allowbreak{}chain.\allowbreak{}py}, \texttt{qside.\allowbreak{}py}, \texttt{RESULTS-\allowbreak{}Qside.\allowbreak{}md}, \texttt{qside\_\allowbreak{}output.\allowbreak{}txt}, \texttt{qside\_\allowbreak{}cache.\allowbreak{}pkl} & Q side for $[2,1,0]$ (opers, Wronskians, $\det N$, regime) & about 75 s; 60/100 digits \\
P2 & \texttt{xxx\_\allowbreak{}chain.\allowbreak{}py}, \texttt{xxx\_\allowbreak{}limit.\allowbreak{}py}, \texttt{RESULTS-\allowbreak{}xxx-\allowbreak{}limit.\allowbreak{}md} & XXX$\to$Gaudin limit, dictionary, $x_0$, localiser prediction & about 2 min (X4 sympy 95 s); exact series \\
P3 & \texttt{qside\_\allowbreak{}310.\allowbreak{}py}, \texttt{RESULTS-\allowbreak{}Qside-\allowbreak{}310.\allowbreak{}md}, \texttt{qside\_\allowbreak{}310\_\allowbreak{}*.\allowbreak{}pkl} & $[3,1,0]$ Q side; rank-4 rule & 4 min warm, 15 min cold; 60/100 digits \\
P4 & \texttt{qside\_\allowbreak{}320.\allowbreak{}py}, \texttt{twovac\_\allowbreak{}fit\_\allowbreak{}320.\allowbreak{}py} (stages T0..T7), \texttt{RESULTS-\allowbreak{}Qside-\allowbreak{}320.\allowbreak{}md}, \texttt{twovac\_\allowbreak{}fit\_\allowbreak{}320\_\allowbreak{}*.\allowbreak{}pkl} & $[3,2,0]$ Q side and jet formula, $[3,1,0]$ comparison & 4 min warm / 8 min cold; 60/80/100 digits \\
D1 & \texttt{twovac\_\allowbreak{}fit.\allowbreak{}py}, \texttt{RESULTS-\allowbreak{}twovac-\allowbreak{}fit.\allowbreak{}md}, \texttt{twovac\_\allowbreak{}fit\_\allowbreak{}*.\allowbreak{}pkl} & direct fit of the adjoint kernel, kernel laws, regime Q-side CG & 2.5 min warm; 80 digits (F4: 120) \\
D2 & \texttt{twovac\_\allowbreak{}limit.\allowbreak{}py}, \texttt{RESULTS-\allowbreak{}twovac-\allowbreak{}limit.\allowbreak{}md}; input \texttt{ref/\allowbreak{}xxx-\allowbreak{}210-\allowbreak{}infinf-\allowbreak{}formula-\allowbreak{}cells.\allowbreak{}txt} & $h\to0$ limit of the ten-product formula, $G$ & about 1 min; exact, 60 digits for L3 \\
D3 & \texttt{gaudin\_\allowbreak{}B.\allowbreak{}py}, \texttt{RESULTS-\allowbreak{}gaudin-\allowbreak{}B.\allowbreak{}md}, \texttt{gaudin\_\allowbreak{}B\_\allowbreak{}*.\allowbreak{}pkl} & $B$ hierarchy, jet basis, K tensor, $G$ from first principles & 6 min warm / 15 min cold; exact Fractions, 80 digits for states; \texttt{GB\_\allowbreak{}THETA=1/\allowbreak{}5,-\allowbreak{}4/\allowbreak{}9} repeats at another $\theta$ \\
G1 & \texttt{general\_\allowbreak{}kernel.\allowbreak{}py} (stages K1, REP, ZC, BLK, SGL, CORNER, ANALYZE), \texttt{general\_\allowbreak{}kernel\_\allowbreak{}bld.\allowbreak{}py}, \texttt{RESULTS-\allowbreak{}general-\allowbreak{}kernel.\allowbreak{}md}, \texttt{general\_\allowbreak{}kernel\_\allowbreak{}rep\_\allowbreak{}<Lambda>.\allowbreak{}pkl} & general-$\Lambda$ kernels, one-site closed form, block theorem & 80 digits; per point $\sim35$ s ($[3,1,0]$, $[3,2,0]$), $[4,1,0]$ about 2 min, $[4,2,0]$ about 4 min \\
U1 & \texttt{untwisted\_\allowbreak{}Q.\allowbreak{}py} (stages T0, T1, EXPORT), \texttt{RESULTS-\allowbreak{}untwisted-\allowbreak{}Q.\allowbreak{}md}, \texttt{untwisted\_\allowbreak{}Q\_\allowbreak{}*.\allowbreak{}pkl} & untwisted Laurent data, oper, $\epsilon$-series & 200/400 digits \\
U2 & \texttt{untwist\_\allowbreak{}limit.\allowbreak{}py} (stages R, N, NU), \texttt{untwist\_\allowbreak{}limit\_\allowbreak{}frob.\allowbreak{}py}, \texttt{RESULTS-\allowbreak{}untwist-\allowbreak{}limit.\allowbreak{}md}, \texttt{untwist\_\allowbreak{}limit\_\allowbreak{}reg\_\allowbreak{}<k>.\allowbreak{}pkl} & analytic untwisting limit for the adjoint & about 1 min with caches; 200 digits, 12 $\epsilon$ points $\sim1$ s each \\
data & \texttt{ref/\allowbreak{}ref-\allowbreak{}*.\allowbreak{}txt} & cell dumps of the XXX notebooks used by the Python agents & -- \\
\end{longtable}}


Other pitfalls recorded by the agents: a numerically vanishing feature ($p_{01}\sim10^{-78}$) must be zeroed explicitly before QR or it is counted as an independent feature; the $100$-digit regime solver needs a kernel threshold of $10^{-50}$ (the default $10^{-35}$ gives spurious lower degrees); a threshold-based degree reading of $q$ fails when the leading coefficient is $\epsilon$-suppressed (read the smallest degree admitting a kernel); the memoised \texttt{laurent} cache holds bad values after a fit at a pole; the mp eigensolver returns $10^{-59}$ imaginary noise in some sectors; the all-stage Python run stalled under contention.
The notebooks were built from JSON and run cell by cell in the shared kernel (wolfbook new-notebook workaround), hence no native outputs.

The sources of this document are the plans \texttt{Gaudin-\allowbreak{}sl3-\allowbreak{}CG-\allowbreak{}plan.\allowbreak{}md}, \texttt{-\allowbreak{}round2-\allowbreak{}}, \texttt{-\allowbreak{}round3-\allowbreak{}}, \texttt{-\allowbreak{}round4-\allowbreak{}}, \texttt{-\allowbreak{}round5-\allowbreak{}plan.\allowbreak{}md}, the ten \texttt{RESULTS-\allowbreak{}*.\allowbreak{}md} files, the four notebooks and the six data files named above, and the repository's \texttt{HANDOFF.\allowbreak{}md} (entries from ``sl(3) GAUDIN'' onwards).
The generator of this document (hand-written prose plus tables built from the data files) lives in \texttt{Paul/\allowbreak{}Mathematica/\allowbreak{}Gaudin/\allowbreak{}notes-\allowbreak{}src/\allowbreak{}}.

\end{document}
